% !TeX program = xelatex
% Payne Zero Paper I 中英文对照译本（附物理背景解读）
\documentclass[11pt,a4paper,UTF8,fontset=fandol]{ctexart}

\usepackage[margin=2.3cm]{geometry}
\usepackage{amsmath,amssymb}
\usepackage{graphicx}
\usepackage{booktabs}
\usepackage{longtable}
\usepackage{array}
\usepackage{algpseudocode}
\usepackage{capt-of}
\usepackage[most]{tcolorbox}
\usepackage{xcolor}
\usepackage{natbib}
\usepackage{hyperref}
\hypersetup{colorlinks=true, linkcolor=black, citecolor=zhblue, urlcolor=zhblue}

\graphicspath{{../arXiv-2607.24141v1/}}

\definecolor{engray}{RGB}{105,105,105}
\definecolor{zhblue}{RGB}{20,60,120}

\newcommand{\paynezero}{\textsc{Payne Zero}}
\newcommand{\tcell}[2]{\parbox[t]{#1}{\raggedright #2}}

% AAS journal macros used in references.bib
\newcommand{\aap}{A\&A}
\newcommand{\aaps}{A\&AS}
\newcommand{\aj}{AJ}
\newcommand{\apj}{ApJ}
\newcommand{\apjl}{ApJL}
\newcommand{\apjs}{ApJS}
\newcommand{\araa}{ARA\&A}
\newcommand{\mnras}{MNRAS}
\newcommand{\pasj}{PASJ}
\newcommand{\pasp}{PASP}
\newcommand{\physscr}{Phys.~Scr.}
\newcommand{\zap}{ZAp}

% 英文原文环境：灰色小字、两端缩进
\newenvironment{EN}{\par\begin{quote}\footnotesize\color{engray}}{\end{quote}}

% 物理背景解读框
\newtcolorbox{physbox}[1]{enhanced, breakable,
  colback=orange!7, colframe=orange!55!black,
  colbacktitle=orange!18, coltitle=black,
  title={物理背景解读｜#1}, fonttitle=\bfseries,
  left=6pt, right=6pt, top=4pt, bottom=4pt,
  boxrule=0.6pt, arc=1.5mm}

% 译者说明框
\newtcolorbox{notebox}[1]{enhanced, breakable,
  colback=zhblue!6, colframe=zhblue!60,
  colbacktitle=zhblue!15, coltitle=black,
  title={#1}, fonttitle=\bfseries,
  left=6pt, right=6pt, top=4pt, bottom=4pt,
  boxrule=0.6pt, arc=1.5mm}

\setlength{\parskip}{2pt}

\title{\textbf{The \paynezero\ Project I\\Stellar Spectra from Physical Models in Seconds}\\[8pt]
\Large \paynezero\ 项目之一：数秒内由物理模型生成恒星光谱\\[6pt]
\large 中英文对照译本 · 附物理背景解读与注释}
\author{原文作者：Yuan-Sen Ting（丁源森）\quad Elliot M. Kim\\
俄亥俄州立大学天文系 / CCAPP / 马普天文研究所；康奈尔大学计算机系}
\date{对照译本整理：2026 年 8 月}

\begin{document}
\maketitle

\begin{notebox}{阅读说明}
\begin{itemize}
  \item 正文段落为\textbf{中文译文}；每段之后的灰色小字引文块为\textbf{英文原文}，便于逐段对照。
  \item 橙色方框为\textbf{物理背景解读}：针对文中涉及的恒星大气、辐射转移、光谱合成等物理概念补充背景知识与直观解释，帮助非恒星物理专业的读者理解。
  \item 公式保持原文形式不变，公式编号沿用原文；图表标题为中英对照，表格内容保留英文并附中文说明。
  \item 本文件为学习用对照译本，非官方译文；如译文与原文有出入，以英文原文为准。原文见 arXiv:2607.24141，代码见 \url{https://github.com/tingyuansen/payne-zero}。
\end{itemize}
\end{notebox}

\begin{notebox}{导读：这篇论文在做什么}
测量一颗恒星的化学成分，本质上是一道“反问题”：给定观测光谱，推断恒星的有效温度、表面重力和各种元素丰度。最可靠的做法是\textbf{正向物理建模}——先用恒星大气理论算出大气结构，再用辐射转移理论算出出射光谱，最后与观测比较。但传统的 Kurucz 程序算一颗星要几十分钟，无法放进需要对每个候选参数反复调用模型的优化器里。于是过去十年学界普遍改用“捷径”：要么预计算一个光谱网格再做插值（光谱模拟器，如 The Payne），要么直接从观测数据学习“标签到光谱”的经验关系（数据驱动模型，如 The Cannon）。这些捷径引入了插值误差和经验相关性，且模型与观测不符时难以定位物理原因。

本文的核心观点很直接：\textbf{现代 GPU/多核 CPU 已经足够快，可以把完整的物理计算重新放回反演循环内部，不再需要任何学习型中间层}（“Zero”即“零层光谱模拟器”）。作者把 Kurucz 的一维 LTE 大气与光谱合成计算重新组织为 GPU 波长并行、CPU 多核大气迭代的形式：300--1000 nm、$R=300{,}000$ 的太阳光谱从约 760 秒降到 14 秒（H100），APOGEE 波段约 1 秒；并训练了一个神经网络专门预测大气迭代的\textbf{初值}（最终大气仍由物理求解器收敛确定）。由此实现三件事：(1) 对 APOGEE 光谱直接做 12 个丰度坐标的物理拟合（每颗星约 40 秒）；(2) 用自动微分对太阳和大角星同时定标超过 $10^5$ 条谱线的振子强度与阻尼参数（约 1 分钟）；(3) 恢复了银河系高/低 $\alpha$ 双化学序列。论文同时诚实地指出了边界：冷星富分子大气的迭代收敛性变差、81 维丰度输入的初始化器采样稀疏、定标结果是模型依赖的有效值而非实验室测量。
\end{notebox}

\section*{关键术语中英对照表}
\begin{longtable}{@{}p{0.38\textwidth}p{0.56\textwidth}@{}}
\toprule
\textbf{English} & \textbf{中文} \\
\midrule
\endfirsthead
\toprule
\textbf{English} & \textbf{中文} \\
\midrule
\endhead
forward model & 前向（正向）模型 \\
model atmosphere & 模型大气 \\
spectral synthesis & 光谱合成 \\
radiative transfer & 辐射转移 \\
LTE / NLTE & 局部热动平衡 / 非局部热动平衡 \\
opacity & 不透明度 \\
continuous / line opacity & 连续 / 谱线不透明度 \\
line blanketing & 谱线覆盖（谱线遮蔽）效应 \\
opacity sampling & 不透明度采样法 \\
Rosseland mean opacity & 罗斯兰平均不透明度 \\
optical depth & 光学深度 \\
column mass & 柱质量 \\
hydrostatic equilibrium & 流体静力学平衡 \\
equation of state & 物态方程 \\
Saha ionization & 萨哈电离（方程） \\
partition function & 配分函数 \\
molecular equilibrium & 分子（化学）平衡 \\
mixing-length convection & 混合长对流 \\
radiative acceleration & 辐射加速度 \\
oscillator strength ($gf$) & 振子强度 \\
damping (van der Waals / Stark / radiative) & 阻尼（范德瓦尔斯 / 斯塔克 / 辐射） \\
Voigt profile & 福格特（谱线）轮廓 \\
Doppler width & 多普勒致宽 \\
microturbulence ($\xi$) & 微观湍流（速度） \\
effective temperature ($T_{\rm eff}$) & 有效温度 \\
surface gravity ($\log g$) & 表面重力（加速度对数） \\
metallicity, [M/H], [X/Fe], [$\alpha$/M] & 金属丰度及相对丰度标记（见正文注释） \\
H$^-$ continuum & 负氢离子连续（吸收）谱 \\
spectral emulator & 光谱模拟器（插值器） \\
data-driven model & 数据驱动模型 \\
label & （恒星）标签，指 $T_{\rm eff}$、$\log g$、丰度等参数 \\
line spread function (LSF) & 线扩展函数 \\
resolving power ($R$) & 分辨本领 \\
fixed point (iteration) & 不动点（迭代） \\
Jacobian / spectral radius & 雅可比矩阵 / 谱半径 \\
principal component analysis (PCA) & 主成分分析 \\
learned initializer & 学习型初始化器 \\
automatic differentiation & 自动微分 \\
L-BFGS / Gauss--Newton & 两种数值优化算法 \\
Fourier transform spectrometer (FTS) atlas & 傅里叶变换光谱仪标准光谱图集 \\
telluric absorption & 地球大气（吸收）污染 \\
APOGEE / ASPCAP & 阿帕奇点天文台银河演化实验巡天 / 其恒星参数与丰度流水线 \\
red giant / red clump & 红巨星 / 红团簇（星） \\
high-/low-$\alpha$ sequence & 高 / 低 $\alpha$（元素丰度）序列 \\
dredge-up & （对流）挖掘上翻 \\
abundance zero point & 丰度零点 \\
equivalent width & 等值宽度 \\
\bottomrule
\end{longtable}

\section*{摘要 Abstract}

现代恒星巡天测量着数以百万计的光谱，然而计算一个自洽的大气模型及其光谱可能需要数十分钟。这一计算成本催生了预计算网格、光谱模拟器和数据驱动模型。我们提出 \paynezero：它将一维 LTE Kurucz 计算重组为 GPU 原生的光谱合成与多核 CPU 大气迭代，并对照原始 Fortran 程序完成验证。在 NVIDIA H100 GPU 上，以 $R_{\rm grid}=300{,}000$ 采样的 300--1000 nm 太阳光谱约需 14 秒；APOGEE 1500--1700 nm 区间约需 1 秒。物理大气迭代在 16 线程 AMD CPU 上需 2--5 秒，学习型初始化器进一步减少了收敛所需的迭代次数。在所测试的矮星与巨星参数区间内，最终光谱与原版计算保持实用层面的一致。这样的速度使直接合成可以被放进优化器内部，而无需“标签到流量”的光谱模拟器。我们演示了对约减后 APOGEE 光谱的多元素直接拟合，恢复出与巡天星表大体一致的多元素丰度趋势。GPU 常驻的视向速度平移、致宽、线扩展函数卷积与探测器采样相对合成的额外开销可以忽略。直接合成搜索在 H100 上每颗星耗时不到一分钟，而大气验证在多核 CPU 上独立运行。同一计算图还能针对太阳和大角星联合定标超过 $10^5$ 个振子强度与阻尼改正，在 H100 上约需一分钟。因此，\paynezero\ 把直接物理拟合与原子数据定标推进到了巡天规模。代码开源于 \url{https://github.com/tingyuansen/payne-zero}。

\begin{EN}
Modern stellar surveys measure millions of spectra, yet one self-consistent
atmosphere and spectrum can require tens of minutes. This cost has motivated
grids, spectral emulators, and data-driven models. We present \paynezero, which
reorganizes one-dimensional LTE Kurucz calculations for GPU-native synthesis
and multicore atmosphere iteration, and validate it against the original
Fortran programs. A
300--1000~nm solar spectrum sampled at $R_{\rm grid}=300{,}000$ takes about
14~s on an NVIDIA H100 GPU, while the APOGEE 1500--1700~nm interval takes about
1~s. Physical atmosphere iterations take 2--5~s on 16 AMD CPU threads, and
learned initializers reduce the iterations required for convergence. Final
spectra remain in practical parity across the tested dwarf and giant regimes.
These speeds place direct synthesis inside an optimizer without a
label-to-flux spectral emulator. We demonstrate direct many-element fitting of
reduced APOGEE spectra and recover multi-element abundance trends broadly
consistent with the survey catalog.
GPU-resident velocity shifts, broadening, line-spread-function convolution,
and detector sampling add negligible cost relative to synthesis.
The direct-synthesis search takes less than one minute per star on an H100, while
atmosphere verification runs independently on multicore CPUs. The same
computational graph calibrates more than $10^5$ oscillator-strength and
damping corrections jointly against the Sun and
Arcturus in about one minute on an H100. \paynezero\ therefore brings direct
physical fitting and atomic-data calibration to survey scale. The code is
available at \url{https://github.com/tingyuansen/payne-zero}.
\end{EN}

\noindent\textbf{关键词（原文）：}恒星大气；恒星光谱线；天文软件；天文数据分析；计算方法。


\section{引言 Introduction}\label{sec:introduction}

大型光谱巡天已经测量、正在测量或即将测量数十万到数百万颗恒星的光谱。主要的银河系巡天项目包括 APOGEE、GALAH、Gaia--ESO、LAMOST、Gaia RVS、SDSS-V、DESI、WEAVE、4MOST 和 PFS \citep{Cui2012,Dalton2012,Gilmore2012,Takada2014,DeSilva2015,Kollmeier2017,Majewski2017,Chiappini2019,Cooper2023,Creevey2023}。这些巡天的科学回报取决于能否把每条光谱转化为恒星性质，包括有效温度、表面重力、化学元素丰度和视向速度。而前向模型的速度与精度为这一分析设定了基本上限。

\begin{EN}
Large spectroscopic surveys have measured, are measuring, or will soon measure
spectra for hundreds of thousands to millions of stars. Major Milky Way
programs include APOGEE, GALAH, Gaia--ESO, LAMOST, Gaia RVS, SDSS-V, DESI,
WEAVE, 4MOST, and PFS
\citep{Cui2012,Dalton2012,Gilmore2012,Takada2014,DeSilva2015,
Kollmeier2017,Majewski2017,Chiappini2019,Cooper2023,Creevey2023}.
Their scientific return depends on turning
each spectrum into stellar properties, including effective temperatures,
surface gravities, chemical abundances, and radial velocities.
The speed and accuracy of the forward model set a basic limit on that analysis.
\end{EN}

一个物理前向模型由三个阶段组成。模型大气描述温度、压强、电子密度和化学粒子布居随深度的变化。光谱合成把这一结构与谱线表结合，计算出射辐射。仪器模型再把内禀光谱映射到观测像素上。大气为合成提供热力学与化学状态，因此同一份元素组成同时控制着大气结构和出射光谱。

\begin{EN}
A physical forward model combines three stages. A model atmosphere describes
the temperature, pressure, electron density, and chemical populations with
depth. Spectral synthesis combines this structure with a line list to calculate
the emergent radiation. An instrument model maps the intrinsic spectrum to the
observed pixels. The atmosphere supplies the thermodynamic and chemical state
used by synthesis, so the same composition controls both the atmospheric
structure and the emergent spectrum.
\end{EN}

\begin{physbox}{为什么大气结构与光谱必须自洽}
恒星大气不是一块均匀发光体：从表面向内部，温度、压强和密度都在变化，不同深度的气体对光谱的贡献不同。大气模型要回答的就是“每一层有多热、多密、由什么粒子组成”；光谱合成再回答“这样的分层结构向外发出什么样的光”。二者被同一份元素丰度耦合在一起——改变镁的丰度，不仅改变镁的谱线，还会改变大气里的自由电子供应，从而间接改变其他元素谱线的强度。因此严谨的分析不能只调某条谱线附近的局部模型，而要让丰度变化同时贯穿大气结构和光谱合成。这正是本文反复强调的“自洽”（self-consistent）的含义，也是直接物理拟合区别于局部经验拟合的关键。
\end{physbox}

所要求的丰度模式作为整体进入物态方程和不透明度。例如，镁既改变自身的谱线，也改变电子供应；而电子供体的变化会改变 H$^-$ 连续谱，进而影响其他元素产生的谱线特征 \citep{John1988,Meszaros2012,Matsuno2024}。碳、氮、氧会重新分配 CO、CN、OH 等分子中的原子，因此可以改变并不属于被扰动元素的光谱特征 \citep{Tsuji1973,Ting2018Oxygen}。来自众多金属的谱线覆盖（line blanketing）改变辐射场，从而改变温度结构 \citep{Gustafsson2008}。不透明度采样大气求解器会针对所要求的化学混合物重新计算粒子布居和不透明度 \citep{Kurucz1996}。

\begin{EN}
The requested abundance pattern enters the equation of state and opacity as a
whole. Magnesium, for example, changes both its own lines and the electron
supply, while electron donors alter the H$^-$ continuum and thereby features
from other species \citep{John1988,Meszaros2012,Matsuno2024}. Carbon, nitrogen,
and oxygen redistribute atoms among molecules such as CO, CN, and OH and can
therefore change spectral features not belonging to the perturbed element
\citep{Tsuji1973,Ting2018Oxygen}. Line blanketing from many metals changes the
radiation field and therefore the temperature structure
\citep{Gustafsson2008}. The opacity-sampling atmosphere solver recomputes the
populations and opacity for the requested mixture \citep{Kurucz1996}.
\end{EN}

\begin{physbox}{丰度为什么会“牵一发而动全身”}
这一段列出了丰度耦合的三条物理通道，值得逐条理解：
\begin{itemize}
\item \textbf{电子供应与 H$^-$ 连续谱。}在太阳型及更冷的恒星中，氢几乎全部呈中性，自由电子主要来自镁、硅、铁等容易电离的金属。这些自由电子与中性氢结合形成负氢离子 H$^-$，而 H$^-$ 是光学—近红外波段连续吸收的主要来源。于是增加金属丰度 $\Rightarrow$ 电子增多 $\Rightarrow$ H$^-$ 吸收增强 $\Rightarrow$ 连续谱被压低，所有元素的谱线对比度都会改变——哪怕那些元素本身的丰度没有变。
\item \textbf{分子平衡。}碳、氮、氧通过 CO、CN、OH 等分子互相“抢夺”原子（CO 结合能极高，会锁死丰度较低的那种元素）。改变 C 或 O 的丰度会重组整个分子体系，从而改变 CN、OH 等分子带的光谱特征。
\item \textbf{谱线覆盖（line blanketing）。}数百万条金属谱线叠加起来像一块“毛玻璃”，把短波辐射挡住、重新以长波辐射放出，改变大气的辐射场和温度结构。这意味着丰度变化会反过来改变大气本身，而不只是光谱。
\end{itemize}
这就是为什么本文坚持“丰度必须作为整体进入计算”，也是逐条谱线孤立测量的经典方法所遗漏的物理。
\end{physbox}

长久以来，运行时间决定了这项计算的哪些部分能够进入实际分析。经典的等值宽度工作可以在预计算网格内插值得到模型大气，并只围绕一条孤立谱线合成一个很小的波段 \citep{CastelliKurucz2003}。长期沿用与现代的代码和分析框架包括 MOOG、Turbospectrum、FERRE、iSpec、Korg、SME 和 PySME \citep{ValentiPiskunov1996,Plez2012,MOOG2012,BlancoCuaresma2014,PiskunovValenti2017,FERRE2023,Wehrhahn2023,Wheeler2023}。巡天级的全光谱分析则不同，它从宽阔波长区间中提取信息，其中包含孤立谱线、混合（blend）谱线，以及冷星中的分子带 \citep{ValentiFischer2005,Ting2017Prospects}。

\begin{EN}
Runtime has long determined which parts of this calculation enter an analysis.
Classical equivalent-width work can obtain a model atmosphere by interpolation
within a precomputed grid and synthesize a small region around an isolated
line \citep{CastelliKurucz2003}. Long-standing and modern codes and analysis
frameworks include MOOG, Turbospectrum, FERRE, iSpec, Korg, SME, and PySME
\citep{ValentiPiskunov1996,Plez2012,MOOG2012,BlancoCuaresma2014,
PiskunovValenti2017,FERRE2023,Wehrhahn2023,Wheeler2023}. Full-spectrum survey
analyses instead draw information from broad wavelength intervals
containing isolated lines, blends, and, in cool stars, molecular bands
\citep{ValentiFischer2005,Ting2017Prospects}.
\end{EN}

对于每一个新的恒星状态，原始的 Kurucz 工作流程以串行程序和中间文件的方式分阶段执行大气与合成计算 \citep{Kurucz2005}。\citet{KimTing2026} 经过验证的 Python 重实现为我们的重新设计提供了直接的软件与物理学谱系。在下文的基准测试中，原始 Fortran 工作流程中依赖标签的大气设置、物理迭代和宽波段合成合计需要数十分钟。若把这一计算放进全光谱优化器中反复执行，单颗星的分析就可能需要数小时。

\begin{EN}
For each new stellar state, the
original Kurucz workflow stages atmosphere and synthesis through serial
programs and intermediate files \citep{Kurucz2005}. The validated Python
reimplementation of \citet{KimTing2026} provides the direct
software and physics lineage for our redesign. In the benchmarks below, the
label-dependent atmosphere setup, physical iterations, and broad synthesis in
the original Fortran workflow together take tens of minutes. Repeating that
calculation inside a full-spectrum optimizer can then take hours for one star.
\end{EN}

光谱模拟器通过对合成光谱库进行插值来避开这一成本；其大气插值误差和流量插值误差都可以被直接测量 \citep{MeszarosAllendePrieto2013}。实现方式包括概率插值（如 Starfish）和神经插值（如 The Payne）\citep{Czekala2015,Ting2019}。但模拟器仍然有有限的训练范围和插值误差，更换物理模型或标签集合就需要重新生成网格 \citep{Rozanski2025a}。当许多种元素的丰度各自独立变化时，标签空间的体积迅速膨胀，训练问题随之变得更难 \citep{Casey2016,Ting2016}。

\begin{EN}
Spectral emulators avoid this cost by interpolating synthetic libraries, a
practice whose atmosphere- and flux-interpolation errors can be measured
directly \citep{MeszarosAllendePrieto2013}. Their
implementations include probabilistic interpolation, as in Starfish, and
neural interpolation, as in The Payne \citep{Czekala2015,Ting2019}. An emulator still
has a finite training range and interpolation error, and a new physical model
or label set requires a new grid \citep{Rozanski2025a}. The training problem
becomes harder when many individual abundances vary because the volume of label
space grows rapidly \citep{Casey2016,Ting2016}.
\end{EN}

数据驱动模型——例如 The Cannon 及其正则化、跨巡天扩展，StarNet 和 astroNN——则改为从观测参考光谱中学习“标签—流量”或“标签推断”关系 \citep{Ness2015,Casey2016,Ho2017,Fabbro2018,LeungBovy2019}。训练标签之间的相关性会把物理敏感性与天体物理相关性混在一起 \citep{Ting2017LAMOST,Xiang2019}。诸如 Cycle-StarNet 的领域自适应方法在合成与观测光谱域之间建立映射，以减少可重复的系统性差异 \citep{OBriain2021}。但无论隐空间坐标还是域映射，都不一定能指认某个光谱特征究竟对应哪种原子物理输入。

\begin{EN}
Data-driven models such as The Cannon and
its regularized and cross-survey extensions, StarNet, and astroNN instead learn
label--flux or label-inference relations
from observed reference spectra
\citep{Ness2015,Casey2016,Ho2017,Fabbro2018,LeungBovy2019}.
Correlated training labels can mix physical sensitivity with astrophysical
correlations \citep{Ting2017LAMOST,Xiang2019}. Related domain-adaptation
methods such as Cycle-StarNet map synthetic and observed spectral domains to
reduce repeatable discrepancies \citep{OBriain2021}. Neither a latent
coordinate nor a domain mapping necessarily identifies the atomic input
responsible for a feature.
\end{EN}

\begin{physbox}{模拟器、数据驱动模型与“归因”问题}
这一段是全文的方法论分水岭。三条路线分别是：
\begin{itemize}
\item \textbf{光谱模拟器（The Payne、Starfish）}：先用物理模型算一个光谱网格，再用神经网络或高斯过程插值。快，但有插值误差，且每增加一个自由丰度维度，训练网格的体积呈指数增长（“维度灾难”）。
\item \textbf{数据驱动模型（The Cannon、StarNet、astroNN）}：直接学习“观测光谱 $\rightarrow$ 恒星标签”的经验关系。问题在于\textbf{归因}：银河系里恒星的年龄、金属丰度、$\alpha$ 丰度天然相关，模型可能靠这种统计相关性“猜对”标签，而不是真正识别出元素谱线的物理响应——训练集之外就可能失效。
\item \textbf{直接物理拟合（本文）}：每次试验都实时计算物理光谱。模型与观测不符时，差异可以被明确归因到具体的物理输入（某条谱线的振子强度、某个阻尼参数、大气结构或缺失的谱线），并逐一检验。
\end{itemize}
\end{physbox}

我们可以改用观测光谱来修正具名的物理输入。振子强度决定原子谱线的强度，阻尼参数决定谱线的翼部 \citep{ValentiPiskunov1996,Gray2005}。这些量的来源包括实验室测量、理论计算、数据库评估以及相对基准恒星的定标 \citep{Piskunov1995,Lobel2011,Jofre2014,Ryabchikova2015,Laverick2019,Smith2021}。一个快速且可微分的计算允许许多这类输入被联合调整，同时每项修正始终绑定在一个物理参数上。

\begin{EN}
We can instead use observed spectra to correct named physical inputs.
Oscillator strengths set atomic line strengths, while damping parameters set
their wings \citep{ValentiPiskunov1996,Gray2005}. These quantities combine
laboratory measurements, theory, database evaluation, and calibration against
benchmark stars
\citep{Piskunov1995,Lobel2011,Jofre2014,Ryabchikova2015,Laverick2019,Smith2021}.
A fast differentiable calculation allows
many such inputs to be adjusted together while each correction remains tied to
a physical parameter.
\end{EN}

当模型与数据不符时，这一区分尤为重要。一个足够灵活的流量模型可以把残差“吸收”掉，却不判断它究竟来自振子强度、阻尼翼、大气结构还是一条缺失的跃迁。直接的物理求值则让这些可能性保持显式且可分别检验。它还允许新的丰度坐标或原子参数直接进入控制方程，而无需重新生成合成光谱网格。

\begin{EN}
This distinction is important when a model misses the data. A flexible flux
model may absorb the residual without deciding whether it came from an
oscillator strength, a damping wing, the atmospheric structure, or a missing
transition. Direct physical evaluation keeps these possibilities explicit and
separately testable. It also allows a new abundance coordinate or atomic
parameter to enter the governing calculation without regenerating a synthetic
flux grid.
\end{EN}

因此，直接物理推断要求一个快得多的计算。大气计算通过相继的物理迭代推进一个深度耦合的状态。结构一旦确定，数万到数百万个波长采样就可以同时求值。Kurucz、MARCS、PHOENIX 和 TLUSTY 等成熟框架做出了不同的物理与数值选择，包括 LTE 冷星网格和非 LTE 热星网格 \citep{HubenyLanz1995,LanzHubeny2003,Kurucz2005,LanzHubeny2007,Gustafsson2008,Husser2013}。我们聚焦于原始的 Kurucz 程序，其串行各阶段通过中间文件交换数据。

\begin{EN}
Direct physical inference therefore requires a much shorter calculation. An
atmosphere calculation advances a depth-coupled state through sequential
physical iterations. Once that structure is known, tens of thousands to millions of
wavelength samples can be evaluated together. Established frameworks such as
Kurucz, MARCS, PHOENIX, and TLUSTY make different physical and numerical
choices, including LTE cool-star and non-LTE hot-star grids
\citep{HubenyLanz1995,LanzHubeny2003,Kurucz2005,LanzHubeny2007,
Gustafsson2008,Husser2013}. We focus on the original Kurucz programs, whose
serial stages exchange intermediate files.
\end{EN}

\paynezero\ 将其支持的一维 LTE Kurucz 物理重组为面向 GPU 合成与多核 CPU 大气迭代的形式。固定数据常驻内存，大型编译算子把相互独立的工作暴露出来。由此得到的前向模型可以在优化器内部被求值。

\begin{EN}
\paynezero\ reorganizes the supported one-dimensional LTE
Kurucz physics for GPU synthesis and multicore CPU atmosphere iteration. Fixed
data remain resident, and large compiled operations expose the independent
work. The resulting forward model can be evaluated inside an optimizer.
\end{EN}

更短的运行时间使直接光谱拟合与原子数据定标成为可能。在建立合成、大气迭代和学习型初始化之后，我们首先测试一个归一化光谱拟合器，其最终候选解用收敛的物理大气重新计算；随后针对太阳和大角星定标原子谱线参数；最后把这两部分与仪器响应、速度、致宽、连续谱和像素不确定度结合起来，完成 APOGEE 应用。

\begin{EN}
The shorter runtime enables direct spectrum fitting and atomic-data
calibration. After establishing synthesis, atmosphere iteration, and learned
initialization, we first test a normalized-spectrum fitter whose final
candidate is recalculated with a converged physical atmosphere. We then
calibrate atomic line parameters against the Sun and Arcturus and combine both
pieces with the instrumental response, velocity, broadening, continuum, and
pixel uncertainties for the APOGEE application.
\end{EN}


\section{光谱合成 Spectral Synthesis}\label{sec:synthesis}

光谱合成把模型大气和谱线表变成出射光谱。它计算贯穿大气的连续与谱线不透明度，然后求解辐射转移以确定离开恒星表面的辐射。\paynezero\ 用一种新的波长并行计算实现了其所支持的 Kurucz 合成物理 \citep{Kurucz2005}。

\begin{EN}
Spectral synthesis turns a model atmosphere and a line list into an emergent
spectrum. It calculates continuous and line opacity through the atmosphere,
then solves radiative transfer to determine the radiation that leaves the
stellar surface. \paynezero\ implements the supported
Kurucz synthesis physics \citep{Kurucz2005} in a new
wavelength-parallel calculation.
\end{EN}

深度状态一旦固定，不同波长之间的相互依赖就微乎其微。因此，宽波段光谱可以把一个长长的波长循环替换成同一短深度计算的许多副本。正是这种并行结构，使下文中的宽波段 Kurucz 合成基准从十余分钟降到数秒——而无需任何模拟器。我们先讲合成，因为它最直接地暴露了本文的核心计算思想。在展开计算之前，表~\ref{tab:atmosphere_symbols} 汇总了全文使用的符号。

\begin{EN}
Once the depth state is fixed, different wavelengths have little dependence on
one another. A broad spectrum can therefore replace one long wavelength loop
with many copies of the same short depth calculation. This parallel structure
is why the broad Kurucz synthesis benchmark below moves from more than ten
minutes to seconds without an emulator. We begin with synthesis because it exposes the central
computational idea directly.
Before developing that calculation, Table~\ref{tab:atmosphere_symbols}
collects the notation used throughout the paper.
\end{EN}

\begin{physbox}{为什么光谱合成天然适合 GPU}
辐射转移只在同一波长上耦合不同的大气深度（光子在某一频率上穿过各层被吸收和再发射），但一个波长的计算几乎不影响另一个波长。于是“一百万个波长”就是一百万份互相独立的短计算——这正是 GPU 最擅长的大规模并行模式。原 Kurucz 程序诞生于内存极小的年代，被迫逐深度写盘、逐波长串行求解；本文只是换了执行方式，让同一个物理方程以大批量数组运算的形式在 GPU 上完成。物理内容一行未改，验证残差在 $10^{-3}$ 量级以下。
\end{physbox}

\begin{notebox}{表~\ref{tab:atmosphere_symbols} 说明}
下表为原文 Table 1 的中英对照精简版，按用途分组列出全文符号；定义栏先中文后英文（英文有所节略）。
\end{notebox}

\begingroup
\footnotesize
\begin{longtable}{@{}p{0.16\textwidth}p{0.28\textwidth}p{0.48\textwidth}@{}}
\caption{全文符号与定义（中英对照）。Symbols and definitions used throughout the paper.}
\label{tab:atmosphere_symbols}\\
\toprule
\textbf{角色 Role} & \textbf{符号 Symbol} & \textbf{定义 Definition} \\
\midrule
\endfirsthead
\toprule
\textbf{角色 Role} & \textbf{符号 Symbol} & \textbf{定义 Definition} \\
\midrule
\endhead
恒星标签 & $T_{\rm eff}$, $\log g$, $g$, $\xi$, $[\mathrm{M}/\mathrm{H}]$, $[\alpha/\mathrm{M}]$, $[\mathrm{Fe}/\mathrm{H}]$, $\{[\mathrm{X}/\mathrm{H}]\}$, $\{[\mathrm{X}/\mathrm{Fe}]\}$, $A(\mathrm{X})$ & 有效温度、表面重力标签、物理加速度 $g=10^{\log g}$（cgs）、微观湍流与丰度。$[\mathrm{M}/\mathrm{H}]$ 整体移动所有金属，$[\alpha/\mathrm{M}]$ 额外叠加于 O、Ne、Mg、Si、S、Ca、Ti。$[\mathrm{X}/\mathrm{H}]\equiv\log_{10}(n_{\rm X}/n_{\rm H})-\log_{10}(n_{\rm X}/n_{\rm H})_\odot$；$[\mathrm{X}/\mathrm{Fe}]=[\mathrm{X}/\mathrm{H}]-[\mathrm{Fe}/\mathrm{H}]$；$A(\mathrm{X})=\log_{10}(n_{\rm X}/n_{\rm H})+12$。显式元素坐标优先于整体或分组缩放。Effective temperature, gravity, microturbulence and abundances; an explicit element coordinate overrides bulk scaling. \\
深度坐标 & $m$, $\tau_{\rm R}$ & 柱质量向内增加；罗斯兰光学深度满足 $\mathrm{d}\tau_{\rm R}=\kappa_{\rm R}\,\mathrm{d}m$。Column mass and Rosseland optical depth. \\
大气状态 & $T$, $\rho$, $P_{\rm gas}$, $n_e$, $\kappa_{\rm R}$, $g_{\rm rad}$, $\xi$ & 各深度层的温度、质量密度、气体压强、电子密度、罗斯兰不透明度与辐射加速度，外加与深度无关的微观湍流速度。 \\
压强项 & $P_{\rm rad}$, $P_0$ & 辐射压强，以及无湍流压强分支中 $P_{\rm gas}+P_{\rm rad}=gm+P_0$ 的表面常数。 \\
粒子布居 & $n_{s,r}$, $U_{s,r}$, $n_{\rm mol}$, $n_{\rm pert}$, $\mathcal{I}_{s,r}$, $\Delta\mathcal{I}_{s,r}$ & 物种 $s$ 第 $r$ 电离级的数密度与配分函数、分子数密度、中性碰撞代理量、电离能及其随密度变化的下降修正。 \\
能量平衡 & $F_{\rm rad}$, $F_{\rm conv}$, $\sigma_{\rm SB}$, $\epsilon_F$ & 辐射流量与对流流量满足 $F_{\rm rad}+F_{\rm conv}=\sigma_{\rm SB}T_{\rm eff}^4$；$\epsilon_F$ 为相对残差。 \\
迭代 & $n$, $j$, $\Delta T^{(n)}$, $\epsilon_T$ & 迭代序号、深度序号、温度改正量，以及深层最大相对温度变化。 \\
不动点 & $\mathbf{x}^{(n)}$, $\mathbf{p}^{(n)}$, $\mathcal{G}$, $\mathbf{J}_{\mathcal{G}}$, $\mathcal{R}$, $\mathbf{x}_\star$, $\mathbf{e}^{(n)}$ & 数值状态及其六廓线投影、一个大气循环、局部雅可比与残差、收敛不动点与迭代误差。 \\
隐变量初始化器 & $\mathbf{u}$, $\overline{\mathbf u}$, $\mathbf{s}_u$, $\widehat{\mathbf z}$, $\widehat{\mathbf c}$, $\widehat{\mathbf u}$, $\mathbf{B}_K$, $K$, $\mathcal{D}$, $\Delta m_j$, $T_{\rm grey}$, $s_g$ & 变换后的廓线及其训练均值与标度、标准化网络预测、反标准化 PCA 系数、解码廓线、基与维度、解码器，以及质量、温度与加速度变换。 \\
训练损失 & $\mathcal{L}$, $\mathcal{L}_{\rm prof}$, $\mathcal{L}_{\nabla}$, $\mathcal{L}_{\tau}$, $\mathcal{L}_{\rm hse}$, $\lambda$ & 初始化器总目标及其解码廓线、深度导数、光学深度与流体静力学各项及其权重。 \\
约化求解器 & $T(\tau_{\rm R})$, $m(\tau_{\rm R})$, $\mathcal{R}_T$, $\mathcal{R}_m$ & 独立的温度与柱质量廓线及其能量平衡、光学深度残差。 \\
常数 & $c$, $h$, $k_{\rm B}$, $e$, $m_e$ & 光速、普朗克与玻尔兹曼常数、元电荷、电子质量。 \\
原子谱线 & $l$, $\nu_l$, $(gf)_l$, $E_l$, $m_l$ & 跃迁序号、频率、统计权重振子强度、下能级能量、吸收体质量。 \\
致宽 & $\gamma_{\rm rad}$, $\gamma_{\rm S}$, $\gamma_{\rm vdW}$ & 辐射阻尼率，以及每个扰动粒子的斯塔克、范德瓦尔斯系数（分别乘电子与中性碰撞密度）。 \\
谱线定标 & $q(l)$, $\beta$, $\delta_{q,gf}$, $\delta_{q,\beta}$ & 谱线分量 $l$ 所属的改正组、阻尼族序号、对振子强度与阻尼的有界改正；撇号表示定标后的值。 \\
采样 & $\nu$, $\lambda_{i_\lambda}$, $\lambda_0$, $\lambda_{\min}$, $\lambda_{\max}$, $R_{\rm grid}$, $N_\lambda$ & 频率、波长坐标、首个采样与端点。合成网格满足 $\lambda_{i_\lambda+1}/\lambda_{i_\lambda}=1+R_{\rm grid}^{-1}$，共 $N_\lambda$ 个采样。 \\
轮廓 & $\Delta v_{\rm D}$, $\Delta\nu_{\rm D}$, $a_l$, $\phi_l$ & 多普勒速度宽度与频率宽度、阻尼比，以及按 $\int\phi_l(\nu)\,\mathrm{d}\nu=1$ 归一的谱线轮廓。 \\
不透明度 & $\kappa_{\nu,{\rm c}}$, $\sigma_{\nu,{\rm c}}$, $\kappa_{\nu,l}$, $\chi_\nu$ & 连续吸收、连续散射与谱线吸收合成 $\chi_\nu=\kappa_{\nu,{\rm c}}+\sigma_{\nu,{\rm c}}+\sum_l\kappa_{\nu,l}$。 \\
谱线筛选 & $b$, $\nu_b$, $\eta_{\rm sel}$, $C_{jb}$, $Q_{srb}$, $C_0$, $R_l$, $T_*$ & 频率 bin 序号与参考频率、筛选比例、连续谱阈值、布居/连续谱代理最大值、固定常数与中心线强代理；$T_*$ 为最深层温度。 \\
辐射转移 & $\tau_\nu$, $\epsilon_\nu$, $\mu$, $B_\nu$, $J_\nu$, $S_\nu$, $I_\nu$ & 光学深度、热吸收占比、外向方向余弦、普朗克函数、平均强度、源函数与比强度。 \\
输出 & $H_\nu$, $H_\lambda$, $F_\nu$, $F_\lambda$, $F_{\lambda,{\rm c}}$, $f_{\rm norm}$ & $H$ 为一阶角矩，$F=4\pi H$。表面流量即光谱，$F_{\lambda,{\rm c}}$ 为其连续谱，$f_{\rm norm}$ 为二者之比。 \\
仪器 & $\mathcal{K}$ & 从内禀总流量与连续流量到观测网格的线扩展函数线性算子。 \\
光谱拟合 & $\boldsymbol{\theta}$, $y_p$, $w_p$, $f_p$, $f_{\rm fast}$, $f_{\rm phys}$, $\Delta f$, $Q$, $\delta Q_{\min}$ & 拟合标签、观测归一化流量、逆方差权重、像素 $p$ 处模型流量。$f_{\rm fast}$/$f_{\rm phys}$ 为大气收敛前/后的光谱，$Q$ 为每有效像素的平均卡方。 \\
巡天冗余参数 & $\boldsymbol{\psi}$, $v_r$, $v_b$, $\mathcal{S}$, $\mathcal{B}$ & 联合拟合状态、残余速度、有效高斯致宽、多普勒算子与致宽算子。 \\
\bottomrule
\end{longtable}
\endgroup

\subsection{物理计算 The physical calculation}

我们在对数波长上均匀采样合成网格：
\begin{equation}
\lambda_{i_\lambda} =
\lambda_0\left(1+R_{\rm grid}^{-1}\right)^{i_\lambda} .
\end{equation}

\noindent 参数 $R_{\rm grid}$ 设定的是模型采样密度，而非仪器分辨本领。仪器模型通过 $\mathcal{K}$ 施加线扩展函数，一个分辨率元内可以容纳多个模型采样点。采样数约为
\begin{equation}
N_\lambda \simeq
\frac{\ln(\lambda_{\max}/\lambda_{\min})}
     {\ln(1+R_{\rm grid}^{-1})}
\simeq R_{\rm grid}\ln\left(\frac{\lambda_{\max}}{\lambda_{\min}}\right).
\end{equation}

\begin{EN}
We sample the synthesis grid uniformly in logarithmic wavelength,
\begin{equation}
\lambda_{i_\lambda} =
\lambda_0\left(1+R_{\rm grid}^{-1}\right)^{i_\lambda} .
\end{equation}

\noindent The parameter $R_{\rm grid}$ sets the model sampling, not the instrumental
resolving power. The instrument model applies the line spread function through
$\mathcal{K}$, and several model samples can lie within one resolution element.
The number of samples is approximately
\begin{equation}
N_\lambda \simeq
\frac{\ln(\lambda_{\max}/\lambda_{\min})}
     {\ln(1+R_{\rm grid}^{-1})}
\simeq R_{\rm grid}\ln\left(\frac{\lambda_{\max}}{\lambda_{\min}}\right).
\end{equation}
\end{EN}

在每个波长处，质量消光系数是连续吸收、连续散射与谱线吸收之和：
\begin{equation}
\chi_\nu = \kappa_{\nu,{\rm c}} + \sigma_{\nu,{\rm c}}
          + \sum_l \kappa_{\nu,l} .
\end{equation}

\noindent 连续项包括来自氢、氦和金属的主要束缚—自由与自由—自由过程 \citep{Mihalas1978,Gray2005}。在冷星的光学与近红外连续谱中，H$^-$ 占据主导地位 \citep{John1988}。瑞利散射与电子散射单独处理，因为它们的来源不是纯热辐射。

\begin{EN}
At each wavelength the mass extinction coefficient is the sum of continuous
absorption, continuous scattering, and line absorption,
\begin{equation}
\chi_\nu = \kappa_{\nu,{\rm c}} + \sigma_{\nu,{\rm c}}
          + \sum_l \kappa_{\nu,l} .
\end{equation}

\noindent The continuous terms include the principal bound-free and free-free
processes from hydrogen, helium, and metals \citep{Mihalas1978,Gray2005}.
H$^-$ dominates much of the optical and near-infrared continuum in cool stars
\citep{John1988}. Rayleigh and electron scattering remain separate because
their source is not purely thermal.
\end{EN}

\begin{physbox}{不透明度与 H$^-$ 的故事}
不透明度 $\chi_\nu$ 描述气体在某频率上“挡光”的能力，分三部分：连续吸收（光致电离、自由—自由跃迁）、连续散射（瑞利散射、电子散射）、以及所有谱线吸收之和。其中最有趣的是 H$^-$：一个中性氢原子临时“挂靠”一个额外电子形成负离子，它可以吸收很宽频率范围的光子。在冷星光球里氢几乎不电离，自由电子稀缺且主要来自金属——这解释了前文“金属丰度改变连续谱”的通道：金属 $\rightarrow$ 电子 $\rightarrow$ H$^-$ $\rightarrow$ 连续谱。H$^-$ 是 1938 年才被 Wildt 提出、并被证认为太阳光学连续谱的主要来源，是恒星大气物理中“看似次要、实则主导”的经典案例。
\end{physbox}

对于具有频率归一轮廓的标准原子跃迁，LTE 下单位质量的谱线不透明度为 \citep{Gray2005}
\begin{equation}
\begin{aligned}
\kappa_{\nu,l}={}&
\frac{\pi e^2}{m_e c}\,
\frac{n_{s,r}}{\rho U_{s,r}}\,(gf)_l
\exp\left(-\frac{E_l}{k_{\rm B}T}\right) \\
&{}\times
\left[1-\exp\left(-\frac{h\nu}{k_{\rm B}T}\right)\right]\phi_l(\nu), \\
&\text{with}\quad \int\phi_l(\nu)\,\mathrm{d}\nu=1 .
\end{aligned}
\end{equation}

\noindent 其中 $n_{s,r}/U_{s,r}$ 是物种 $s$ 第 $r$ 电离级经配分函数归一的数密度，$(gf)_l$ 包含下能级统计权重。轮廓宽度与阻尼尺度可概括为
\begin{equation}
\begin{aligned}
\Delta v_{\rm D}^{2} &=
\frac{2k_{\rm B}T}{m_l}+\xi^2, \\
\Delta\nu_{\rm D} &= \frac{\nu_l}{c}\Delta v_{\rm D}, \\
a_l &=
\frac{\gamma_{\rm rad}+\gamma_{\rm S}n_e+
\gamma_{\rm vdW}n_{\rm pert}}{4\pi\Delta\nu_{\rm D}} .
\end{aligned}
\end{equation}

\noindent 这里 $n_{\rm pert}$ 是范德瓦尔斯处方使用的有效中性碰撞密度代理。得到的福格特轮廓决定中心不透明度在波长上的分布。在布居固定时，振子强度缩放频率积分的不透明度，而阻尼把不透明度从线心重新分配到翼部 \citep{Barklem2000,Gray2005}。温度、电子密度、丰度和电离状态既改变吸收体的数量，也改变它们所处的致宽环境。这些直接联系使得恒星标签与原子数据可以在光谱空间中被优化。

\begin{EN}
For a standard atomic transition with a frequency normalized profile, the LTE
line opacity per unit mass is \citep{Gray2005}
\begin{equation}
\begin{aligned}
\kappa_{\nu,l}={}&
\frac{\pi e^2}{m_e c}\,
\frac{n_{s,r}}{\rho U_{s,r}}\,(gf)_l
\exp\left(-\frac{E_l}{k_{\rm B}T}\right) \\
&{}\times
\left[1-\exp\left(-\frac{h\nu}{k_{\rm B}T}\right)\right]\phi_l(\nu), \\
&\text{with}\quad \int\phi_l(\nu)\,\mathrm{d}\nu=1 .
\end{aligned}
\end{equation}

\noindent Here $n_{s,r}/U_{s,r}$ is the partition-normalized number density of
ion stage $r$ of species $s$, and $(gf)_l$ includes the lower-level statistical
weight. The profile width and
damping scale can be summarized by
\begin{equation}
\begin{aligned}
\Delta v_{\rm D}^{2} &=
\frac{2k_{\rm B}T}{m_l}+\xi^2, \\
\Delta\nu_{\rm D} &= \frac{\nu_l}{c}\Delta v_{\rm D}, \\
a_l &=
\frac{\gamma_{\rm rad}+\gamma_{\rm S}n_e+
\gamma_{\rm vdW}n_{\rm pert}}{4\pi\Delta\nu_{\rm D}} .
\end{aligned}
\end{equation}

\noindent Here $n_{\rm pert}$ is the effective neutral-collision density proxy
used by the van der Waals prescription.
The resulting Voigt profile sets how the central opacity is distributed in
wavelength. At fixed population, the oscillator strength scales the
frequency-integrated opacity, while damping redistributes opacity from the
core into the wings \citep{Barklem2000,Gray2005}. Temperature, electron density, abundance,
and ionization change both the number of absorbers and their broadening
environment. These direct connections permit stellar labels and atomic data
to be optimized in spectrum space.
\end{EN}

\begin{physbox}{一条谱线是怎么形成的：振子强度、致宽与福格特轮廓}
上式是全文最核心的物理公式，拆解如下：
\begin{itemize}
\item \textbf{吸收体数量}：$\exp(-E_l/k_{\rm B}T)$ 是玻尔兹曼因子——下能级能量 $E_l$ 越高，能处于该能级的原子越少；配分函数 $U_{s,r}$ 则归一化各能级的布居。温度既改变激发也改变电离，所以同一条谱线的强度随恒星类型剧烈变化，这正是光谱型分类的物理基础。
\item \textbf{振子强度 $(gf)$}：量子力学给出的跃迁“固有强度”，缩放整条谱线的积分吸收。它是本文第 6 节定标的主角——许多谱线的 $gf$ 值理论或实验室误差不小，要靠太阳等基准星反推。
\item \textbf{受激辐射修正}：$[1-\exp(-h\nu/k_{\rm B}T)]$ 是 LTE 下扣掉受激发射的因子。
\item \textbf{轮廓 $\phi_l(\nu)$}：多普勒致宽（热运动 $\propto\sqrt{T/m}$ 加微观湍流 $\xi$）给出高斯线心；阻尼（自然宽度、电子碰撞斯塔克致宽、中性原子碰撞范德瓦尔斯致宽）给出洛伦兹翼部。高斯与洛伦兹的卷积就是\textbf{福格特轮廓}。阻尼参数决定强线翼部的形状，是测量表面重力的重要探针（巨星气压低、阻尼翼弱，矮星反之）。
\end{itemize}
\end{physbox}

标准福格特轮廓并不足以描述所有跃迁。我们实现了氢的精细结构、斯塔克轮廓 \citep{VidalCooperSmith1973}，以及通过占据概率物理把高序谱线并入连续谱的处理 \citep{HummerMihalas1988}。独立的 Kurucz 分支处理 He~I、$^3$He、He~II 以及不对称的 Shore--Fano 自电离轮廓 \citep{Fano1961,Shore1967}。这些专门轮廓遵循 \citet{KuruczAvrett1981}。分子谱线不透明度使用分子平衡给出的布居 \citep{Tsuji1973}，分子质量随谱线数据存储并用于设定其多普勒宽度。

\begin{EN}
The standard Voigt profile is not sufficient for every transition. We
implement hydrogen fine structure, Stark profiles
\citep{VidalCooperSmith1973}, and the merging of high series members into the
continuum through occupation-probability physics \citep{HummerMihalas1988}.
Separate Kurucz branches treat He~I, $^3$He, He~II, and asymmetric
Shore--Fano autoionizing profiles \citep{Fano1961,Shore1967}. The specialized
profiles follow \citet{KuruczAvrett1981}.
Molecular line opacity uses populations from molecular equilibrium
\citep{Tsuji1973} and molecular masses stored with the line data
to set their Doppler widths.
\end{EN}

消光通过 $\mathrm{d}\tau_\nu=\chi_\nu\,\mathrm{d}m$ 定义光学深度，其中 $m$ 为柱质量。LTE 下热辐射源是普朗克函数，谱线不透明度按热吸收处理；连续散射的源则依赖平均强度 \citep{Mihalas1978,Sobeck2011}。因此源函数可写为
\begin{equation}
\epsilon_\nu =
\frac{\kappa_{\nu,{\rm c}}+\sum_l\kappa_{\nu,l}}{\chi_\nu},
\qquad
S_\nu = \epsilon_\nu B_\nu(T)
      + \left(1-\epsilon_\nu\right)J_\nu .
\end{equation}

对每个波长，把热源与散射占比映射到光学深度之后，我们求解
\begin{equation}
\mu\frac{\mathrm{d}I_\nu}{\mathrm{d}\tau_\nu}=I_\nu-S_\nu .
\end{equation}

\noindent 转移计算对散射贡献做迭代。由于各波长保持独立，整个波长网格可以并行求解。取 $\mu$ 相对向外法线计量，出射的爱丁顿流量经下式换算：
\begin{equation}
F_\lambda
=F_\nu\left|\frac{\mathrm{d}\nu}{\mathrm{d}\lambda}\right|
=4\pi H_\nu\frac{c}{\lambda^2}
=4\pi H_\lambda .
\end{equation}
\noindent 报告的绝对流量为每纳米的 $F_\lambda$，其中 $c$ 与 $\lambda$ 以一致的纳米单位取值。我们对每个波长携带两条转移分支：一条包含谱线与连续不透明度，另一条只含连续谱。二者一起求解，给出总流量 $F_\lambda$ 及其物理连续谱 $F_{\lambda,{\rm c}}$。存在仪器算子时，先对两者分别施加仪器算子再取比值：
\begin{equation}
f_{\rm norm} =
\frac{\mathcal{K}F_\lambda}
     {\mathcal{K}F_{\lambda,{\rm c}}} .
\end{equation}

\noindent 只要连续谱在一个分辨率元内有变化，这个顺序就很重要。它还让优化器可以直接比较模型与观测像素，而不必在合成与卷积之间把密集采样的光谱搬回 CPU 内存。以上方程定义了物理计算本身；实现层面改变的只是这些运算的组织与执行方式。

\begin{EN}
The extinction defines optical depth through
$\mathrm{d}\tau_\nu=\chi_\nu\,\mathrm{d}m$, where $m$ is column mass. In LTE the thermal source
is the Planck function and line opacity is treated as thermal absorption. The
continuum scattering source depends on the mean intensity
\citep{Mihalas1978,Sobeck2011}. The source function
can therefore be written as
\begin{equation}
\epsilon_\nu =
\frac{\kappa_{\nu,{\rm c}}+\sum_l\kappa_{\nu,l}}{\chi_\nu},
\qquad
S_\nu = \epsilon_\nu B_\nu(T)
      + \left(1-\epsilon_\nu\right)J_\nu .
\end{equation}

For each wavelength, after mapping the thermal source and scattering fraction
to optical depth, we solve
\begin{equation}
\mu\frac{\mathrm{d}I_\nu}{\mathrm{d}\tau_\nu}=I_\nu-S_\nu .
\end{equation}

\noindent The transfer calculation iterates the scattering contribution. Since
wavelengths remain independent, the full wavelength grid can be solved in
parallel. With $\mu$ measured from the outward surface normal, the emergent
Eddington flux is converted through
\begin{equation}
F_\lambda
=F_\nu\left|\frac{\mathrm{d}\nu}{\mathrm{d}\lambda}\right|
=4\pi H_\nu\frac{c}{\lambda^2}
=4\pi H_\lambda .
\end{equation}
\noindent The reported absolute flux is $F_\lambda$ per nanometer, with $c$
and $\lambda$ evaluated in consistent nanometer units.
We carry two transfer branches for each wavelength. One
contains line and continuum opacity. The other contains only the continuum.
We solve them together to produce the total flux $F_\lambda$ and its physical
continuum $F_{\lambda,{\rm c}}$. When an instrument operator is present, we
apply it to both before taking their ratio,
\begin{equation}
f_{\rm norm} =
\frac{\mathcal{K}F_\lambda}
     {\mathcal{K}F_{\lambda,{\rm c}}} .
\end{equation}

\noindent This order matters whenever the continuum varies across a resolution element.
It also lets an optimizer compare the model directly with observed pixels
without moving a densely sampled spectrum back to CPU memory between synthesis
and convolution. These equations define the physical calculation. The
implementation changes how the same operations are arranged and executed.
\end{EN}

\begin{physbox}{源函数、辐射转移方程与 LTE}
辐射转移方程 $\mu\,\mathrm{d}I_\nu/\mathrm{d}\tau_\nu=I_\nu-S_\nu$ 的含义很直观：一束光穿过气体时，强度的净变化 = 被吸收的（$-I_\nu$）加上新发射的（$+S_\nu$）。全部物理都装在\textbf{源函数} $S_\nu$ 里。在 LTE（局部热动平衡）近似下，热吸收对应的辐射由当地温度的普朗克函数 $B_\nu(T)$ 给出——即假设气体局部就像一个温度为 $T$ 的小黑体，碰撞足够频繁，使能级布居由温度唯一决定，无需逐个求解辐射与能级的耦合（那是 NLTE 计算的昂贵所在）。散射部分则不同：散射光子只是“换了个方向”，其辐射强度取决于各方向的平均强度 $J_\nu$，所以要做迭代。我们看到的光谱，就是各深度辐射沿视线方向累积出来的结果：光学深度 $\tau_\nu\sim1$ 的层次贡献最大——谱线中心不透明度大，其 $\tau=1$ 面更高、更冷，所以谱线表现为连续谱背景上的吸收凹陷。
\end{physbox}

\begin{figure}[t]
\centering
\includegraphics[width=0.98\textwidth]{figure_01.pdf}
\caption{与原 Kurucz 计算在 300--1000 nm、$R_{\rm grid}=300{,}000$ 下的一致性。各行依次为热矮星、太阳、红巨星与 K 矮星；各列依次为完整归一化光谱、一个代表性特征，以及 $\Delta f_{\rm norm}$ 残差。黑色与橙色比较两套完整计算；残差面板中蓝色固定大气以单独检验合成，橙色包含在相同标签下独立收敛的大气。残差乘以 $10^3$。自上而下，仅合成与完整计算的 99 百分位绝对残差（单位 $10^{-3}$）分别为 $(2.0,0.7)$、$(2.8,2.2)$、$(5.8,5.8)$ 和 $(3.0,2.9)$。（Agreement with the original Kurucz calculation over 300--1000 nm.）}
\label{fig:synthesis_overview}
\end{figure}

\subsection{从分阶段文件到波长并行执行 From staged files to wavelength-parallel execution}

\begin{figure}[t]
\centering
\includegraphics[width=0.98\textwidth]{figure_02.pdf}
\caption{$R_{\rm grid}=300{,}000$ 下每条光谱的合成运行时间。左图比较四种恒星类型在原始串行 Kurucz 程序与 \paynezero（单张 V100、A100 或 H100 GPU）上的 300--1000 nm 计算。中上、右上面板分解太阳与 K 矮星的 GPU 耗时，下面板给出 1500--1700 nm 太阳光谱的总耗时。原子与分子谱线不透明度主导 GPU 运行时间。（Synthesis runtime per spectrum.）}
\label{fig:synthesis_runtime}
\end{figure}

\begin{figure}[t]
\centering
\includegraphics[width=0.8\textwidth]{figure_03.pdf}
\caption{合成耗时随波长跨度的缩放（$R_{\rm grid}=300{,}000$）。所有窗口从 300 nm 起，止于 400、500、700 或 1000 nm。彩色曲线为 H100 上四类恒星；虚线与阴影带为原 Kurucz 耗时的中位数与范围。更平缓的 H100 缩放使加速比随波长区间与保留跃迁数增大。（Synthesis scaling with wavelength span.）}
\label{fig:synthesis_scaling}
\end{figure}

原始的 Kurucz 计算以一系列 Fortran 程序串行运行。谱线读取程序先准备所请求的波长区间；不透明度阶段随后沿大气深度循环，每一层都重绕谱线记录、累积谱线不透明度并写出按深度排序的记录；直接存取文件把这些记录转置；转移阶段读取转置后的不透明度，逐波长推进，在每个波长点求解连续谱与总量的转移问题。这种组织方式与当年程序编写时的内存条件非常匹配，但也产生了中间输入输出，并让最大的独立维度基本保持串行 \citep{Kurucz2005}。

\begin{EN}
The original Kurucz calculation runs as a sequence of Fortran programs. Line
readers first prepare a requested wavelength interval. The opacity stage then
loops over atmosphere depth, rewinds the line records for each layer,
accumulates line opacity, and writes depth-ordered records. Direct-access files
transpose these records. The transfer stage reads the transposed opacity and
advances through wavelength one point at a time. At every point it solves the
continuum and total transfer problems. This organization was well
matched to the memory available when the programs were written. It also
creates intermediate input and output and leaves the largest independent
dimension largely serial \citep{Kurucz2005}.
\end{EN}

并行执行的机会来自物理依赖关系。辐射转移在一个波长上耦合各层大气，但并不把一个波长耦合到下一个波长。因此自然的工作单元是一大组相互独立的波长列，每列只做一次短深度计算。历史上按深度排序的记录和文件转置掩盖了这一结构；把完整的深度—波长数组保留在内存中，则直接暴露了它。

\begin{EN}
The opportunity for parallel execution follows from the physical dependencies.
Radiative transfer couples the atmosphere depths at one wavelength, but it does
not couple one wavelength to the next. The natural unit of work is therefore a
large set of independent wavelength columns, each with a short calculation
through depth. The historical depth-ordered records and file transpose obscure
this structure. Keeping the complete depth and wavelength array in memory
exposes it directly.
\end{EN}

当一个编译算子同时作用于大量数组元素时，GPU 才能发挥效力，这样的算子常被称为 kernel。在 \paynezero\ 中，我们传给每个 kernel 的是一大组波长或活跃谱线，而非单个波长。物理方程保持不变，但反复的小调用被合并成编译后的数组运算。连续不透明度覆盖整个深度—波长数组，活跃谱线并行地叠加其轮廓，我们对所有波长沿短深度轴积分光学深度。总量与连续谱转移共享同一批次。所有数组和固定的转移几何一直留在 GPU 上，直到 kernel 返回最终流量。

\begin{EN}
A GPU is effective when one compiled operation acts on many array
elements at once. Such an operation is often called a kernel. In \paynezero,
we pass each kernel a large group of wavelengths or active lines rather than
one wavelength. We leave the physical equations unchanged, but repeated
small calls become compiled array operations. Continuum opacity spans the full
depth and wavelength array, active lines add their profiles in parallel, and
we integrate optical depth along the short depth axis for all wavelengths.
Total and continuum transfer share the same batch. We keep the arrays and fixed
transfer geometry on the GPU until the kernel returns the final flux.
\end{EN}

因此，占主导地位的窄谱线与波长计算避免了对单个轮廓或波长采样的 Python 循环。专门的宽线与自电离轮廓以及有界的目录块仍保留少量 Python 层面的调度循环。PyTorch 把大数组运算发往 CUDA 或 CPU kernel，而分子目录由一个缓存的 Numba kernel 准备。测试过的 Apple Metal 路径使合成可以在 Apple silicon 本地运行。所报告的耗时与吞吐建议均基于专用 NVIDIA V100、A100、H100 加速卡。由于 Kurucz 程序本身也是编译代码，速度差异来自并发、更大的 kernel 和常驻输入，而非仅仅来自编译。表~\ref{tab:synthesis_speed} 总结了这些改变如何替代串行循环与中间文件。在反复拟合过程中，固定的窗口、目录与转移几何保持常驻，而依赖大气的不透明度与流量对每个标签向量重新计算。不缓存、不模拟任何合成流量。

\begin{EN}
The dominant narrow-line and wavelength calculations therefore avoid Python
loops over individual profiles or wavelength samples. Specialized broad and
autoionizing profiles and bounded catalog chunks retain small Python-level
dispatch loops. PyTorch sends their large array operations to CUDA or CPU
kernels, while a cached Numba kernel prepares the molecular catalog. The tested Apple Metal path makes synthesis available
locally on Apple silicon. The reported timings and throughput
recommendations use dedicated NVIDIA V100, A100, and H100 accelerators. Since the
Kurucz programs are also compiled, the speed difference comes from concurrency,
larger kernels, and resident inputs rather than compilation alone.
Table~\ref{tab:synthesis_speed} summarizes how these changes replace serial
loops and intermediate files. During repeated fits, the fixed window, catalog,
and transfer geometry remain resident, while atmosphere-dependent opacity and
flux are recomputed for every label vector. No synthetic flux is cached or
emulated.
\end{EN}

\begin{notebox}{表~\ref{tab:synthesis_speed} 说明（原文 Table 2，内容保留英文）}
该表逐阶段对比原 Kurucz 与 \paynezero\ 的执行方式：窗口准备（Numba 编译 + 进程内缓存）、工作布局（PyTorch 批量波长）、谱线不透明度（有界块 + 中心强度测试剔除弱“深度—谱线”对）、内存与 I/O（深度—波长数组常驻设备，取消磁盘中转）、辐射转移（所有波长与两条转移分支一次张量计算）。
\end{notebox}

\begin{table*}[t]
\centering
\caption{Execution of the same synthesis stages in the original Kurucz
workflow and \paynezero. The final column summarizes how resident data and
wavelength-parallel GPU kernels change the computational cost.}
\label{tab:synthesis_speed}
\begingroup
\setlength{\tabcolsep}{4pt}
\renewcommand{\arraystretch}{1.18}
\footnotesize
\begin{tabular}{@{}llll@{}}
\toprule
\tcell{0.12\textwidth}{\textbf{Stage}} &
\tcell{0.25\textwidth}{\textbf{Original Kurucz execution}} &
\tcell{0.28\textwidth}{\textbf{\paynezero\ execution}} &
\tcell{0.25\textwidth}{\textbf{Effect}} \\
\midrule
\tcell{0.12\textwidth}{\textbf{Window preparation}} &
\tcell{0.25\textwidth}{Line readers prepare wavelength-specific sequential
records for the staged programs} &
\tcell{0.28\textwidth}{Numba compiles molecular source records into portable
window-bound arrays. An in-process cache retains the grid, catalogs, tables,
and fixed tensors for one device, dtype, and chunk configuration} &
\tcell{0.25\textwidth}{The retained measurements of repeated calls amortize this
preparation without making the disk products device specific} \\
\tcell{0.12\textwidth}{\textbf{Work layout}} &
\tcell{0.25\textwidth}{Compiled Fortran routines retain serial outer loops over
depth records and wavelengths} &
\tcell{0.28\textwidth}{PyTorch batches independent wavelengths and the dominant
narrow-line and depth contributions. The retained V100, A100, and H100
measurements use CUDA. The same path also supports Metal and CPU execution} &
\tcell{0.25\textwidth}{The bulk work advances concurrently, while Python dispatch is
limited to bounded chunks and specialized profiles} \\
\tcell{0.12\textwidth}{\textbf{Line opacity}} &
\tcell{0.25\textwidth}{Each depth rereads atomic and molecular line records and
advances through profiles in record order} &
\tcell{0.28\textwidth}{Atomic and molecular catalogs are processed in bounded
line chunks. A continuum-referenced center test rejects weak depth--line pairs,
and survivors deposit only their required wing spans} &
\tcell{0.25\textwidth}{Cost and temporary memory follow retained-line scans and
active profile deposits, not the full line-list--wavelength product} \\
\tcell{0.12\textwidth}{\textbf{Memory and I/O}} &
\tcell{0.25\textwidth}{Depth-ordered opacity is written and transposed through
intermediate files} &
\tcell{0.28\textwidth}{Depth--wavelength opacity, source, and transfer arrays
stay on the device between stages, while line and active-pair chunks bound temporary
memory} &
\tcell{0.25\textwidth}{Disk opacity staging and transfers of full arrays through CPU memory are
removed without constructing a depth--line--wavelength array} \\
\tcell{0.12\textwidth}{\textbf{Transfer}} &
\tcell{0.25\textwidth}{Wavelengths advance serially, with separate total and
continuum calls} &
\tcell{0.28\textwidth}{All wavelengths and the stacked total and continuum
branches share one tensor calculation} &
\tcell{0.25\textwidth}{The largest independent dimension is evaluated
together} \\
\addlinespace[2pt]
\bottomrule
\end{tabular}
\endgroup
\end{table*}

\subsection{精度与性能 Accuracy and performance}

我们在 $R_{\rm grid}=300{,}000$ 的对数网格上检验 300--1000 nm 的最终光谱。图~\ref{fig:synthesis_overview} 做了两组比较：共享大气与谱线目录用于单独检验合成精度；在相同恒星标签下各自独立计算则把大气与合成一起检验。两者都以原始 Fortran Kurucz 链条为基准。\footnote{\url{https://github.com/tingyuansen/kurucz}}

\begin{EN}
We test the final spectrum over 300--1000~nm on a logarithmic grid with
$R_{\rm grid}=300{,}000$. Figure~\ref{fig:synthesis_overview} makes two
comparisons. A shared atmosphere and line catalog isolate synthesis accuracy.
Independent calculations at the same stellar labels test atmosphere and
synthesis together. Both use the original Fortran Kurucz chain as the
reference.\footnote{\url{https://github.com/tingyuansen/kurucz}}
\end{EN}

跨四种恒星类型，归一化流量残差的合并均方根在仅合成时为 $9.1\times10^{-4}$，完整计算为 $8.6\times10^{-4}$。对热矮星、太阳、红巨星和 K 矮星，99 百分位绝对残差仅合成时分别为 $(2.0,2.8,5.8,3.0)\times10^{-3}$，完整计算时为 $(0.7,2.2,5.8,2.9)\times10^{-3}$。这些残差反映的是同一物理计算的两种实现之间的微小数值差异。

\begin{EN}
Across the four stellar types, the pooled rms normalized-flux residual is
$9.1\times10^{-4}$ for synthesis alone and $8.6\times10^{-4}$ for the complete
calculation. For the hot dwarf, Sun, red giant, and K dwarf, respectively, the
99th-percentile absolute residuals are $(2.0,2.8,5.8,3.0)\times10^{-3}$ for
synthesis alone and $(0.7,2.2,5.8,2.9)\times10^{-3}$ for the complete
calculation. These residuals reflect small numerical differences between two
implementations of the same physical calculation.
\end{EN}

图~\ref{fig:synthesis_runtime} 测量的是一次性源窗口准备之后的重复前向计算。每次计时都包含连续不透明度、原子与分子谱线、总量与连续谱转移以及返回的光谱。在 $R_{\rm grid}=300{,}000$ 的 300--1000 nm 网格上，太阳光谱的中位耗时约为：原 Kurucz 程序 763 秒、单张 V100 为 32 秒、A100 为 29 秒、H100 为 14 秒，对应加速比分别为 24、26 与 54 倍。

\begin{EN}
Figure~\ref{fig:synthesis_runtime} measures repeated forward calculations after
one-time source-window preparation. Each wall includes continuum opacity,
atomic and molecular lines, total and continuum transfer, and the returned
spectrum. On the 300--1000 nm grid at $R_{\rm grid}=300{,}000$, the solar
medians are about 763~s for the original Kurucz programs, 32~s for one V100,
29~s for one A100, and 14~s for one H100. These correspond to speedups of
$24$, $26$, and $54$ on the V100, A100, and H100, respectively.
\end{EN}

对同一太阳大气、较窄的 APOGEE 1500--1700 nm 区间，V100、A100、H100 的中位耗时约为 3.1、3.1、1.4 秒。全部使用 $R_{\rm grid}=300{,}000$，均为任何仪器卷积或重采样之前的原生网格合成耗时。APOGEE 基准只改变波长区间，不改变合成分辨率。

\begin{EN}
For the same solar atmosphere over the narrower APOGEE 1500--1700~nm interval, the V100,
A100, and H100 medians are about 3.1, 3.1, and 1.4~s.
All use $R_{\rm grid}=300{,}000$ and are native-grid synthesis timings before
any instrumental convolution or resampling. The APOGEE benchmark changes only
the wavelength interval, not the synthesis resolution.
\end{EN}

V100 运行使用 V100S 变体与 2000 条谱线的块，A100 与 H100 分别使用 4000 与 8000 条的块。V100 与 A100 耗时接近，是因为在剖面测试的太阳与 K 矮星计算中，原子与 H/He 不透明度占 80--92\%，而两者在这一主导阶段的实测时间相差不到 10\%；H100 把同一阶段约减半。这些轮廓求值与索引累积 kernel 并未使用稠密张量核矩阵乘法。各 GPU 在其可用的主机、PyTorch 与 CUDA 环境中运行，因此这些是实际应用耗时，而非孤立的架构对比。

\begin{EN}
The V100 runs use the V100S variant and 2,000-line chunks, while the A100 and
H100 runs use 4,000- and 8,000-line chunks. The V100 and A100 walls are close
because atomic and H/He opacity accounts for
80--92 percent of the profiled Sun and K-dwarf calculations, and their measured
times for this dominant stage differ by less than 10 percent. The H100 roughly
halves the same stage. These profile-evaluation and indexed-accumulation kernels
do not use dense tensor-core matrix multiplication. The GPUs ran in their
available host, PyTorch, and CUDA environments. These are therefore achieved
application timings rather than an isolated architecture comparison.
\end{EN}

合成耗时对网格采样只有弱依赖。在固定的 300--1000 nm 区间上，把 $R_{\rm grid}$ 提高 30 倍，四颗星的耗时仅增加 3.1--5.0 倍：更密的采样增加了转移列数和每条活跃轮廓内的像素数，而保留目录的扫描与每次调用的固定开销几乎不变。

\begin{EN}
Synthesis time depends only weakly on grid sampling. Over the fixed
300--1000 nm interval, increasing $R_{\rm grid}$ by a
factor of 30 raises the wall by factors of only 3.1--5.0 across the four stars.
Denser sampling adds transfer columns and pixels within each active profile,
while scanning the retained catalog and the fixed per-call floor remain nearly
unchanged.
\end{EN}

图~\ref{fig:synthesis_scaling} 检验的是波长跨度。在固定的 $R_{\rm grid}=300{,}000$ 下，窗口从 100 nm 扩到 400 nm 跨度使保留目录增大 25 倍，而 H100 耗时仅上升 1.5--1.8 倍，原 Kurucz 则上升 16--18 倍。最短窗口上 GPU 增益较小，是因为工作不足以喂饱加速器；更多波长与轮廓进入批次后，固定的单次调用开销被摊薄。此外，保留条数只是不透明度工作的上界：连续谱截断会剔除大多数“深度—谱线”对，幸存轮廓也只访问有界的像素区间。因此 H100 的加速比从 100 nm 跨度的约 3--6 倍增长到 300--1000 nm 的 36--54 倍。

\begin{EN}
Figure~\ref{fig:synthesis_scaling} instead tests wavelength span. At fixed
$R_{\rm grid}=300{,}000$, expanding the window from a 100 to a 400~nm span
increases the retained catalog by a factor of 25. The H100 wall
rises by only factors of 1.5--1.8, whereas the original Kurucz wall rises by
factors of 16--18. The smaller GPU gain for the shortest window reflects less
work to fill the accelerator. A fixed per-call floor is amortized
as more wavelengths and profiles enter the batch. Moreover, the retained count
is only an upper bound on opacity work because the continuum cutoff rejects most
depth--line pairs and surviving profiles visit bounded pixel ranges. The H100
speedup therefore grows from about $3$--$6$ for the 100~nm span to $36$--$54$
over 300--1000~nm.
\end{EN}


\section{模型大气 Model Atmospheres}\label{sec:atmosphere}

大气计算为合成提供深度结构。在给定有效温度、表面重力、化学组成和微观湍流的条件下，它使压强、粒子布居、不透明度、辐射场与对流达到相互自洽。\paynezero\ 用编译后的数组 kernel 实现了此处所用的一维、平面平行 LTE Kurucz 大气计算 \citep{Kurucz2005}。我们未包含可选的 NLTE 或湍流压强分支。

\begin{EN}
The atmosphere calculation supplies the depth structure used by synthesis. At
specified effective temperature, surface gravity, composition, and
microturbulence, it makes the pressure, populations, opacity, radiation field,
and convection mutually consistent. \paynezero\ implements the
one-dimensional, plane-parallel LTE Kurucz atmosphere calculation used here
\citep{Kurucz2005} with compiled array kernels.
We do not include optional NLTE or
turbulent-pressure branches.
\end{EN}

大气结构必须迭代求解，因为其变量构成一个反馈环：温度和压强决定布居，布居决定不透明度，不透明度决定辐射流量与加热；由此产生的流量不平衡又改变下一次迭代所用的温度和深度结构。因此其计算成本既取决于每一趟物理遍历的工作量，也取决于达到自洽结构所需的遍历次数。

\begin{EN}
Atmospheric structure must be solved iteratively because its variables form a
feedback loop. Temperature and pressure set the populations, the populations
set the opacity, and the opacity sets the radiative flux and heating. The
resulting flux imbalance changes the temperature and depth structure used by
the next iteration. Its computational cost is therefore set by both the work
in each physical pass and the number of passes needed to reach a consistent
structure.
\end{EN}

\begin{physbox}{为什么大气必须“迭代”求解}
这是一个“鸡生蛋”问题：要知道每层的温度，需要知道不透明度；要知道不透明度，又需要知道温度、压强和布居。打破循环的办法是\textbf{迭代}：先猜一个结构（例如灰色大气），计算它产生的辐射流量；若各层流量之和不等于应有的 $\sigma_{\rm SB}T_{\rm eff}^4$，说明温度结构不对，据此修正温度，再重复，直到结构不再变化——即到达\textbf{不动点}。这与解方程 $x=\cos x$ 时反复按“取余弦”键的过程同构。迭代总成本 = 每趟成本 × 趟数；本文第 3 节压缩每趟成本（多核并行），第 4 节压缩趟数（学习型初值）。
\end{physbox}

\subsection{物理计算 The physical calculation}

生产网格在罗斯兰光学深度上取 80 层。我们用六条主廓线概括其物理结构：柱质量 $m$、温度 $T$、气体压强 $P_{\rm gas}$、电子密度 $n_e$、罗斯兰不透明度 $\kappa_{\rm R}$ 和辐射加速度 $g_{\rm rad}$。迭代还携带更新这些廓线所需的辐射压与对流状态。微观湍流 $\xi$ 由请求的标签固定。收敛产物包括合成所需的布居与多普勒宽度。

\begin{EN}
The production grid contains 80 layers in Rosseland optical depth. We summarize
its physical structure with six principal profiles. These are column mass $m$,
temperature $T$, gas pressure $P_{\rm gas}$, electron density $n_e$, Rosseland
opacity $\kappa_{\rm R}$, and radiative acceleration $g_{\rm rad}$. The
iteration also carries the radiation-pressure and convection state needed to
update these profiles. Microturbulence $\xi$ is fixed by the requested labels. The
converged product includes the populations and Doppler widths needed by
synthesis.
\end{EN}

流体静力学平衡把压强与上方物质的重量联系起来。此处采用积分的平面平行形式：
\begin{equation}
P_{\rm gas}(m)+P_{\rm rad}(m)
=g m+P_0 .
\end{equation}
\noindent 其中 $P_{\rm rad}$ 为辐射压强，$P_0$ 为表面积分常数。这是本文支持分支的约定，其中湍流压强被关闭。随后，物态方程在每一层联立求解电荷守恒、萨哈电离、玻尔兹曼布居与分子平衡 \citep{Tsuji1973,Mihalas1978,HubenyMihalas2014}。相邻电离级满足
\begin{equation}
\begin{aligned}
\frac{n_{s,r+1}n_e}{n_{s,r}}
&=\frac{2U_{s,r+1}}{U_{s,r}}
\left(\frac{2\pi m_e k_{\rm B}T}{h^2}\right)^{3/2} \\
&\quad\times
\exp\left[-\frac{\mathcal{I}_{s,r}-\Delta\mathcal{I}_{s,r}}
{k_{\rm B}T}\right].
\end{aligned}
\end{equation}
\noindent 其中 $\Delta\mathcal{I}_{s,r}$ 是所采用处方中孤立电离能随密度下降的修正量，配分函数包含相应的占据修正。电荷守恒封闭电子密度，分子质量作用定律封闭分子布居 \citep{Tsuji1973,BarklemCollet2016}。这些计算共同给出不透明度计算所需的密度与布居。

\begin{EN}
Hydrostatic equilibrium relates pressure to the weight of material above each
layer. In the integrated plane parallel form used here,
\begin{equation}
P_{\rm gas}(m)+P_{\rm rad}(m)
=g m+P_0 .
\end{equation}
\noindent Here $P_{\rm rad}$ is the radiation pressure and $P_0$ is the surface
integration constant. This is the convention for the supported
branch, in which turbulent pressure is disabled. The equation of state then solves charge
conservation, Saha ionization, Boltzmann populations, and molecular equilibrium
at each layer \citep{Tsuji1973,Mihalas1978,HubenyMihalas2014}. Adjacent ion stages obey
\begin{equation}
\begin{aligned}
\frac{n_{s,r+1}n_e}{n_{s,r}}
&=\frac{2U_{s,r+1}}{U_{s,r}}
\left(\frac{2\pi m_e k_{\rm B}T}{h^2}\right)^{3/2} \\
&\quad\times
\exp\left[-\frac{\mathcal{I}_{s,r}-\Delta\mathcal{I}_{s,r}}
{k_{\rm B}T}\right].
\end{aligned}
\end{equation}
\noindent Here $\Delta\mathcal{I}_{s,r}$ is the density-dependent lowering of
the isolated ionization energy in the adopted prescription. Its
partition functions include the corresponding occupation correction. Charge
conservation closes the electron density, while
molecular mass action closes the molecular populations
\citep{Tsuji1973,BarklemCollet2016}. Together these
calculations produce the density and populations used by the opacity calculation.
\end{EN}

\begin{physbox}{流体静力学平衡与萨哈方程}
\textbf{流体静力学平衡}即“每一层气体既不塌缩也不被吹走”：向下的引力被气体压强与辐射压强的梯度抵消。积分形式 $P_{\rm gas}+P_{\rm rad}=gm+P_0$ 中，$gm$ 正是上方质量柱的重量——压强就是“头顶气体的重量”。这也解释了为什么 $\log g$ 能从光谱中测出：$\log g$ 越大，同样光学深度处压强越高，压强致宽（阻尼翼）越强。\textbf{萨哈方程}描述电离平衡：温度越高，指数因子越偏向电离态；电子密度 $n_e$ 越高，复合越快、电离度越低——这解释了巨星（低密度）比矮星电离度更高的现象。电荷守恒把各元素的电离与自由电子总数联立成非线性方程组，即“物态方程”求解，是连接微观物理与宏观结构的枢纽。
\end{physbox}

大气求解器用直接不透明度采样来评估谱线覆盖：这一方法在选定的频率上代表海量的原子与分子谱线集合 \citep{JohnsonKrupp1976}。对每一种所请求的丰度混合物，求解器都重新计算物态方程和每一条被代表的原子与分子不透明度 \citep{Kurucz1996,KuruczAtlas12ASCL}。这些内部波长采样决定大气迭代中的辐射加热与流量。

\begin{EN}
The atmosphere solver evaluates line blanketing by direct opacity sampling, a
method that represents large atomic and molecular line ensembles at selected
frequencies \citep{JohnsonKrupp1976}. For each requested abundance mixture,
the solver recomputes the equation of state and every represented atomic and
molecular opacity \citep{Kurucz1996,KuruczAtlas12ASCL}. These internal wavelength samples
determine radiative heating and flux during atmosphere iteration.
\end{EN}

在每个采样波长处，求解器贯穿所有深度层计算连续与谱线消光 \citep{Mihalas1978}。罗斯兰平均给出与扩散能量输运相关的不透明度：
\begin{equation}
\frac{1}{\kappa_{\rm R}}=
\frac{\displaystyle\int_0^\infty
      \chi_\nu^{-1}\frac{\partial B_\nu}{\partial T}\,\mathrm{d}\nu}
     {\displaystyle\int_0^\infty
      \frac{\partial B_\nu}{\partial T}\,\mathrm{d}\nu},
\qquad \mathrm{d}\tau_{\rm R}=\kappa_{\rm R}\,\mathrm{d}m .
\end{equation}
\noindent 其中 $\chi_\nu$ 为总质量消光。在相同采样波长上的辐射转移给出单色流量、辐射压、辐射加速度与加热积分。特别地，辐射场产生的加速度为
\begin{equation}
g_{\rm rad}(m)=\frac{1}{c}\int_0^\infty
\chi_\nu(m)F_\nu(m)\,\mathrm{d}\nu .
\end{equation}
\noindent 在对流活跃的层次，混合长计算提供对流流量 \citep{BohmVitense1958,Mihalas1978,Gustafsson2008}。目标能量平衡及其相对残差为
\begin{equation}
\begin{aligned}
F_{\rm rad}(m)+F_{\rm conv}(m)
&=\sigma_{\rm SB}T_{\rm eff}^{4}, \\
\epsilon_F(m)
&=\frac{F_{\rm rad}(m)+F_{\rm conv}(m)}
        {\sigma_{\rm SB}T_{\rm eff}^{4}}-1 .
\end{aligned}
\end{equation}

\begin{EN}
At each sampling wavelength, the solver evaluates continuous and line
extinction through all depth layers \citep{Mihalas1978}. The Rosseland
mean gives the opacity relevant to diffusive energy transport,
\begin{equation}
\frac{1}{\kappa_{\rm R}}=
\frac{\displaystyle\int_0^\infty
      \chi_\nu^{-1}\frac{\partial B_\nu}{\partial T}\,\mathrm{d}\nu}
     {\displaystyle\int_0^\infty
      \frac{\partial B_\nu}{\partial T}\,\mathrm{d}\nu},
\qquad \mathrm{d}\tau_{\rm R}=\kappa_{\rm R}\,\mathrm{d}m .
\end{equation}
\noindent Here $\chi_\nu$ is the total mass extinction. Radiative transfer at
the same sampling wavelengths supplies the
monochromatic flux, radiation pressure, radiative acceleration, and heating
integrals. In particular, the acceleration from the radiation field is
\begin{equation}
g_{\rm rad}(m)=\frac{1}{c}\int_0^\infty
\chi_\nu(m)F_\nu(m)\,\mathrm{d}\nu .
\end{equation}
\noindent In layers where convection is active, a mixing-length calculation supplies
the convective flux \citep{BohmVitense1958,Mihalas1978,Gustafsson2008}. The target energy balance and its fractional residual are
\begin{equation}
\begin{aligned}
F_{\rm rad}(m)+F_{\rm conv}(m)
&=\sigma_{\rm SB}T_{\rm eff}^{4}, \\
\epsilon_F(m)
&=\frac{F_{\rm rad}(m)+F_{\rm conv}(m)}
        {\sigma_{\rm SB}T_{\rm eff}^{4}}-1 .
\end{aligned}
\end{equation}
\end{EN}

\begin{physbox}{罗斯兰平均、辐射加速度与能量守恒}
\begin{itemize}
\item \textbf{罗斯兰平均}：在大气深处光子走不远就被吸收，辐射输运类似热传导（扩散近似），此时起控制作用的是罗斯兰平均不透明度。它是 $\chi_\nu^{-1}$ 的加权调和平均——意味着\textbf{不透明度小的“窗口”频率主导能量外流}，因为能量总从最容易漏出去的地方漏。
\item \textbf{辐射加速度 $g_{\rm rad}$}：光子被吸收时把动量交给气体，产生向外的推力。热星中它可接近引力量级（与星风相关）；冷星中很小。它进入流体静力学平衡，等效于减小有效重力。
\item \textbf{能量守恒判据}：稳态大气的每一层，辐射流量加对流流量必须等于 $\sigma_{\rm SB}T_{\rm eff}^4$（斯特藩—玻尔兹曼定律）。残差 $\epsilon_F$ 就是温度修正的依据：某层流量不够，就提高该层温度。
\item \textbf{混合长对流}：当辐射温度梯度超过绝热梯度时，气体元像沸水一样上升下沉输运能量。“混合长”理论用一个自由参数（元胞瓦解前走的距离）近似这一三维过程，是冷星大气模型的重要不确定来源。
\end{itemize}
\end{physbox}

令 $\mathbf{x}^{(n)}$ 表示进入第 $n$ 趟的完整重映射大气状态。其主廓线为 $m$、$T$、$P_{\rm gas}$、$n_e$、$\kappa_{\rm R}$、$g_{\rm rad}$，并携带下一趟所需的压强与对流量。一趟物理求解迭代的示意依赖关系为
\begin{equation}
\begin{aligned}
\mathbf{x}^{(n)}
&\longrightarrow
\left(P_{\rm gas},n_e,\{n_{s,r}\},\{n_{\rm mol}\}\right)^{(n)} \\
&\longrightarrow
\chi_\nu^{(n)} \longrightarrow
\left(J_\nu,F_{\rm rad},g_{\rm rad},\kappa_{\rm R},\tau_{\rm R}\right)^{(n)} \\
&\longrightarrow \left(F_{\rm conv},\Delta T\right)^{(n)}
\longrightarrow \left(T,m\right)^{(n+1)} \\
&\xrightarrow{\ \rm remap\ } \mathbf{x}^{(n+1)},
\qquad T^{(n+1)}=T^{(n)}+\Delta T^{(n)} .
\end{aligned}
\end{equation}
\noindent 即：流体静力学平衡与物态方程先确定压强与布居，随后不透明度与转移计算确定流量、辐射加速度与温度修正。修正同时更新柱质量。两条被修正的廓线连同其他携带字段一起被重映射到标准罗斯兰网格，下一趟再从该状态重算压强与所有布居。

\begin{EN}
Let $\mathbf{x}^{(n)}$ denote the full remapped atmosphere state
entering pass $n$. Its principal profiles are $m$, $T$, $P_{\rm gas}$, $n_e$,
$\kappa_{\rm R}$, and $g_{\rm rad}$. It also carries the pressure and convection
quantities needed by the next pass. One physical solver iteration then has the
schematic dependence
\begin{equation}
\begin{aligned}
\mathbf{x}^{(n)}
&\longrightarrow
\left(P_{\rm gas},n_e,\{n_{s,r}\},\{n_{\rm mol}\}\right)^{(n)} \\
&\longrightarrow
\chi_\nu^{(n)} \longrightarrow
\left(J_\nu,F_{\rm rad},g_{\rm rad},\kappa_{\rm R},\tau_{\rm R}\right)^{(n)} \\
&\longrightarrow \left(F_{\rm conv},\Delta T\right)^{(n)}
\longrightarrow \left(T,m\right)^{(n+1)} \\
&\xrightarrow{\ \rm remap\ } \mathbf{x}^{(n+1)},
\qquad T^{(n+1)}=T^{(n)}+\Delta T^{(n)} .
\end{aligned}
\end{equation}
\noindent Thus hydrostatic balance and the equation of state set pressure and
populations before the opacity and transfer calculation determines the fluxes,
radiative acceleration, and temperature correction. The correction also
updates column mass. Both corrected profiles, together with the other carried
fields, are remapped to the standard Rosseland grid. The next pass recomputes
pressure and every population from this remapped state.
\end{EN}

我们把保留的 Kurucz 工作流程中外部检查器所用的深层结构相继差值范数，改造为逐趟的生产停止判据：
\begin{equation}
\epsilon_T=\max_{j\in {\rm deep}}
\left|\frac{T_j^{(n+1)}-T_j^{(n)}}{T_j^{(n+1)}}\right|
<5\times10^{-4} .
\end{equation}
\noindent 其中 $ {\rm deep}=\{40,\ldots,75\}$ 对应 80 层生产网格。该范数检验深层结构是否已基本停止变化。由于它是局部判据，我们另行单独检查流量平衡与出射光谱。

\begin{EN}
We adapt the deep-layer successive-structure norm used by the external checker
in the retained Kurucz workflow into a per-pass production stopping test,
\begin{equation}
\label{eq:atmosphere_stop}
\epsilon_T=\max_{j\in {\rm deep}}
\left|\frac{T_j^{(n+1)}-T_j^{(n)}}{T_j^{(n+1)}}\right|
<5\times10^{-4} .
\end{equation}
\noindent Here $ {\rm deep}=\{40,\ldots,75\}$ on the 80-layer production grid.
This norm tests whether the deep structure has ceased changing
appreciably. Because it is local, we check flux balance and the emergent
spectrum separately.
\end{EN}

\begin{figure}[t]
\centering
\includegraphics[width=0.98\textwidth]{figure_05.pdf}
\caption{AMD EPYC 9B45 CPU 上的大气性能。左图比较四类恒星在匹配的三趟运行中的每趟平均耗时（\paynezero\ 用 16 线程）；中图为 1 至 32 线程的每趟耗时；右图把红巨星的一趟分解为各物理阶段，对比原 Kurucz 与 16 线程 \paynezero。后者在四颗星上快 6--7 倍，谱线不透明度是剩余最大开销。（Atmosphere performance on an AMD EPYC 9B45 CPU.）}
\label{fig:atmosphere_performance}
\end{figure}

\subsection{从串行迭代到 CPU 并行执行 From serial iteration to parallel CPU execution}

原 Kurucz 大气程序以串行外层迭代执行这一物理循环。每趟迭代中，它先沿深度层循环更新压强与布居、构建谱线不透明度，然后逐个推进 30,000 个不透明度采样波长：在每个波长处计算连续不透明度、贯穿所有层求解辐射转移，并把结果累加进罗斯兰、辐射压与温度修正积分。只有完成这一频率循环之后，才能计算对流和下一个温度结构。

\begin{EN}
The original Kurucz atmosphere program carries this physical cycle through a serial outer
iteration. Within each iteration it loops through the depth layers to update
pressure and populations, constructs line opacity, and then advances through
the 30,000 opacity sampling wavelengths. At each wavelength it evaluates the
continuous opacity, solves radiative transfer through all layers, and adds the
result to the Rosseland, radiation pressure, and temperature correction
integrals. Only after this frequency loop is complete can convection and the
next temperature structure be calculated.
\end{EN}

相继迭代之间无法重叠，因为第 $n$ 趟的温度修正会改变第 $n+1$ 趟的布居与不透明度。但在一趟之内，沿深度、物种、跃迁和不透明度采样波长的工作在归约之前是相互独立的。\paynezero\ 并行化这些块、复用固定输入，并去除重复的解析与中转。

\begin{EN}
Successive iterations cannot overlap because the temperature correction from
pass $n$ changes the populations and opacity of pass $n+1$. Within one pass,
however, work over depth, species, transitions, and opacity-sampling wavelengths
is independent until the partial results are reduced. \paynezero\ parallelizes
these blocks, reuses fixed inputs, and removes repeated parsing and staging.
\end{EN}

\paynezero\ 保持外层大气循环的物理顺序与修正方程不变。在所保留的原程序基准中，求解器精确执行所请求的迭代次数，每批 30 趟；再由单独的检查器比较最后两个结构，必要时启动下一批。\paynezero\ 改为在三趟最低趟数之后，每趟都评估公式~\ref{eq:atmosphere_stop}，到达阈值或配置上限即停止。

\begin{EN}
\paynezero\ keeps the physical ordering and correction equations of the outer
atmosphere cycle. In the retained original-program benchmark, the solver executes
the requested number of iterations exactly. Each batch contains 30 iterations.
A separate checker then compares the last two structures and launches another
batch when needed. \paynezero\ instead evaluates
Equation~\ref{eq:atmosphere_stop} after every pass following a three-pass
minimum, stopping at the threshold or the configured cap.
\end{EN}

在每趟内部，性能关键的循环用 Python 编写。Numba 的即时编译把标记为 \texttt{njit} 的函数变成机器码，\texttt{prange} 把独立循环迭代分配给各 CPU 线程。在启用分子的生产路径中，布居与分子平衡 kernel 保持深度有序，而连续频率、谱线不透明度与转移工作使用线程化执行。大多数固定物理表与完整谱线源目录常驻于带类型的进程数组。由于参考求解器已是编译后的 Fortran，仅凭 JIT 编译无法解释提速；更大的 kernel、常驻输入与中间文件中转的消除提供了其余的加速。

\begin{EN}
Within each pass, the performance-critical loops are written in Python.
Numba's just-in-time compiler turns functions marked \texttt{njit} into
machine code, while \texttt{prange} divides independent loop iterations among
CPU threads. In the molecules-enabled production path, population and
molecular-equilibrium kernels remain depth-ordered, while continuum-frequency,
line-opacity, and transfer work use this threaded execution. Most fixed physics tables and the full line-source
catalogs remain in typed process arrays. Since the reference solver is already
compiled Fortran, JIT compilation alone does not explain the gain. Larger
kernels, resident inputs, and removal of intermediate file staging supply the
remaining reduction.
\end{EN}

谱线筛选可以说明这些改变。源目录固定，但为不透明度采样保留的子集依赖大气。保留测试只需要每个物种—频率 bin 内布居与连续不透明度的紧凑摘要，以及激发随温度的依赖。令 $C_{jb}$ 为深度 $j$、频率 bin $b$ 处的连续谱阈值：
\begin{equation}
C_{jb}=\eta_{\rm sel}
\frac{\kappa_{\nu_b,{\rm c},j}+\sigma_{\nu_b,{\rm c},j}}
{1-\exp[-h\nu_b/(k_{\rm B}T_j)]}.
\end{equation}
\noindent 其中 $\eta_{\rm sel}=10^{-3}$ 为固定筛选比例。对物种 $s$ 的第 $r$ 电离级，深度依赖被压缩为
\begin{equation}
Q_{srb} =
\max_j\left[
\frac{(n_{s,r}/U_{s,r})_j}
{\rho_j\,(\Delta\nu_{\rm D}/\nu)_{jsr}\,C_{jb}}
\right] .
\end{equation}
\noindent 在固定单位换算之内，谱线 $l$ 的中心强度代理为
\begin{equation}
R_l=C_0\,\frac{(gf)_l\,Q_{s(l)r(l)b(l)}}{\nu_{b(l)}} .
\end{equation}
\noindent 其中 $s(l)$、$r(l)$、$b(l)$ 为谱线 $l$ 所属的物种、电离级与频率 bin，$C_0$ 包含固定的谱线不透明度常数与单位换算。当 $R_l\geq 1$ 且 $R_l\exp[-E_l/(k_{\rm B}T_*)]\geq 1$（$T_*$ 为最深层温度）时保留该谱线。这一保守测试避免丢弃可能在任何一层起作用的谱线。

\begin{EN}
Line selection illustrates these changes. The source catalog is fixed, but the
subset retained for opacity sampling depends on the atmosphere. The keep test
needs only a compact summary of the population and continuum opacity in each
species and frequency bin, together with the temperature dependence of
excitation. Let $C_{jb}$ be the continuum threshold at depth $j$ and frequency
bin $b$,
\begin{equation}
C_{jb}=\eta_{\rm sel}
\frac{\kappa_{\nu_b,{\rm c},j}+\sigma_{\nu_b,{\rm c},j}}
{1-\exp[-h\nu_b/(k_{\rm B}T_j)]}.
\end{equation}
\noindent Here $\eta_{\rm sel}=10^{-3}$ is the fixed selection fraction. For ion stage
$r$ of species $s$, the depth dependence is reduced to
\begin{equation}
Q_{srb} =
\max_j\left[
\frac{(n_{s,r}/U_{s,r})_j}
{\rho_j\,(\Delta\nu_{\rm D}/\nu)_{jsr}\,C_{jb}}
\right] .
\end{equation}
\noindent Up to fixed unit conversions, the central strength proxy for line $l$ is
\begin{equation}
R_l=C_0\,\frac{(gf)_l\,Q_{s(l)r(l)b(l)}}{\nu_{b(l)}} .
\end{equation}
\noindent Here $s(l)$, $r(l)$, and $b(l)$ are the species, ion stage, and
frequency bin assigned to line $l$, and $C_0$ contains fixed line-opacity
constants and unit conversions.
The line is retained when both $R_l\geq 1$ and
$R_l\exp[-E_l/(k_{\rm B}T_*)]\geq 1$, where $T_*$ is the deepest layer temperature.
This conservative test avoids discarding a line that can matter in any layer.
\end{EN}

编译后的遍历对固定源目录每个大气只做一次这些测试，并为后续迭代保留筛选出的记录。同一模式贯穿整个物理循环：独立工作被编译与并行化，可复用输入保持常驻，物理顺序保持不变。表~\ref{tab:atmosphere_speed} 总结了这些执行变化。

\begin{EN}
Compiled passes apply these tests once per atmosphere to the fixed source
catalog and retain the selected records for subsequent iterations. The same
pattern extends across the physical cycle. Independent work is compiled and
parallelized, reusable inputs remain resident, and the physical sequence is
preserved. Table~\ref{tab:atmosphere_speed} summarizes these execution changes.
\end{EN}

\begin{notebox}{表~\ref{tab:atmosphere_speed} 说明（原文 Table 3，内容保留英文）}
逐阶段对比大气计算：物理输入（固定物理表与谱线目录常驻内存）、谱线筛选（并行保留测试一次完成，取消筛选文件往返）、布居与连续谱（连续谱 kernel 按频率并行，布居保持深度有序）、谱线不透明度（编译并行块累积）、转移（3 万个频率分块并行求解后归约）。
\end{notebox}

\begin{table*}[t]
\centering
\caption{Execution of the same atmosphere stages in the original Kurucz
workflow and \paynezero. The final column summarizes how resident data,
compiled kernels, and multicore frequency parallelism change the computational
cost.}
\label{tab:atmosphere_speed}
\begingroup
\setlength{\tabcolsep}{4pt}
\renewcommand{\arraystretch}{1.18}
\footnotesize
\begin{tabular}{@{}llll@{}}
\toprule
\tcell{0.12\textwidth}{\textbf{Stage}} &
\tcell{0.25\textwidth}{\textbf{Original Kurucz execution}} &
\tcell{0.28\textwidth}{\textbf{\paynezero\ execution}} &
\tcell{0.25\textwidth}{\textbf{Effect}} \\
\midrule
\tcell{0.12\textwidth}{\textbf{Physical inputs}} &
\tcell{0.25\textwidth}{Input decks and line catalogs use fixed-width or
Fortran-unformatted records. Many numerical tables are compiled into the source} &
\tcell{0.28\textwidth}{Most fixed physics tables and the full line-source
catalogs are retained as typed in-process arrays} &
\tcell{0.25\textwidth}{Cached tables and line catalogs avoid repeated parsing,
although isotope and molecular-equilibrium inputs are still reloaded each pass} \\
\tcell{0.12\textwidth}{\textbf{Line selection}} &
\tcell{0.25\textwidth}{A separate first-pass scan writes the
atmosphere-dependent selected-line file} &
\tcell{0.28\textwidth}{Once per atmosphere, fused parallel keep tests over each
catalog family reuse cached columns and bin assignments and return the selected
catalog in memory. Later passes reuse it} &
\tcell{0.25\textwidth}{Selection remains atmosphere dependent, without a
selected-line file round trip} \\
\tcell{0.12\textwidth}{\textbf{Populations and continuum}} &
\tcell{0.25\textwidth}{Population and continuum calculations advance through
serial nested depth, species, and frequency loops} &
\tcell{0.28\textwidth}{Population and molecular-equilibrium kernels are compiled
but depth-ordered, while continuum kernels parallelize over sampling frequency} &
\tcell{0.25\textwidth}{Continuum frequency work uses the configured CPU threads,
whereas population work remains ordered} \\
\tcell{0.12\textwidth}{\textbf{Line opacity}} &
\tcell{0.25\textwidth}{Selected and detailed line records are scanned serially
and deposited through nested depth and wavelength loops} &
\tcell{0.28\textwidth}{Compiled parallel chunks accumulate selected and
detailed transitions on the opacity-sampling grid} &
\tcell{0.25\textwidth}{The dominant line-opacity work uses the configured CPU
threads} \\
\tcell{0.12\textwidth}{\textbf{Transfer}} &
\tcell{0.25\textwidth}{Each of 30,000 wavelengths is solved and integrated
before the next} &
\tcell{0.28\textwidth}{Compiled frequency chunks solve transfer in parallel into
private accumulators, which are reduced after the parallel region} &
\tcell{0.25\textwidth}{Independent frequency work is parallel, while the
outer physical cycle remains ordered} \\
\addlinespace[2pt]
\bottomrule
\end{tabular}
\endgroup
\end{table*}

大气计算提供的 GPU 并行性不如合成：频率并行的块与物态方程、归约、对流、修正和重映射交错排列，这些阶段交换很小的深度状态且必须保持有序，而可用的波长批次远小于宽波段合成。因此生产求解器使用多核 CPU kernel。当内存流量与有序阶段开始主导时，线程扩展趋于饱和，饱和点由处理器决定。

\begin{EN}
The atmosphere offers less GPU parallelism than synthesis. Frequency-parallel
blocks are interleaved with the equation of state, reductions, convection,
correction, and remapping. These stages exchange a small depth state and must
remain ordered, while the available wavelength batch is much smaller than in
broad spectral synthesis. The production solver therefore uses multicore CPU
kernels. Thread scaling saturates when memory traffic and the ordered stages
dominate, with the saturation point set by the processor.
\end{EN}

\subsection{精度与性能 Accuracy and performance}

图~\ref{fig:synthesis_overview} 中的第一组比较固定大气以单独检验合成；第二组比较中，两个大气求解器在相同恒星标签下各自产生结构再做合成，检验大气差异在最终可观测量中是否仍可忽略。对热矮星、太阳、红巨星和 K 矮星，300--1000 nm 上 99 百分位绝对归一化流量差为 $(0.70,\,2.18,\,5.76,\,2.90)\times10^{-3}$。因此在所测试的参数区间内，大气加合成的完整计算与原 Kurucz 程序保持实用的光谱一致性。

\begin{EN}
The first comparison in Figure~\ref{fig:synthesis_overview} holds the atmosphere
fixed to isolate synthesis. In the second, each atmosphere solver produces its
own structure at the same stellar labels before synthesis, testing whether
atmosphere differences remain negligible in the final observable. For the hot
dwarf, Sun, red giant, and K dwarf over 300--1000~nm,
the 99th-percentile absolute normalized-flux differences are
$(0.70,\,2.18,\,5.76,\,2.90)\times10^{-3}$. The combined atmosphere and synthesis
calculation is thus in practical spectral parity with the original Kurucz
programs across the tested regimes.
\end{EN}

为了把一趟物理遍历的成本与初始化、停止策略分离开，两种实现从同一结构出发，在同一颗 AMD EPYC 9B45 CPU 上以完全相同的物理输入执行三趟迭代。图~\ref{fig:atmosphere_performance} 总结了匹配的单趟耗时、线程扩展与红巨星阶段分解。该设计测量的是一趟物理遍历的执行，而非初始状态或停止策略的选择。

\begin{EN}
To isolate the cost of one physical pass from initialization and stopping, both
implementations begin from the same structure and execute three iterations with
identical physical inputs on the same AMD EPYC 9B45 CPU.
Figure~\ref{fig:atmosphere_performance} summarizes the matched iteration walls,
thread scaling, and red-giant stage breakdown. This design measures the
execution of one physical pass rather than the choice of starting state or
stopping policy.
\end{EN}

对热矮星、太阳、红巨星、K 矮星，原 Kurucz 每趟耗时分别为 15、13、34、38 秒；对应的 16 线程 \paynezero\ 耗时为 2.1、2.2、4.7、5.3 秒，加速 6--7 倍。原程序没有线程并行路径，因此其单核测量是正常运行方式而非外加限制。\paynezero\ 的扩展在约 16 线程处饱和，因为内存流量与有序阶段开始主导。在所测配置中，CPU 推进有序的大气计算，GPU 处理宽得多的合成并行。红巨星的分解显示，原程序的耗时主要由目录输入与谱线筛选主导，而 \paynezero\ 剩余的最大成本是不透明度。因此增益来自常驻数据、并行保留测试与编译不透明度 kernel 在同一 CPU--GPU 工作流中的组合。单趟加速之后，迭代趟数成为剩余的乘性成本。

\begin{EN}
For the hot dwarf, Sun, red giant, and K dwarf, the original Kurucz walls are
15, 13, 34, and 38~s per iteration. The corresponding 16-thread \paynezero\
walls are 2.1, 2.2, 4.7, and 5.3~s, giving factors of 6--7. The original
program has no thread-parallel execution path, so its single-core measurement
is its normal operating mode rather than an imposed restriction. Scaling in
\paynezero\ saturates near 16 threads because memory traffic and ordered stages
begin to dominate. In the measured configuration, the CPUs advance the ordered
atmosphere calculation while the GPU handles the much broader synthesis
parallelism.

The red-giant breakdown shows that the original wall is dominated by catalog
input and line selection, while opacity is the largest remaining \paynezero\
cost. The gain therefore combines resident data, a parallel keep test, and
compiled opacity kernels within one CPU--GPU workflow. Per-iteration
acceleration leaves the number of iterations as the remaining multiplicative
cost.
\end{EN}

\section{学习型大气初始化 Learned Atmosphere Initialization}\label{sec:atmosphere_emulator}

我们通过学习初始状态来削减上述剩余成本，同时保留物理修正循环。一趟循环重建压强与布居、不透明度、辐射转移、对流以及重映射后的温度和柱质量结构。它不构建任何全局导数矩阵，因此每次迭代的成本就是一趟物理遍历。这种经济的更新方式在热矮星、太阳、红巨星和 K 矮星测试中表现良好。在我们的测试中，更通用的求根算法虽然提高了稳健性，但需要多得多的物理遍历。

\begin{EN}
We reduce this remaining cost by learning the starting state while retaining
the physical correction cycle. One cycle
rebuilds pressure and populations, opacity, transfer, convection, and the
remapped temperature and column-mass structure. It forms no global derivative
matrix, so each iteration costs one physical pass. This economical update
works well across the hot-dwarf, solar, red-giant, and K-dwarf tests. In our
tests, more general root solvers improved robustness but required many physical
passes.
\end{EN}

收敛与否仍取决于从初始状态出发的反复更新是否趋近物理解。因此我们用\textbf{物理重启}（而非仅看廓线精度）来评判初始化器。所有比较都保持目标标签与丰度不变，只改变初始大气。预测结果不会被直接当作大气接受：物理求解器必须通过其收敛检查，另有独立检验考察流量闭合与合成光谱。

\begin{EN}
Convergence still depends on whether repeated updates from the starting state
approach the physical solution. We therefore judge an initializer by its
physical restart rather than profile accuracy alone. Every comparison keeps
the target labels and abundances fixed and changes only the starting atmosphere.
The prediction is not accepted as an atmosphere. The physical solver must pass
its convergence checks, while independent tests examine flux closure and the
synthesized spectrum.
\end{EN}

\begin{physbox}{学习型初始化器的定位：只猜初值，不做结论}
这是本文方法论中最“克制”的一处机器学习应用。神经网络不预测光谱、不替代物理，只做一件事：根据恒星标签猜一个足够好的大气初始结构，让物理迭代少跑几趟。最终结果被接受的充要条件仍是物理求解器收敛、流量平衡闭合、光谱检验通过。这与 The Payne 等“标签直接到光谱”的模拟器有本质区别：模拟器的误差会直接进入最终科学产物，而初始化器的误差只会影响收敛速度，不影响收敛到的物理解本身（只要初值落在收敛域内）。这体现了一条重要原则：\textbf{让机器学习加速探索，让物理负责裁决}。
\end{physbox}

\begin{notebox}{表~\ref{tab:initializer_families} 说明（原文 Table 4，内容保留英文）}
三个初始化器家族：\textbf{五标签}（$T_{\rm eff}$、$\log g$、$[\mathrm{M}/\mathrm{H}]$、$[\alpha/\mathrm{M}]$、$\xi$；52,199 个收敛大气）；\textbf{八标签 CNO}（增加 C、N、O 独立丰度；53,824 个）；\textbf{直接丰度}（$T_{\rm eff}$、$\log g$、$\xi$ 加 81 个独立 $[\mathrm{X}/\mathrm{H}]$，内部以 $[\mathrm{Fe}/\mathrm{H}]$ 加 80 个非铁 $[\mathrm{X}/\mathrm{Fe}]$ 表示；82,016 个）。$N_{\rm corpus}$ 为划分训练/验证/测试前保留的独立收敛大气数。
\end{notebox}

\begin{table*}[t]
\centering
\caption{Learned atmosphere initializer families, sample sizes, input labels,
and training envelopes. $N_{\rm corpus}$ is the number of independently
converged atmospheres retained before splitting. The next column gives the
training, validation, and held-out test counts. The public direct-abundance
input contains 81 $[\mathrm{X}/\mathrm{H}]$ values, represented internally by
$[\mathrm{Fe}/\mathrm{H}]$ and 80 non-iron $[\mathrm{X}/\mathrm{Fe}]$
coordinates.}
\label{tab:initializer_families}
\begingroup
\setlength{\tabcolsep}{5pt}
\renewcommand{\arraystretch}{1.20}
\footnotesize
\begin{tabular}{@{}lllll@{}}
\toprule
\tcell{0.11\textwidth}{\textbf{Family}} &
\tcell{0.08\textwidth}{$N_{\rm corpus}$} &
\tcell{0.16\textwidth}{\textbf{Training / validation / test}} &
\tcell{0.27\textwidth}{\textbf{Input labels}} &
\tcell{0.30\textwidth}{\textbf{Training envelope}} \\
\midrule
\tcell{0.11\textwidth}{\textbf{Five-label}} &
\tcell{0.08\textwidth}{52,199} &
\tcell{0.16\textwidth}{50,000 / 2,000 / 199} &
\tcell{0.27\textwidth}{$T_{\rm eff}$, $\log g$, $[\mathrm{M}/\mathrm{H}]$,
$[\alpha/\mathrm{M}]$, $\xi$} &
\tcell{0.30\textwidth}{$4000\leq T_{\rm eff}\leq10500$ K,
$0.7\leq\log g\leq5.3$, $-2.5\leq[\mathrm{M}/\mathrm{H}]\leq0.5$,
$-0.1\leq[\alpha/\mathrm{M}]\leq0.5$,
$0.5\leq\xi\leq4.0$ km s$^{-1}$} \\
\tcell{0.11\textwidth}{\textbf{Eight-label CNO}} &
\tcell{0.08\textwidth}{53,824} &
\tcell{0.16\textwidth}{51,625 / 2,000 / 199} &
\tcell{0.27\textwidth}{$T_{\rm eff}$, $\log g$, $[\mathrm{M}/\mathrm{H}]$,
$[\alpha/\mathrm{M}]$, $\xi$, $[\mathrm{C}/\mathrm{M}]$,
$[\mathrm{N}/\mathrm{M}]$, $[\mathrm{O}/\mathrm{M}]$} &
\tcell{0.30\textwidth}{Five-label bounds, with $-0.5\leq[\mathrm{C}/\mathrm{M}],
[\mathrm{N}/\mathrm{M}],[\mathrm{O}/\mathrm{M}]\leq0.5$} \\
\tcell{0.11\textwidth}{\textbf{Direct abundance}} &
\tcell{0.08\textwidth}{82,016} &
\tcell{0.16\textwidth}{71,660 / 9,733 / 623} &
\tcell{0.27\textwidth}{$T_{\rm eff}$, $\log g$, $\xi$, and 81
$[\mathrm{X}/\mathrm{H}]$ abundances, represented by
$[\mathrm{Fe}/\mathrm{H}]$ and 80 non-iron
$[\mathrm{X}/\mathrm{Fe}]$ coordinates} &
\tcell{0.30\textwidth}{$4000\leq T_{\rm eff}\leq10500$ K,
$0.7\leq\log g\leq5.3$, $0.5\leq\xi\leq4.0$ km s$^{-1}$,
$-2.5\leq[\mathrm{Fe}/\mathrm{H}]\leq0.5$, and
$-0.5\leq[\mathrm{X}/\mathrm{Fe}]\leq0.5$} \\
\addlinespace[2pt]
\bottomrule
\end{tabular}
\endgroup
\end{table*}

\begin{figure}[t]
\centering
\includegraphics[width=0.98\textwidth]{figure_06.pdf}
\caption{五标签与八标签初始化器的精度与迭代成本。左图比较 $K=160$ PCA 压缩下限（虚线、空心）与完整的留出标签到大气误差（实线、实心）；中图给出六个对照在三趟物理求解后的保守光谱误差上限（均低于 $5.1\times10^{-3}$）；右图比较“邻近存档大气”起点与两种学习型起点的迭代次数，横线标记中位数约 12、5.5、4 趟。（Accuracy and iteration cost of the five- and eight-label atmosphere initializers.）}
\label{fig:initializer_convergence}
\end{figure}

\subsection{作为不动点问题的初始化 Initialization as a fixed point problem}

固定所请求的 $T_{\rm eff}$、$\log g$、$\xi$ 与丰度 $\{[\mathrm{X}/\mathrm{H}]\}$，令 $\mathbf{x}^{(n)}$ 为开始第 $n$ 趟迭代所需的完整数值状态。其物理投影包含 80 层生产网格上的六条廓线：
\begin{equation}
\mathbf{p}^{(n)}
=\left\{m_j,T_j,P_{{\rm gas},j},n_{e,j},
\kappa_{{\rm R},j},g_{{\rm rad},j}\right\}_{j=1}^{80}
\end{equation}
\noindent 而完整求解器状态还携带上一次温度修正、累积的罗斯兰不透明度信息与迭代控制量。记 $\mathcal{G}$ 为作用于该增广状态的一个完整 Kurucz 循环：
\begin{equation}
\mathbf{x}^{(n+1)}=\mathcal{G}\left(\mathbf{x}^{(n)}\right).
\end{equation}
\noindent 它返回更新后的状态，其物理投影是同一恒星标签下重映射后的六场大气。收敛的求解器状态即其不动点：
\begin{equation}
\mathbf{x}_\star=\mathcal{G}\left(\mathbf{x}_\star\right).
\end{equation}
\noindent 初始化就是选择 $\mathbf{x}^{(0)}$。在不动点附近，误差 $\mathbf{e}^{(n)}=\mathbf{x}^{(n)}-\mathbf{x}_\star$ 按下式演化：
\begin{equation}
\mathbf{e}^{(n+1)}
\simeq
\mathbf{J}_{\mathcal{G}}(\mathbf{x}_\star)\mathbf{e}^{(n)} .
\end{equation}
\noindent 局部收敛要求 $\mathbf{J}_{\mathcal{G}}$ 的谱半径小于 1，即一趟遍历缩小而非放大小误差。这个条件定义了一个局部收敛域（basin）。它的边界可能很复杂，因为温度与压强改变布居和不透明度，不透明度改变下一步修正，重映射把光学深度与柱质量耦合，对流边界还可能移动。

\begin{EN}
For fixed requested values of $T_{\rm eff}$, $\log g$, $\xi$, and the
abundances $\{[\mathrm{X}/\mathrm{H}]\}$, let $\mathbf{x}^{(n)}$ denote the
complete numerical state needed to begin iteration $n$. Its physical projection
contains the six profiles on the 80-layer production grid,
\begin{equation}
\mathbf{p}^{(n)}
=\left\{m_j,T_j,P_{{\rm gas},j},n_{e,j},
\kappa_{{\rm R},j},g_{{\rm rad},j}\right\}_{j=1}^{80}
\end{equation}
\noindent while the complete solver state also carries the previous temperature
correction, accumulated Rosseland-opacity information, and iteration controls. We use
$\mathcal{G}$ for one complete Kurucz cycle acting on this augmented state,
\begin{equation}
\mathbf{x}^{(n+1)}=\mathcal{G}\left(\mathbf{x}^{(n)}\right).
\end{equation}
\noindent It returns an updated state whose physical projection is the remapped
six-field atmosphere at the same stellar labels. The converged solver state is
its fixed point,
\begin{equation}
\mathbf{x}_\star=\mathcal{G}\left(\mathbf{x}_\star\right).
\end{equation}
\noindent Initialization means choosing $\mathbf{x}^{(0)}$. Near the fixed point, the error
$\mathbf{e}^{(n)}=\mathbf{x}^{(n)}-\mathbf{x}_\star$ evolves as
\begin{equation}
\mathbf{e}^{(n+1)}
\simeq
\mathbf{J}_{\mathcal{G}}(\mathbf{x}_\star)\mathbf{e}^{(n)} .
\end{equation}
\noindent Local convergence requires the spectral radius of
$\mathbf{J}_{\mathcal{G}}$ to be below unity, so one pass reduces rather than
amplifies a small error. This condition defines a local convergence basin.
Its boundary can be complicated because temperature and pressure change the
populations and opacity, opacity changes the next correction, remapping couples
optical depth to column mass, and the convective boundary can move.
\end{EN}

不动点残差
\begin{equation}
\mathcal{R}(\mathbf{x})
=\mathcal{G}(\mathbf{x})-\mathbf{x}
\end{equation}
\noindent 衡量完整携带状态在一趟遍历下是否自洽。起始与收敛的六条廓线很接近，并不保证完整状态在下一趟下变化很小：小的廓线误差仍可能叠加成大的压强、不透明度或对流响应。因此一个有用的初始化器必须落在物理迭代的收敛域之内，而不仅仅是最小化廓线误差。

\begin{EN}
The fixed point residual
\begin{equation}
\mathcal{R}(\mathbf{x})
=\mathcal{G}(\mathbf{x})-\mathbf{x}
\end{equation}
\noindent measures whether the complete carried state is consistent under one
pass. Close agreement between the six starting and converged profiles does not
guarantee that the complete state changes little under the next pass. Small
profile errors can still combine into a large pressure, opacity, or convection
response. A useful initializer must therefore lie inside the basin of the
physical iteration, not merely minimize profile error.
\end{EN}

\begin{physbox}{不动点、雅可比谱半径与收敛域}
把大气迭代写成 $\mathbf{x}^{(n+1)}=\mathcal{G}(\mathbf{x}^{(n)})$ 后，收敛问题就变成标准的\textbf{不动点迭代}理论：解 $\mathbf{x}_\star$ 满足 $\mathcal{G}(\mathbf{x}_\star)=\mathbf{x}_\star$；在解附近线性化，误差每趟乘以雅可比矩阵 $\mathbf{J}_{\mathcal{G}}$。若其\textbf{谱半径}（最大特征值的模）小于 1，误差逐趟衰减——初值就在收敛域内；若大于 1，再接近的初值也可能被弹开。这给“什么样的初值算好”一个精确判据：不是和答案长得像，而是落在收敛域里。后文讨论冷星发散时（第 8 节），物理图像就是小温度扰动会重组分子平衡与对流边界，等效于谱半径超过 1，固定点迭代失去局部稳定性。
\end{physbox}

爱丁顿灰色大气提供了一种通用的初始结构，但它缺乏对谱线覆盖、分子平衡与对流的响应 \citep{Mihalas1978,Gustafsson2008}，因此可能落在有用的收敛域之外。预计算大气网格能提供邻近的收敛起始模型 \citep{CastelliKurucz2003}，但存档模型只有在物理迭代仍处于目标解收敛域内时才有用。

\begin{EN}
An Eddington grey atmosphere provides a generic starting structure, but it
lacks the response to line blanketing, molecular equilibrium, and convection
\citep{Mihalas1978,Gustafsson2008}. It can therefore lie outside the useful
basin. A precomputed atmosphere grid supplies nearby converged starting models
\citep{CastelliKurucz2003}, although a stored model is useful only when the
physical iteration remains in the target solution's basin.
\end{EN}

\begin{physbox}{灰色大气是什么}
“灰色大气”是最简单的解析大气模型：假设不透明度与波长无关（“灰色”），且能量全靠辐射扩散输运，可解出 $T(\tau_{\rm R})=T_{\rm eff}[\frac{3}{4}(\tau_{\rm R}+\frac{2}{3})]^{1/4}$。它抓住了“越深越热”的整体趋势，但真实大气的不透明度强烈依赖波长（谱线覆盖），冷星还有分子与对流。灰色初值与真实结构差得远，物理迭代要跑很多趟甚至失败——这就是需要一个“懂物理的”学习型初值的原因。后文把温度归一化到灰色廓线上再做预测，正是为了让网络只学习相对灰色背景的\textbf{偏离}部分，降低学习难度。
\end{physbox}

训练初始化器首先需要一批收敛的物理大气。更快的物理迭代使这批初始集合可以负担。之后交替进行物理求解与重训练来扩充语料，只有收敛的解才被保留为训练目标。

\begin{EN}
Training an initializer first requires a set of converged physical atmospheres.
Faster physical iteration makes this initial set affordable. Alternating
physical solutions and retraining then expands the corpus, with only converged
solutions retained as training targets.
\end{EN}

\subsection{紧凑的物理大气表示 A compact physical atmosphere representation}

核心设计选择是同时预测全部六个场。它们的原始数值跨越多个数量级，且满足正性与单调性约束。我们把第 $j$ 层变换为
\begin{equation}
\begin{aligned}
\mathbf{u}_j=\bigg(&
\log_{10}\Delta m_j,\,
\log_{10}\frac{T_j}{T_{{\rm grey},j}},\,
\log_{10}P_{{\rm gas},j},\\
&\log_{10}n_{e,j},\,
\log_{10}\kappa_{{\rm R},j},\,
\operatorname{asinh}\frac{g_{{\rm rad},j}}{s_g}
\bigg),
\end{aligned}
\end{equation}
\noindent 其中 $\Delta m_1=m_1$、$j>1$ 时 $\Delta m_j=m_j-m_{j-1}$，且
\begin{equation}
T_{\rm grey}(\tau_{\rm R})=
T_{\rm eff}
\left[\frac{3}{4}\left(\tau_{\rm R}+\frac{2}{3}\right)\right]^{1/4}.
\end{equation}
\noindent 正的质量增量保证单调性，对数保证正性，灰色归一化去掉了大部分整体温度趋势。家族专属的标度 $s_g$ 是训练语料中辐射加速度绝对值的中位数；带符号变换保留了辐射加速度的两个方向。

\begin{EN}
The central design choice is to predict all six fields together. Their raw
values span many orders of magnitude and obey positivity and monotonicity
constraints. We therefore transform layer $j$ to
\begin{equation}
\begin{aligned}
\mathbf{u}_j=\bigg(&
\log_{10}\Delta m_j,\,
\log_{10}\frac{T_j}{T_{{\rm grey},j}},\,
\log_{10}P_{{\rm gas},j},\\
&\log_{10}n_{e,j},\,
\log_{10}\kappa_{{\rm R},j},\,
\operatorname{asinh}\frac{g_{{\rm rad},j}}{s_g}
\bigg),
\end{aligned}
\end{equation}
\noindent where $\Delta m_1=m_1$ and
$\Delta m_j=m_j-m_{j-1}$ for $j>1$, and
\begin{equation}
T_{\rm grey}(\tau_{\rm R})=
T_{\rm eff}
\left[\frac{3}{4}\left(\tau_{\rm R}+\frac{2}{3}\right)\right]^{1/4}.
\end{equation}
\noindent Positive mass increments enforce monotonicity, logarithms preserve
positivity, and the grey scaling removes most of the global temperature trend.
The family-specific scale $s_g$ is the median absolute radiative acceleration
in the training corpus. The signed transform retains both directions of
radiative acceleration.
\end{EN}

变换后的 $80\times6$ 大气共 480 个坐标。联合主成分分析（PCA）\citep{JolliffeCadima2016} 把它压缩为 $K=160$ 个系数，同时保留超过 99.999\% 的标准化训练方差。

\begin{EN}
The transformed $80\times6$ atmosphere contains 480 coordinates. A joint
principal component analysis, or PCA \citep{JolliffeCadima2016}, reduces it to
$K=160$ coefficients while retaining more than 99.999 percent of the
standardized training variance.
\end{EN}

一个前馈网络把表~\ref{tab:initializer_families} 中的五标签、八标签输入映射到标准化的 PCA 系数。对直接丰度家族，一个元素感知的集合编码器先把每种丰度与其元素身份结合 \citep{Zaheer2017}，再把该摘要与 $T_{\rm eff}$、$\log g$、$\xi$（有效温度、表面重力、微观湍流）拼接，然后预测同样的系数表示。若 $\widehat{\mathbf z}$ 为标准化网络输出，解码则逆转系数与廓线的两步标准化：
\begin{equation}
\begin{aligned}
\widehat{\mathbf c}
&=\overline{\mathbf c}+\mathbf{s}_{c}\odot\widehat{\mathbf z},
\\
\operatorname{vec}(\widehat{\mathbf u})
&=\overline{\mathbf u}
+\mathbf{s}_{u}\odot\mathbf{B}_{K}^{\mathsf T}\widehat{\mathbf c},
\\
\mathbf{p}^{(0)}&=\mathcal{D}(\widehat{\mathbf{u}}).
\end{aligned}
\end{equation}
\noindent 其中 $\overline{\mathbf c}$、$\mathbf{s}_c$ 为 PCA 系数的训练均值与标度；对应的廓线量为 $\overline{\mathbf u}$、$\mathbf{s}_u$；$\mathbf{B}_K$ 为保留的 $K$ 分量 PCA 基；帽子标记网络预测；$\operatorname{vec}$ 把深度×场数组拉平；$\odot$ 为逐元素乘；$\mathsf T$ 为转置。解码器 $\mathcal{D}$ 逆转物理变换，返回六条起始廓线 $\mathbf{p}^{(0)}$，求解器再围绕这些廓线构建其余携带状态。联合基保留了跨场、跨深度的相关变化——这很重要，因为物理求解器对六条廓线是整体响应的。

\begin{EN}
A feedforward network maps the five- and eight-label inputs in
Table~\ref{tab:initializer_families} to standardized PCA coefficients. For the
direct-abundance family, an element-aware set encoder first combines each
abundance with its element identity \citep{Zaheer2017}, then joins that summary to
$T_{\rm eff}$, $\log g$, and $\xi$, which denote effective temperature,
surface gravity, and microturbulence, before
predicting the same coefficient representation. If $\widehat{\mathbf z}$ is
the standardized network output, decoding reverses both coefficient and
profile standardization,
\begin{equation}
\begin{aligned}
\widehat{\mathbf c}
&=\overline{\mathbf c}+\mathbf{s}_{c}\odot\widehat{\mathbf z},
\\
\operatorname{vec}(\widehat{\mathbf u})
&=\overline{\mathbf u}
+\mathbf{s}_{u}\odot\mathbf{B}_{K}^{\mathsf T}\widehat{\mathbf c},
\\
\mathbf{p}^{(0)}&=\mathcal{D}(\widehat{\mathbf{u}}).
\end{aligned}
\end{equation}
\noindent Here $\overline{\mathbf c}$ and $\mathbf{s}_c$ are the training mean
and scale of the PCA coefficients. The corresponding profile quantities are
$\overline{\mathbf u}$ and $\mathbf{s}_u$, while $\mathbf{B}_K$ is the retained
$K$-component PCA basis. Hats identify network predictions.
$\operatorname{vec}$ flattens the depth-by-field profile array, $\odot$
denotes elementwise multiplication, and $\mathsf T$ denotes transpose. The
decoder $\mathcal{D}$ reverses the physical transforms and returns the six
starting profiles $\mathbf{p}^{(0)}$. The solver constructs the remaining
carried state around these profiles.
A joint basis preserves correlated variations across fields and depth, which
matters because the physical solver responds to all six profiles together.
\end{EN}

\begin{physbox}{为什么要做这些变换和 PCA}
\begin{itemize}
\item \textbf{变换的目的}：神经网络最擅长拟合量级相近、平滑变化的量。大气各场跨越十几个数量级（如 $n_e$ 从 $10^8$ 到 $10^{18}$），直接预测会被大数主导、小数失真。取对数把乘法变加法；预测 $\Delta m$ 而非 $m$ 自动保证柱质量单调递增；温度除以灰色廓线则把“必然存在的整体趋势”剥离，网络只学物理上有趣的偏离。
\item \textbf{PCA 的作用}：480 个坐标并非独立——相邻层温度相关，温度与压强被流体静力学绑定。PCA 找到这些相关性的主轴，用 160 个系数即可重建 99.999\% 的方差。网络只需预测 160 个数，训练更稳、更省样本。
\item \textbf{物理约束损失}（下式）：损失函数不只看廓线像不像，还检查深度导数、$\mathrm{d}\tau_{\rm R}/\mathrm{d}m=\kappa_{\rm R}$ 和流体静力学平衡——即预测结构要满足大气成立的基本方程。这让初值更可能落在物理迭代的收敛域内。
\end{itemize}
\end{physbox}

我们在 PCA 解码之后评估训练损失，因此它约束的是完整深度结构而非仅系数误差。四项分别对应不同的重启误差：$\mathcal{L}_{\rm prof}$ 匹配六条解码廓线的数值；$\mathcal{L}_{\nabla}$ 匹配相邻层之间的变化，防止整体水平正确但深度依赖错误；$\mathcal{L}_{\tau}$ 检查柱质量与罗斯兰光学深度的映射；$\mathcal{L}_{\rm hse}$ 检查流体静力学压强平衡。系数 $\lambda$ 设定后三项的相对权重：
\begin{equation}
\mathcal{L}
=\mathcal{L}_{\rm prof}
+\lambda_{\nabla}\mathcal{L}_{\nabla}
+\lambda_{\tau}\mathcal{L}_{\tau}
+\lambda_{\rm hse}\mathcal{L}_{\rm hse}.
\end{equation}
\noindent 廓线与梯度项对变换后的场使用稳健的 smooth-L1 失配；光学深度与流体静力学项在以物理单位重建温度、压强、不透明度与柱质量之后求值，衡量对下式的偏离：
\begin{equation}
\frac{\mathrm{d}\tau_{\rm R}}{\mathrm{d}m}=\kappa_{\rm R},
\qquad
\frac{\mathrm{d}P_{\rm gas}}{\mathrm{d}m}\simeq g-g_{\rm rad}.
\end{equation}
\noindent 光学深度项把积分的 $\kappa_{\rm R}\,\mathrm{d}m$ 关系与收敛目标比较；流体静力学项把气压梯度与 $g-g_{\rm rad}$ 比较，并惩罚气压随深度下降的情况。这些选择保留了强烈影响物理重启的廓线形状与两条结构平衡，但不替代辐射与对流平衡——那仍须由物理求解器建立。五标签与八标签模型取 $\lambda_{\nabla}=0.1$、$\lambda_{\tau}=\lambda_{\rm hse}=0.05$；直接丰度模型取 $\lambda_{\nabla}=0.2$，物理一致性项权重相同。

\begin{EN}
We evaluate the training loss after PCA decoding, so it constrains the full
depth structure rather than coefficient errors alone. Four terms address
distinct restart errors. $\mathcal{L}_{\rm prof}$ matches the values of all six
decoded profiles. $\mathcal{L}_{\nabla}$ matches their changes between adjacent
layers, which prevents a profile with the correct overall level from retaining
the wrong depth dependence. $\mathcal{L}_{\tau}$ checks the mapping between
column mass and Rosseland optical depth, while $\mathcal{L}_{\rm hse}$ checks
hydrostatic pressure balance. The coefficients $\lambda$ set the relative
weight of the last three terms,
\begin{equation}
\label{eq:initializer_loss}
\mathcal{L}
=\mathcal{L}_{\rm prof}
+\lambda_{\nabla}\mathcal{L}_{\nabla}
+\lambda_{\tau}\mathcal{L}_{\tau}
+\lambda_{\rm hse}\mathcal{L}_{\rm hse}.
\end{equation}
\noindent The profile and gradient terms use a robust smooth-L1 mismatch on
the transformed fields. The optical-depth and hydrostatic terms are evaluated
after reconstructing temperature, pressure, opacity, and column mass in
physical units. They measure departures from
\begin{equation}
\frac{\mathrm{d}\tau_{\rm R}}{\mathrm{d}m}=\kappa_{\rm R},
\qquad
\frac{\mathrm{d}P_{\rm gas}}{\mathrm{d}m}\simeq g-g_{\rm rad}.
\end{equation}
\noindent The optical-depth term compares the integrated
$\kappa_{\rm R}\,\mathrm{d}m$ relation with the converged target. The
hydrostatic term compares the gas-pressure gradient with $g-g_{\rm rad}$ and
also penalizes decreasing gas pressure. These choices preserve the profile
shape and two structural balances that strongly affect a physical restart.
They do not replace radiative and convective equilibrium, which the physical
solver must still establish. The five- and eight-label
models use $\lambda_{\nabla}=0.1$ and
$\lambda_{\tau}=\lambda_{\rm hse}=0.05$. The direct-abundance model uses
$\lambda_{\nabla}=0.2$ with the same weights on the physical-consistency
terms.
\end{EN}

\begin{figure}[t]
\centering
\includegraphics[width=0.98\textwidth]{figure_04.pdf}
\caption{物理大气结构与学习型初始化。各列分别改变 $T_{\rm eff}$、$\log g$、$[\mathrm{M}/\mathrm{H}]$、$[\alpha/\mathrm{M}]$ 或微观湍流（按色标）；各行为公共罗斯兰光学深度网格上的六条变换后深度廓线。虚线为收敛的物理大气，实线为用于初始化物理求解器的八标签预测。（Physical atmosphere structure and learned initialization.）}
\label{fig:initializer_profiles}
\end{figure}

\subsection{自举、标签覆盖与验证 Bootstrap, label coverage, and validation}\label{sec:atmosphere_emulator_bootstrap}

我们用 Kurucz-A1 初始化器 \citep{Li2025} 的起始猜测播下第一个标签块的种子。只有经物理求解器收敛的大气才进入语料；这些解训练出一个临时初始化器，其预测减少了下一块所需的物理趟数；如此反复求解与重训练，得到表~\ref{tab:initializer_families} 中的语料。预测值从未被用作训练目标。

\begin{EN}
We seeded a first label block with starting guesses from the Kurucz-A1
initializer \citep{Li2025}. Only atmospheres that converged with the physical
solver entered the corpus. Those solutions trained a provisional initializer,
whose predictions reduced the physical passes needed for the next block.
Repeating physical solves and retraining produced the corpora summarized in
Table~\ref{tab:initializer_families}. A prediction never served as a training
target.
\end{EN}

发布的三个家族是实用的默认选择，而非对方法的限制。每个训练大气只需数秒到数十秒，巡天项目完全可以针对自己的标签范围或丰度参数化并行生成数万个。

\begin{EN}
The three released families are practical defaults rather than restrictions on
the method. With each training atmosphere taking seconds to tens of seconds, a
survey can generate tens of thousands in parallel for its own label range or
abundance parameterization.
\end{EN}

五标签家族覆盖以整体金属丰度和 $\alpha$ 模式足以描述大气的广泛恒星工作。拉丁超立方边际抽样 \citep{McKay1979} 在不使用指数级张量网格的情况下对每个标签范围分层。对 $T_{\rm eff}$ 与 $\log g$ 秩次施加的弱相关性让常见的矮星与巨星结构被更密地采样，而不把语料限制在某条演化轨迹上。

\begin{EN}
The five-label family covers broad stellar work in which a bulk metallicity and an
alpha pattern describe the atmosphere adequately. Latin-hypercube
marginals \citep{McKay1979} stratify every label range without an exponentially
large tensor grid. A small imposed correlation between the $T_{\rm eff}$ and
$\log g$ ranks samples common dwarf and giant structures more densely without
restricting the corpus to an evolutionary track.
\end{EN}

八标签家族面向演化恒星的丰度应用：挖掘上翻（dredge-up）使表面 C、N 独立于整体金属与 $\alpha$ 模式而变化 \citep{Salaris2015,Martig2016}。这些变化重新分配 CO、CN、OH 中的原子，可改变分子不透明度与大气结构 \citep{Tsuji1973,Gustafsson2008}。C、N、O 的条件拉丁超立方面板附着在抽样得到的五标签母状态上，每个八标签大气独立求解。在两个低维家族中，$[\alpha/\mathrm{M}]$ 在整体 $[\mathrm{M}/\mathrm{H}]$ 缩放之上叠加于 $\alpha$ 元素；CNO 家族中显式的 $[\mathrm{O}/\mathrm{M}]$ 坐标取代氧的这一缩放，而 $[\alpha/\mathrm{M}]$ 继续设定 Ne、Mg、Si、S、Ca、Ti。

\begin{EN}
The eight-label family addresses abundance applications in evolved stars,
where dredge-up changes surface C and N independently of the bulk metal and
alpha patterns \citep{Salaris2015,Martig2016}. These changes redistribute
atoms among CO, CN, and OH and can alter molecular opacity and atmospheric
structure \citep{Tsuji1973,Gustafsson2008}. Conditional Latin-hypercube
panels in C, N, and O are attached to sampled five-label parent states, and
every resulting eight-label atmosphere is solved independently.
In both lower-dimensional families, $[\alpha/\mathrm{M}]$ adds to the bulk
$[\mathrm{M}/\mathrm{H}]$ scaling for the alpha elements. The explicit
$[\mathrm{O}/\mathrm{M}]$ coordinate replaces that scaling for oxygen in the
CNO family, while $[\alpha/\mathrm{M}]$ continues to set Ne, Mg, Si, S, Ca,
and Ti.
\end{EN}

\begin{physbox}{为什么巨星需要单独的 C、N、O 坐标}
恒星离开主序后会经历对流“挖掘上翻”：对流包层加深，把核心氢燃烧（CNO 循环）的加工产物翻到表面，使 C 减少、N 增多——这颗星的表面 C、N 不再反映其出生时的整体金属丰度，而是携带其质量与演化阶段的信息（C/N 比甚至被用作红巨星的“时钟”）。因此对演化恒星，一个整体 $[\mathrm{M}/\mathrm{H}]$ 加 $[\alpha/\mathrm{M}]$ 的两参数模式无法描述真实大气；C、N、O 必须成为独立坐标。又因为 CO 分子的结合能极高，C 与 O 谁少谁就被锁进 CO，剩下的参与 CN、OH 等分子——所以 C/O 比的细微变化会显著改变分子不透明度，这也是大气模型必须自洽处理分子平衡的原因。
\end{physbox}

直接丰度家族接受 81 个独立变化的 $[\mathrm{X}/\mathrm{H}]$ 丰度，覆盖 Li 到 Bi（除 Tc、Pm）以及 Th、U，没有整体金属丰度或 $\alpha$ 坐标。网络把同一请求混合物重述为 $[\mathrm{Fe}/\mathrm{H}]$ 与 80 个非铁 $[\mathrm{X}/\mathrm{Fe}]$ 值，但这一坐标变换不改变传给物理求解器的丰度。

\begin{EN}
The direct-abundance family accepts 81 independently variable
$[\mathrm{X}/\mathrm{H}]$ abundances. It covers Li through Bi except Tc and Pm,
together with Th and U. It has no bulk
metallicity or alpha coordinate. The network re-expresses the same requested
mixture as $[\mathrm{Fe}/\mathrm{H}]$ and 80 non-iron
$[\mathrm{X}/\mathrm{Fe}]$ values, but this change of coordinates does not alter
the abundances passed to the physical solver.
\end{EN}

我们首先在六个光谱对照上验证五标签与八标签家族：热矮星、太阳、红巨星，以及富碳、富金属、富氧混合物。单元素覆盖值是绝对的 $[\mathrm{X}/\mathrm{H}]$，替换该元素的整体缩放。

\begin{EN}
We first validate the five- and eight-label families on six spectral controls spanning a hot
dwarf, the Sun, a red giant, and C-enhanced, metal-rich, and O-enhanced
mixtures. Individual abundance
overrides are absolute $[\mathrm{X}/\mathrm{H}]$ values and replace the bulk
scaling for that element.
\end{EN}

图~\ref{fig:initializer_convergence} 左图分离了两种误差。空心符号是 $K=160$ PCA 压缩下限，其标准化坐标 RMS 中位数五标签为 $6.3\times10^{-4}$、八标签为 $4.9\times10^{-4}$——即把 480 个廓线坐标压到 160 个系数再解码，丢弃的大气结构很少。实心符号包含从恒星标签经网络与 PCA 解码的完整预测，对应误差为 $1.5\times10^{-2}$ 与 $1.1\times10^{-2}$，比压缩下限大一个数量级以上。因此当前精度由学习到的“标签到大气”关系决定，而非 PCA 压缩。

\begin{EN}
The left panel of Figure~\ref{fig:initializer_convergence} separates two errors.
Open symbols show the $K=160$ PCA compression floor, whose median
standardized-coordinate RMS is $6.3\times10^{-4}$ for the five-label family and
$4.9\times10^{-4}$ for the eight-label family. Thus compressing the 480 profile
coordinates to 160 coefficients and decoding them discards little of the
atmosphere structure. Filled symbols include the complete prediction from
stellar labels through the network and PCA decoder. Their corresponding errors
are $1.5\times10^{-2}$ and $1.1\times10^{-2}$, more than an order of magnitude
larger than the compression floor. The learned label-to-atmosphere relation,
rather than PCA compression, therefore sets the current accuracy.
\end{EN}

中图问：从每个五/八标签预测出发做三趟求解后，还可能残留多少光谱误差？记 $f_3$、$f_5$、$f_{\rm ref}$ 分别为三趟、五趟后与存档参考的光谱，则 $|f_3-f_{\rm ref}|\leq |f_3-f_5|+|f_5-f_{\rm ref}|$。我们报告这一保守上限的最大值：全部六个对照对两个发布的初始化器都低于 $5.1\times10^{-3}$。三趟只是进度诊断，不是停止规则。

\begin{EN}
The middle panel asks how much spectral error can remain after starting from
each five- and eight-label prediction and taking three solver passes. If
$f_3$, $f_5$, and $f_{\rm ref}$ are the spectra after three passes, after five
passes, and from the stored reference, then
$|f_3-f_{\rm ref}|\leq |f_3-f_5|+|f_5-f_{\rm ref}|$. We report the maximum of
this conservative bound. All six controls lie below $5.1\times10^{-3}$ for
both released initializers. Three passes are a progress diagnostic rather than
a stopping rule.
\end{EN}

右图在 16 个匹配的留出案例中单独考察初始大气的作用。为模拟从预计算网格的邻近节点出发，我们按预先声明的偏移（$T_{\rm eff}$ $\pm250$ K、$\log g$ $\pm0.25$ dex、金属丰度 $\pm0.20$ dex）选取附近的收敛结构，仅用存档结构初始化，然后每颗星都在精确目标标签与丰度下求解。

\begin{EN}
The right panel isolates the effect of the starting atmosphere in 16 matched
held-out cases. To emulate starting from a nearby node in a precomputed grid,
we selected converged structures near two sets of predeclared offsets. They are
$\pm250$ K in $T_{\rm eff}$, $\pm0.25$ dex in $\log g$, and $\pm0.20$ dex in
metallicity. Only the stored structure
initializes the calculation. Every atmosphere is then solved at the exact
target labels and abundances.
\end{EN}

两类网格式起点的迭代中位数约 12 趟，而五标签初始化器为 5.5 趟、八标签为 4 趟——学习型起点把迭代次数分别压缩约 2 倍和 3 倍。结合实测的 6--7 倍单趟加速，八标签计算在这些匹配案例中总共和快约 20 倍。所保留的原程序基准每批请求 30 趟迭代；4 趟物理遍历加更快的单趟，相对该基准约有 50 倍增益，而邻近网格大气给出的是上文中更小、更现实的匹配比较。

\begin{EN}
The median of the two grid-like starts is about 12 iterations, compared with
5.5 for the five-label initializer and 4 for the eight-label initializer. The
learned starts therefore reduce the iteration count by factors of about two and
three. Combined with the measured six- to seven-fold per-iteration gain, the
eight-label calculation is about 20 times faster in these matched cases. The
retained original-program benchmark requests 30 iterations per batch. Four
physical passes and the faster iteration give a gain of about 50 relative to
that benchmark, while a nearby grid atmosphere gives the smaller and more realistic
matched comparison above.
\end{EN}

图~\ref{fig:initializer_profiles} 在廓线空间直接检查八标签初始化器：实线预测与虚线收敛物理廓线对比了 25 个大气（其 C、N、O 随整体混合物变化）。每列改变一个恒星标签，六行展示耦合的温度、压强、密度、不透明度与深度响应。该比较检验从标签经网络与 PCA 解码的完整映射。把标准化残差对所有廓线、大气与深度点合并，全深度 RMS 为 0.028，$\log\tau_{\rm R}\geq1$ 处为 0.027。

\begin{EN}
Figure~\ref{fig:initializer_profiles} provides a direct profile-space check of
the eight-label initializer. Solid predictions are compared with dashed
converged physical profiles for 25 atmospheres whose C, N, and O values follow
the corresponding bulk mixture. Each column varies one stellar label while the
six rows trace the coupled temperature, pressure, density, opacity, and depth
responses. This comparison tests the complete mapping from labels through the
network and PCA decoder. Pooling the standardized residuals over all profiles,
atmospheres, and depth points gives an RMS of 0.028 over the full depth range
and 0.027 for $\log\tau_{\rm R}\geq1$.
\end{EN}

对直接丰度家族，623 个测试对象含 313 个内部与 310 个外部大气。内部大气未参与网络拟合，其温度与气压误差的 95 百分位为 0.67\% 与 0.034 dex。另做的 16 案例重启诊断在 8--14 趟物理遍历内收敛，中位数约 9 趟。在下文应用中，该网络提供初始状态，每个被采用的大气随后都经物理求解与验证。

\begin{EN}
For the direct-abundance family, the 623-object test split contains 313
internal and 310 external atmospheres. The internal atmospheres were excluded
from network fitting and have 95th-percentile temperature and gas-pressure
errors of 0.67 percent and 0.034 dex. A separate 16-case restart diagnostic
converges in 8--14 physical passes, with a median near 9. In the application
below, this network supplies the starting state and every retained atmosphere
is then solved and validated physically.
\end{EN}

这个家族的插值问题更难：3 个恒星参数加 81 个独立丰度构成 84 维输入，训练语料对该空间的采样比五/八标签语料稀疏得多。按化学或不透明度响应对丰度分组的坐标，有望在未来初始化器中降低有效维度。

\begin{EN}
This family presents a harder interpolation problem because three stellar
parameters and 81 independently varying abundances define an 84-dimensional
input. Its training corpus samples that space much more sparsely than the
five- and eight-label corpora. Coordinates that group abundances by their
chemical or opacity response may reduce the effective dimension in future
initializers.
\end{EN}

\section{归一化光谱拟合 Normalized Spectrum Fitting}\label{sec:fitter}

如果大气与合成计算足够快，物理前向模型本身就可以放进光谱优化器中求值。直接合成可以反复探索标签；但若每次试探都收敛一个新大气，计算就会被大气求解主导。因此我们用初始化器提供初始深度结构，每次试探都重新计算布居、谱线不透明度与辐射转移，而把大气收敛保留给候选极小点及任何实质性修正。\textbf{拟合中不进入任何“标签到流量”的光谱模拟器。}

\begin{EN}
If atmosphere and synthesis calculations are fast enough, the physical forward
model itself can be evaluated inside a spectrum optimizer. Direct synthesis can
explore labels repeatedly, while converging a new atmosphere at every trial
would dominate the calculation. We therefore use the initializer to supply the
starting depth structure, recompute populations, line opacity, and radiative
transfer at every trial, and reserve atmosphere convergence for the candidate
minimum and any material correction. No label-to-flux spectral emulator enters
the fit.
\end{EN}

我们首先在一个受控的归一化光谱上测试这一策略，其生成模型与拟合模型完全相同。该测试在引入谱线数据与巡天仪器失配之前，先单独检验标签恢复与大气修正。

\begin{EN}
We first test this strategy on a controlled normalized spectrum whose
generating and fitted models are identical. The test isolates label recovery
and atmosphere correction before introducing line-data and survey-instrument
mismatch.
\end{EN}

\begin{physbox}{“快速搜索 + 物理收敛”的两级策略}
这是全文拟合策略的枢纽，值得看清其分工：\textbf{快速阶段}（GPU 上每次约 1 秒）从学习型初值直接合成光谱，在标签空间里做局部最小二乘下山；它用到的“大气”未经完整物理迭代，因此存在小误差。\textbf{裁决阶段}在快速搜索收敛后，于该候选点真正收敛物理大气、重新合成，并量化快速阶段造成的光谱差 $\Delta f$；若修正预测的收益超过阈值 $\delta Q_{\min}$，就再收敛一次大气确认。这样，成千上万次试探都走便宜的快速通道，昂贵的物理收敛只发生一到两次——而最终被接受的解一定来自收敛大气。这是“机器学习加速探索、物理负责裁决”原则在拟合环节的又一次体现。
\end{physbox}

快速搜索从初始标签向量出发。一条直接光谱加每个标签一个有界扰动定义局部光谱响应；加权线性最小二乘提出有界的标签更新，直接合成线搜索只在目标下降时才接受 \citep{NocedalWright2006}。发布的拟合器每次迭代重建该响应；图~\ref{fig:fitter_convergence} 的短诊断在连续接受的步长之间复用初始响应。响应用有限差分获得，因为完整的大气到光谱计算尚未构成单一 PyTorch 计算图；解析响应可以替代这些差分而不改变物理大气计算。其中 $N_{\rm good}$ 为保留像素数，目标 $Q$ 为每有效像素的平均卡方。当没有任何有界线搜索试探能降低 $Q$ 时，快速循环停止。

\begin{EN}
The fast search begins at the initial label vector. One direct spectrum and one
bounded perturbation per label define a local spectral response. A weighted
linear least-squares step proposes a bounded label update, and a direct-synthesis
line search retains it only when the objective improves
\citep{NocedalWright2006}. The released fitter rebuilds this response at each
iteration. The short Figure~\ref{fig:fitter_convergence} diagnostic reuses its
initial response between successive accepted steps. We obtain the response by
finite differences because the complete atmosphere-to-spectrum calculation
does not yet form one PyTorch computational graph. An analytic response could
replace these differences without changing the physical atmosphere calculation.
Here $N_{\rm good}$ is the number of retained pixels. The objective $Q$ is the
mean chi-square per valid pixel.
The fast loop stops when no bounded line-search trial lowers $Q$.
\end{EN}

第一个收敛大气的作用不只是检验该候选点。记 $f_{\rm fast}$ 与 $f_{\rm phys}$ 为快速极小点处大气收敛前、后的光谱。我们把 $\Delta f=f_{\rm phys}-f_{\rm fast}$ 在一个有界局部步长内固定，并为每个活跃标签重建光谱响应；随后用一个高斯—牛顿步求解加权线性最小二乘，给出预计能降低像素残差的标签变化组合。每个活跃标签都参与这一修正；暴露了逐元素丰度坐标的拟合因此可以在第一次物理计算之后调整这些坐标。

\begin{EN}
The first converged atmosphere does more than test this candidate. Let
$f_{\rm fast}$ and $f_{\rm phys}$ denote the spectra before and after physical
atmosphere convergence at the fast minimum. We hold
$\Delta f=f_{\rm phys}-f_{\rm fast}$ fixed over one bounded local step and
rebuild the spectral response for every active label. A Gauss--Newton step then
solves the weighted linear least-squares problem for the combination of label
changes predicted to reduce the pixel residuals. Every active label participates
in this correction. Fits that expose element-by-element abundance coordinates
may therefore adjust them after the first physical calculation.
\end{EN}

这一局部修正只假设 $\Delta f$ 在局部信赖域内变化缓慢。只有当预测的 $Q$ 下降超过 $\delta Q_{\min}$ 时，才把提案送入第二个收敛大气；本受控诊断取 $\delta Q_{\min}=0.005$（每有效像素平均卡方）。只有当该物理光谱真的降低了同一目标时才接受提案。算法~1 总结了此受控诊断使用的单次修正路径；发布的精细化版本只在预测下降超过 $\delta Q_{\min}$ 时重复该计算，并可在物理评估次数上限内继续做有界修正。

\begin{EN}
This local correction assumes only
that $\Delta f$ varies slowly across the local trust region. We send the
proposal to a second converged atmosphere only when its predicted decrease in
$Q$ exceeds $\delta Q_{\min}$. For this controlled diagnostic,
$\delta Q_{\min}=0.005$ in mean chi-square per valid pixel. We accept the
proposal only when that physical spectrum lowers the same objective.
Algorithm~1 summarizes the
single-correction path used for this controlled diagnostic. The released
refinement repeats this calculation only when the predicted decrease in $Q$
exceeds $\delta Q_{\min}$ and can continue bounded corrections up to its cap on
physical evaluations.
\end{EN}

受控案例是一颗留出的非太阳丰度 K 巨星，五个初始标签均有偏移，噪声为 $\mathrm{S/N}=100$ 的独立高斯噪声。它使用 $R_{\rm grid}=20{,}000$ 对数网格上的 1500--1700 nm 归一化光谱，无仪器线扩展函数或其他致宽。拟合最小化每有效像素的平均卡方，并省略观测光谱所需的掩模、连续谱、速度、致宽与变化逆方差。其无噪声真值来自收敛大气加直接合成，因此标签恢复与快速阶段近似都已知。最终物理模型把 $T_{\rm eff}$ 恢复到 2 K 以内，$\log g$ 与两个丰度坐标到 0.003 dex 以内，微观湍流到 0.04 km s$^{-1}$ 以内。

\begin{EN}
The controlled example is a held-out non-solar K giant with offsets in all five
initial labels and independent Gaussian noise at $\mathrm{S/N}=100$. It uses a
normalized 1500--1700~nm spectrum on a logarithmic grid with
$R_{\rm grid}=20{,}000$, with no instrumental line-spread function or other
broadening. The fit minimizes the mean chi-square per valid pixel and omits
the masks, continuum, velocity, broadening, and varying inverse variances needed
for observed spectra. Its noiseless truth comes from a converged atmosphere and
direct synthesis, so both label recovery and the fast-stage approximation are
known. The final physical model recovers $T_{\rm eff}$ within 2~K, $\log g$ and both
abundance coordinates within
0.003~dex, and microturbulence within 0.04~km~s$^{-1}$.
\end{EN}

图~\ref{fig:fitter_convergence} 上图追踪被接受的标签轨迹，下图比较最终物理光谱与模拟观测，残差与注入噪声相符。第一个收敛大气已很接近模拟真值；为演示修正路径，该诊断仍走了一个有界标签步并用第二个收敛大气确认。常规使用中，只有预测下降超过 $\delta Q_{\min}$ 时才走这一步。五标签搜索在单张 H100 上约 20 秒；收敛大气检查在多核 CPU 上运行，占去剩余的主要耗时。这一近乎精确的恢复在已知物理模型上确立了优化器。下一步处理观测光谱中存在的模型—数据失配，即谱线表定标。

\begin{EN}
Figure~\ref{fig:fitter_convergence} traces the accepted label trajectory in
the upper panels and compares the final physical spectrum with the mock
observation in the lower panels, where the residuals match the injected noise.
The first converged atmosphere is already close to the mock truth. To
demonstrate the correction path, this diagnostic nevertheless takes one bounded
label step and confirms it with a second converged atmosphere. In routine use,
the refinement takes such a step only when its predicted decrease in $Q$
exceeds $\delta Q_{\min}$.
The five-label search
takes about 20~s on one H100, while the converged-atmosphere checks run on
multicore CPUs and dominate the remaining wall time. This near-exact recovery
establishes the optimizer on a known physical model. We next address the
model--data mismatch present in observed spectra through line-list calibration.
\end{EN}

\begin{physbox}{卡方、响应矩阵与高斯—牛顿步}
拟合的数学骨架：目标 $Q$ 是加权残差平方和（卡方）除以像素数。在每个当前点，通过对每个标签做小扰动、重新合成光谱，测得“标签 $\rightarrow$ 光谱”的局部线性响应（即导数矩阵/雅可比的列）。把光谱残差投影到这组响应上解线性最小二乘，就得到预计最能降低残差的标签步长——这就是\textbf{高斯—牛顿}思想：用一阶导数近似二阶曲率，避免昂贵的 Hessian 计算。每一步都用真实合成验证目标确实下降（线搜索），步长有界（信赖域），保证不会跳到模型不可信的区域。有限差分响应每维需要一次额外合成——12 维就是 13 条光谱——所以“每次试探都便宜”是该策略可行的前提。
\end{physbox}

\begin{figure}[t]
\centering
\includegraphics[width=0.98\textwidth]{figure_07.pdf}
\caption{对 $\mathrm{S/N}=100$、真值已知的归一化模拟光谱做受控五标签拟合。上图在 $T_{\rm eff}$--$\log g$ 与 $[\mathrm{M}/\mathrm{H}]$--$[\alpha/\mathrm{M}]$ 平面追踪被接受的快速搜索候选（圆点）、由第一个收敛大气给出的可选修正（菱形）以及收敛大气加合成的最终解（方块）；颜色为累计模型时间，空心星为真值。下图在代表性波长区间比较模拟观测与最终物理模型，阴影为 $1\sigma$ 噪声带。恢复的光谱与标签使用直接合成，无任何标签到流量模拟器。（Controlled five-label fit to a normalized mock spectrum.）}
\label{fig:fitter_convergence}
\end{figure}

\begin{notebox}{算法 1 说明（内容保留英文伪代码）}
受控直接光谱拟合流程：定义每有效像素平均卡方 $Q$ → 从初始标签反复做“直接合成 + 有界最小二乘步 + 线搜索”直到无法下降 → 在快速极小点收敛物理大气 $\mathbf{x}_{\rm phys}$ 并合成 $f_{\rm phys}$，暂定为解 → 固定 $\Delta f=f_{\rm phys}-f_{\rm fast}$ 重建响应并提出有界修正 $\boldsymbol{\theta}_{\rm p}$ → 仅当预测 $Q$ 下降超过 $\delta Q_{\min}$ 才在 $\boldsymbol{\theta}_{\rm p}$ 收敛第二个大气并合成 $f_{\rm prop}$，若 $Q(f_{\rm prop})<Q(f_{\rm phys})$ 则接受 → 返回 $(\boldsymbol{\theta}_\star,\mathbf{x}_\star,f_\star)$。
\end{notebox}

\begin{table}[t]
\centering
\begin{minipage}{\columnwidth}
\centering
\textbf{Algorithm 1. Controlled direct spectrum fit}\par
\vspace{2pt}
\noindent\rule{\columnwidth}{0.5pt}
\vspace{2pt}
\begin{algorithmic}[1]
\Require flux $y_p$, weights $w_p$, initial labels $\boldsymbol{\theta}_0$,
threshold $\delta Q_{\min}$
\Ensure labels $\boldsymbol{\theta}_\star$, atmosphere state $\mathbf{x}_\star$,
spectrum $f_\star$
\State Define $Q(f)\gets N_{\rm good}^{-1}\sum_{p\in{\rm good}}w_p(y_p-f_p)^2$
\State $\boldsymbol{\theta}\gets\boldsymbol{\theta}_0$
\Repeat
  \State Evaluate direct $f_{\rm fast}$ and its local response
  \State Propose a bounded least-squares step
  \State Retain it only if direct synthesis lowers $Q$
\Until{no bounded trial lowers $Q$}
\State $\boldsymbol{\theta}_{\rm f}\gets\boldsymbol{\theta}$
\State Solve $\mathbf{x}_{\rm phys}$ at $\boldsymbol{\theta}_{\rm f}$ and synthesize $f_{\rm phys}$
\State Set $(\boldsymbol{\theta}_\star,\mathbf{x}_\star,f_\star)
  \gets(\boldsymbol{\theta}_{\rm f},\mathbf{x}_{\rm phys},f_{\rm phys})$
\State Fix $\Delta f\gets f_{\rm phys}-f_{\rm fast}$
\State Rebuild the response
\State Propose bounded correction $\boldsymbol{\theta}_{\rm p}$
\If{its predicted decrease exceeds $\delta Q_{\min}$}
  \State Solve $\mathbf{x}_{\rm prop}$ at $\boldsymbol{\theta}_{\rm p}$ and synthesize $f_{\rm prop}$
  \If{$Q(f_{\rm prop})<Q(f_{\rm phys})$}
    \State Set $(\boldsymbol{\theta}_\star,\mathbf{x}_\star,f_\star)
      \gets(\boldsymbol{\theta}_{\rm p},\mathbf{x}_{\rm prop},f_{\rm prop})$
  \EndIf
\EndIf
\State \Return $(\boldsymbol{\theta}_\star,\mathbf{x}_\star,f_\star)$
\end{algorithmic}
\vspace{2pt}
\noindent\rule{\columnwidth}{0.5pt}
\end{minipage}
\end{table}


\section{谱线表定标 Line List Calibration}\label{sec:linecal}

受控模拟确立了前向模型内部自洽时的恢复能力。观测光谱还会进一步检验模型的物理输入。沿用成熟的近红外天体物理谱线定标传统，我们考察可微分计算能否联合更新振子强度与三个阻尼族 \citep{MelendezBarbuy1999,Shetrone2015,Smith2021}。每项提议的修正都被约束为只能通过谱线不透明度与辐射转移来改变光谱。定标的评估对象是太阳与大角星的高分辨率傅里叶变换光谱仪（FTS）图集。

\begin{EN}
The controlled mock establishes recovery when the forward model is internally
consistent. Observed spectra also test the physical inputs to that model.
Following established near-infrared astrophysical line calibration, we ask
whether the differentiable calculation can update oscillator strength and the
three damping families jointly
\citep{MelendezBarbuy1999,Shetrone2015,Smith2021}. Each proposed
correction is constrained to change the spectrum through line opacity and
radiative transfer. We evaluate the corrections against high-resolution
Fourier transform spectrometer (FTS) atlases of the Sun and Arcturus.
\end{EN}

\begin{physbox}{为什么用太阳和大角星定标原子数据}
振子强度 $gf$ 与阻尼常数原则上可由实验室或原子结构理论给出，但许多谱线（尤其近红外、过渡族元素）的数值误差很大甚至缺失。补救办法是\textbf{天体物理定标}：用光谱质量最高的恒星反推这些参数。太阳 FTS 图集来自傅里叶变换光谱仪，分辨本领可达 $10^6$、信噪比极高，且太阳的物理参数（$T_{\rm eff}=5777$ K、$\log g=4.44$、丰度）是所有恒星中测得最准的。大角星（Arcturus）是 K 型红巨星的标准星，气压低、分子带丰富，与太阳处于互补的物理区间。两者联合定标能约束不同压强与分子环境下的谱线行为。本文的新意在于：由于整个合成在 GPU 上可微分，超过 $10^5$ 个谱线参数可以\textbf{一次性联合优化}，而传统做法是逐条谱线手工调整。
\end{physbox}

\begin{figure}[t]
\centering
\includegraphics[width=0.98\textwidth]{figure_08.pdf}
\caption{对太阳与大角星 FTS 图集的全波段定标，分六个连续波长条带。每条带上图为归一化流量、下图为图集减模型。黑色为图集，灰虚线为原始谱线数据，橙虚线为对该标准星的独立四参数谱线拟合。残差面板在同一颗星内共用固定范围。大角星列显示冬季图集，其拟合对夏、冬两历元等权。定标使整段区间的“残差森林”在两颗星上都明显收缩。（Full-band calibration of the solar and Arcturus FTS atlases.）}
\label{fig:linecal_overview}
\end{figure}

对分配给修正组 $q(l)$ 的谱线分量 $l$，定标量写为
\begin{equation}
\begin{aligned}
(gf)'_l &= (gf)_l 10^{\delta_{q(l),gf}},\\
\gamma'_{l,\beta} &= \gamma_{l,\beta}10^{\delta_{q(l),\beta}},
\qquad \beta\in\{\mathrm{vdW},\mathrm{rad},\mathrm{S}\}.
\end{aligned}
\end{equation}
\noindent 其中 $\gamma_{l,\beta}$ 为谱线分量 $l$、阻尼族 $\beta$ 的原始阻尼系数；三个族为范德瓦尔斯、辐射与斯塔克阻尼；撇号表示定标后的值。第四类拟合参数是 $\log(gf)$。在布居固定时，振子强度缩放频率积分的不透明度，阻尼把不透明度从线心重新分配到翼部 \citep{Barklem2000,Gray2005}；二者在饱和与混合特征中的响应有重叠。

\begin{EN}
For line component $l$ assigned to correction group $q(l)$, we write the
calibrated quantities as
\begin{equation}
\begin{aligned}
(gf)'_l &= (gf)_l 10^{\delta_{q(l),gf}},\\
\gamma'_{l,\beta} &= \gamma_{l,\beta}10^{\delta_{q(l),\beta}},
\qquad \beta\in\{\mathrm{vdW},\mathrm{rad},\mathrm{S}\}.
\end{aligned}
\label{eq:atomic_calibration_parameters}
\end{equation}
\noindent Here $\gamma_{l,\beta}$ is the original damping coefficient for line
component $l$ and family $\beta$. The three families are van der Waals,
radiative, and Stark damping, while the prime denotes the calibrated value.
The fourth fitted atomic-parameter type is $\log(gf)$. At fixed population, oscillator strength
scales the frequency-integrated opacity, while damping redistributes opacity
from the core into the wings \citep{Barklem2000,Gray2005}. Their responses
overlap in saturated and blended features.
\end{EN}

PyTorch 记录从每项原子修正经不透明度、辐射转移、致宽、采样、连续谱轮廓化直到最终目标的完整路径。一次反向传播即可同时给出所有活跃修正的导数，无需有限差分 \citep{Baydin2015}。逐条谱线的独立辐射转移工作在光谱推断中演示过同样的计算原理 \citep{Kawahara2022}。固定的大气、连续与分子不透明度以及目录几何全程常驻 GPU。

\begin{EN}
PyTorch records the path from each atomic correction through opacity,
radiative transfer, broadening, sampling, continuum profiling, and the final
objective. A reverse pass gives derivatives for every active correction
together, without finite differences \citep{Baydin2015}. Independent
line-by-line radiative-transfer work demonstrates the same computational
principle in spectral inference \citep{Kawahara2022}. The fixed atmosphere, continuous and
molecular opacity, and catalog geometry remain resident on the GPU.
\end{EN}

\begin{physbox}{自动微分为什么关键}
要优化 101,124 个参数，必须知道目标函数对每个参数的导数。用有限差分需要每个参数各算一条扰动光谱——十万余次合成，不可行。\textbf{自动微分}（反向传播）利用链式法则，沿计算图反向走一遍就能同时得到所有参数的精确导数，成本仅为一次正向合成的数倍。这正是深度学习框架（PyTorch）被引入光谱计算的红利：只要把所有运算写成框架内的张量操作，梯度就“免费”获得。由此，逐条谱线人工调参的传统定标方式被一个联合辐射转移优化问题取代。
\end{physbox}

由于谱线中心通过整数网格放置进入不透明度叠加，计算图尚不能提供有效的波长导数，因此波长保持固定。低激发能也在本四参数概念验证之外——本次只定标公式~\ref{eq:atomic_calibration_parameters} 中的量。

\begin{EN}
Because line centers enter opacity deposition through integer grid placement,
the graph does not yet provide a valid wavelength derivative. Wavelengths are
therefore fixed. Lower excitation energies also remain outside this
four-parameter proof of concept, which calibrates the quantities in
Equation~\ref{eq:atomic_calibration_parameters}.
\end{EN}

我们使用 1500--1700 nm 区间以直接衔接 APOGEE 波长范围，并针对 FTS 图集而非 APOGEE 像素做定标。合成采用 $R_{\rm grid}=300{,}000$，然后采样到各图集网格。太阳用 Livingston 与 Wallace 的日面中心图集，大角星用 Hinkle 等人的两个历元 \citep{LivingstonWallace1991,Hinkle1995,Smith2021}。太阳取 $T_{\rm eff}=5777$ K、$\log g=4.44$、$\xi=0.7$ km s$^{-1}$；大角星取 $T_{\rm eff}=4286$ K、$\log g=1.66$、$[\mathrm{M}/\mathrm{H}]=-0.52$、$[\alpha/\mathrm{M}]=0.30$、$\xi=1.7$ km s$^{-1}$ \citep{Ramirez2011,Smith2021}，逐元素混合物采用 \citet{Smith2021} 的汇总表（多数值取自 \citet{Ramirez2011} 并辅以其他专门研究）。

\begin{EN}
We use the 1500--1700 nm interval to connect directly to the APOGEE wavelength
range. We perform the calibration against these FTS atlases rather than APOGEE
pixels. The synthesis calculation uses
$R_{\rm grid}=300{,}000$ before sampling onto each atlas grid. We use the
Livingston and Wallace solar center-of-disk atlas and two Hinkle et al.
Arcturus epochs \citep{LivingstonWallace1991,Hinkle1995,Smith2021}. For the
solar calculation we adopt $T_{\rm eff}=5777$ K,
$\log g=4.44$, and $\xi=0.7$ km s$^{-1}$. For Arcturus we adopt
$T_{\rm eff}=4286$ K, $\log g=1.66$, $[\mathrm{M}/\mathrm{H}]=-0.52$,
$[\alpha/\mathrm{M}]=0.30$, and $\xi=1.7$ km s$^{-1}$
\citep{Ramirez2011,Smith2021}.

We adopt the
element-by-element Arcturus mixture tabulated by
\citet{Smith2021}. That compilation draws most values from
\citet{Ramirez2011} and supplements them from other element-specific
studies.
\end{EN}

每个标准星的收敛物理大气在内层定标中保持固定。太阳取 1.55 km s$^{-1}$ 的高斯致宽；大角星取 2.5 km s$^{-1}$ 高斯致宽加 $v\sin i=2.0$ km s$^{-1}$ 自转。给定的多普勒因子与残余配准修正把图集放入模型静止系。每次求值都用图集与基线模型中靠近连续谱的像素轮廓化一条平滑的乘性连续谱，紧凑的系数边界防止该调整吸收谱线轮廓失配。配套的地球大气透过率数组定义像素相对权重并掩去不可靠波长。

\begin{EN}
Each standard's converged physical atmosphere remains fixed during the inner calibration. We adopt a
Gaussian dispersion of 1.55~km~s$^{-1}$ for the Sun. For Arcturus we adopt a Gaussian
dispersion of 2.5~km~s$^{-1}$ and $v\sin i=2.0$~km~s$^{-1}$. The supplied Doppler
factors and a residual registration correction place the atlases in the model
rest frame.

At every evaluation we profile a smooth multiplicative continuum using pixels
near the continuum in both the atlas and the baseline model. Tight coefficient
bounds prevent this adjustment from absorbing line-profile mismatch.
Companion telluric-transmission arrays define relative pixel weights and mask
unreliable wavelengths.
\end{EN}

对每个图集，我们在地球大气降权与连续谱轮廓化之后最小化均方流量差。独立的大角星解对夏、冬历元等权，独立的太阳解只用太阳图集——这些拟合衡量各标准星的定标容量。随后拟合一个共享的太阳—大角星解，对两颗星等总权、对大角星两历元等权；APOGEE 应用即使用这一公共定标。

\begin{EN}
For each atlas, we minimize the mean squared flux difference after telluric
downweighting and continuum profiling. The independent Arcturus solution fits
the summer and winter epochs with equal weight, while the independent solar
solution uses only the solar atlas. These fits measure the calibration capacity
of each standard. We then fit one shared Sun--Arcturus solution, giving equal
total weight to each star and equal weight to the two Arcturus epochs. We use
this common calibration for APOGEE.
\end{EN}

我们纳入该区间内所有在丰度向量中出现的物种的中性与一次电离原子记录。先保留一小组精选的锚定谱线群，再按强度排序的贪心遍历把每个剩余分量与同物种、同电离级、下激发能相差 2.0 cm$^{-1}$、波长相差 0.035 nm 以内的未使用谱线分组。42,886 条目录记录形成 25,281 个组，共 101,124 个拟合标量修正；同组所有分量共享其四个修正。

\begin{EN}
We include neutral and singly ionized atomic records in the interval for every
species represented in the adopted abundance vector. We first preserve a small
set of curated anchor complexes. A strength-ordered greedy pass then groups
each remaining component with unused lines of the same species and ion stage
within 2.0~cm$^{-1}$ in lower excitation and 0.035~nm in wavelength. The 42,886
catalog records form 25,281 groups and give 101,124 fitted scalar corrections.
All components in one group share its four corrections.
\end{EN}

我们用 PyTorch L-BFGS 联合优化完整 200 nm 区间 \citep{NocedalWright2006}。这种基于梯度的方法用近期梯度估计曲率，强 Wolfe 线搜索只接受以降低目标且局部斜率一致的步长。平滑参数变换保证物理边界。联合优化让混合像素同时约束多条跃迁，避免把一条谱线轮廓拆给各自独立拟合的波长窗口。

\begin{EN}
We optimize the complete 200 nm interval jointly with PyTorch L-BFGS
\citep{NocedalWright2006}. This gradient-based method uses recent gradients to estimate
curvature, while a strong-Wolfe line search accepts steps that lower the
objective with a consistent local slope. Smooth parameter transforms enforce
the physical bounds. Joint optimization lets blended pixels constrain several
transitions and avoids splitting a line profile among independently fitted
wavelength windows.
\end{EN}

\begin{notebox}{表~\ref{tab:linecal_ablation} 与表~\ref{tab:linecal_overlay} 说明（原文 Table 5、6，内容保留英文）}
表 5：相对原始谱线数据，地球大气降权后“图集−模型”均方流量差的下降百分比。只调 $\log(gf)$ 即可下降 78\%（太阳）/57\%（大角星）；四类参数全调达 90\%/65\%；共享太阳—大角星解为 84\%/56\%。各参数类型的贡献不可相加，因为它们能解释部分相同的残差。表 6：发布的公共定标文件的内容与用法——源目录摘要与分量映射锁定谱线身份；每个链接组存四个对数修正：$\log(gf)$ 直接相加，三类阻尼乘以 $10^{\delta}$。
\end{notebox}

\begin{table}[t]
\centering
\captionof{table}{Reduction relative to the original line data in the
telluric-downweighted mean squared atlas--model flux difference. The first five
rows fit each standard independently. The final row fits the Sun and Arcturus
with one common set of line-parameter corrections.}
\label{tab:linecal_ablation}
\begingroup
\setlength{\tabcolsep}{6pt}
\renewcommand{\arraystretch}{1.12}
\footnotesize
\begin{tabular}{@{}lrr@{}}
\toprule
\textbf{Calibration case} & \textbf{Sun} & \textbf{Arcturus} \\
\midrule
$\log(gf)$ & 78.2\% & 56.8\% \\
van der Waals damping & 36.9\% & 7.7\% \\
radiative damping & 19.0\% & 5.4\% \\
Stark damping & 26.6\% & 2.7\% \\
all four types & 90.1\% & 64.9\% \\
\midrule
shared Sun--Arcturus fit & 83.5\% & 56.0\% \\
\addlinespace[2pt]
\bottomrule
\end{tabular}
\endgroup
\end{table}

\begin{table}[t]
\centering
\captionof{table}{Contents and application of the released common
Sun--Arcturus line-list calibration. Source identities bind the corrections to
the intended catalog. Each linked transition group stores four logarithmic
corrections, from which the calibrated component values are obtained.}
\label{tab:linecal_overlay}
\begingroup
\setlength{\tabcolsep}{4pt}
\renewcommand{\arraystretch}{1.12}
\footnotesize
\begin{tabular}{@{}lll@{}}
\toprule
\tcell{0.24\linewidth}{\textbf{Quantity}} &
\tcell{0.38\linewidth}{\textbf{Application}} &
\tcell{0.27\linewidth}{\textbf{Physical role}} \\
\midrule
\tcell{0.24\linewidth}{source identity} & \tcell{0.38\linewidth}{catalog digest and\\component map} & \tcell{0.27\linewidth}{selects the exact\\source transitions} \\
\tcell{0.24\linewidth}{linked group} & \tcell{0.38\linewidth}{component $\rightarrow$\\fitted group} & \tcell{0.27\linewidth}{shares one correction\\among linked lines} \\
\midrule
\tcell{0.24\linewidth}{$\log(gf)$} & \tcell{0.38\linewidth}{$\log(gf)_{\rm cal}=$\\$\log(gf)_{\rm base}+\delta_{q,gf}$} & \tcell{0.27\linewidth}{integrated line\\opacity} \\
\tcell{0.24\linewidth}{van der Waals\\damping} & \tcell{0.38\linewidth}{$\gamma_{\rm cal}=$\\$\gamma_{\rm base}10^{\delta_{q,\rm vdW}}$} & \tcell{0.27\linewidth}{neutral-collision\\wings} \\
\tcell{0.24\linewidth}{radiative damping} & \tcell{0.38\linewidth}{$\gamma_{\rm cal}=$\\$\gamma_{\rm base}10^{\delta_{q,\rm rad}}$} & \tcell{0.27\linewidth}{natural line\\width} \\
\tcell{0.24\linewidth}{Stark damping} & \tcell{0.38\linewidth}{$\gamma_{\rm cal}=$\\$\gamma_{\rm base}10^{\delta_{q,\rm S}}$} & \tcell{0.27\linewidth}{electron-pressure\\wings} \\
\bottomrule
\end{tabular}
\endgroup
\end{table}

表~\ref{tab:linecal_ablation} 把单参数对照与每个标准星的独立四参数拟合相比较。下降量不可相加，因为各参数类型可解释部分相同的残差。振子强度在两颗星上都贡献了主要增益，四参数拟合则进一步改进。阻尼在太阳光谱中杠杆更大——更高的光球压强产生更强的压致翼 \citep{Barklem2000,Gray2005}；在大角星中，拟合经验上把大部分增益归于线强，其拥挤的分子光谱提供的孤立翼部约束更少。

\begin{EN}
Table~\ref{tab:linecal_ablation} compares the single-parameter controls with
each standard's independent four-parameter fit. The reductions are not additive because
the parameter types can explain part of the same residual. Oscillator strength
carries most of the gain in both stars, while the four-parameter fit improves
further in each case. Damping adds more leverage in the solar
spectrum, where the higher photospheric pressure produces stronger
pressure-broadened wings \citep{Barklem2000,Gray2005}. In Arcturus, the fit
empirically assigns most of the gain to line strength. Its crowded molecular
spectrum provides fewer isolated wing constraints.
\end{EN}

\begin{physbox}{为什么阻尼在太阳上更重要}
斯塔克与范德瓦尔斯致宽都来自碰撞：分别由带电粒子（电子、离子）和中性原子的扰动造成。碰撞频率正比于粒子密度，也就是气压。太阳（矮星，$\log g=4.44$）光球密度比大角星（巨星，$\log g=1.66$）高两个数量级以上，因此强线的阻尼翼在太阳光谱中宽阔显著，约束阻尼参数的信息多；在巨星中翼部窄，拟合自然把改进主要分配给线心强度（$gf$）。这也是为什么好的原子数据定标需要\textbf{多颗处于不同物理区间的标准星}——单颗星无法把线强与阻尼的退化解开。
\end{physbox}

有 101,124 个活跃修正时，反向传播避免了“每参数一条扰动光谱”。从独立解的键对齐均值出发，40 步共享太阳—大角星更新在单张 H100 上用 64 秒。

\begin{EN}
With 101,124 active corrections, the reverse pass avoids synthesizing one
perturbed spectrum per parameter. Starting from the key-aligned mean of the
independent solutions, 40 shared Sun--Arcturus updates take 64~s on one H100.
\end{EN}

图~\ref{fig:linecal_overview} 显示目标下降并不局限于少数孤立特征：残差森林在两个图集的整个展示区间内收缩。太阳的改进更强，大角星在其更拥挤的光谱中保留更大的结构化残差。图~\ref{fig:linecal_residual_distribution} 给出太阳图集与大角星冬季历元的未加权残差分布；表~\ref{tab:linecal_ablation} 报告的则是包含大角星两历元、经地球大气加权的均方下降。独立与共享拟合都收缩了长残差尾。

\begin{EN}
Figure~\ref{fig:linecal_overview} shows that the objective reduction is not
confined to a few isolated features. The residual forest contracts across the
displayed interval for both atlases. The change is stronger for the Sun, while
Arcturus retains larger structured residuals in its more crowded spectrum.

Figure~\ref{fig:linecal_residual_distribution} shows unweighted residual
distributions for the solar atlas and the displayed Arcturus winter epoch.
Table~\ref{tab:linecal_ablation} instead reports telluric-weighted mean-square
reductions over the calibration data, including both Arcturus epochs. Both the
independent and shared fits contract the long residual tail.
\end{EN}

图~\ref{fig:linecal_profiles} 放大六条单独的谱线轮廓。这些太阳示例由算法从共享太阳—大角星拟合中选取，未做局部重拟合；它们是可归属、波长上分开的不同物种跃迁。共享定标在这些示例中同时改进了线深与轮廓形状。

\begin{EN}
Figure~\ref{fig:linecal_profiles} zooms in on six individual line profiles. The
solar examples are selected algorithmically from the
shared Sun--Arcturus fit and are not refitted locally. They are attributable,
wavelength-separated transitions from different species. The shared
calibration improves both line depths and profile shapes across these examples.
\end{EN}

我们把共享太阳—大角星定标作为紧凑的修正文件发布。把它施加到匹配的基底目录即得到表~\ref{tab:linecal_overlay} 总结的定标谱线参数；\texttt{write\_substituted\_catalog} 辅助程序可把这些值写成完整的替换目录，默认目录保持不变。独立的太阳与大角星拟合作为诊断产物保留。

\begin{EN}
We distribute the shared Sun--Arcturus calibration as a compact correction
file. Applying it to the matching base catalog produces the calibrated line
parameters summarized in Table~\ref{tab:linecal_overlay}. The
\texttt{write\_substituted\_catalog} helper can write those values as a complete
replacement catalog, while the default catalog remains unchanged. Independent
solar and Arcturus fits remain available as diagnostic products.
\end{EN}

这些是\textbf{模型依赖的标准星定标，而非实验室测量} \citep{Heiter2015,Heiter2021}，且内层拟合中每个标准星大气保持固定。生产级定标应在多个标准星之间共享信息 \citep{Heiter2015,Jofre2015,Smith2021}，并在大幅度不透明度更新之后重复大气与谱线表求解。

\begin{EN}
These are model-dependent standard-star calibrations, not laboratory
measurements \citep{Heiter2021}, and each standard-star atmosphere remains fixed
during the inner fit. A production calibration should share information across
several standards \citep{Heiter2015,Jofre2015,Smith2021}. After a substantial
opacity update, it should also repeat the atmosphere and line-list solutions.
\end{EN}

部分大残差仍然存在，因为当前参数化不能移动谱线中心，也不能创造缺失的跃迁。我们探索过一个由残差引导的预言原子跃迁子集，但留出测试不支持其转正。因此预言跃迁在进入通用谱线表之前，需要在多颗独立恒星上验证。

\begin{EN}
Some large residuals remain because the current parameterization cannot move
line centers or create missing transitions. We explored a residual-guided
subset of predicted atomic transitions, but held-out tests did not justify
promotion. Predicted transitions therefore require validation across
independent stars before inclusion in a general line list.
\end{EN}

\begin{figure}[t]
\centering
\includegraphics[width=0.9\textwidth]{figure_09.pdf}
\captionof{figure}{太阳与大角星的归一化流量绝对残差分布：横轴为残差阈值，纵轴为绝对残差超过该值的保留像素占比。灰虚线为原始谱线数据，蓝点划线只调振子强度，橙实线为各标准星独立四参数拟合，品红点线为共享太阳—大角星拟合。（Absolute normalized-flux residual distributions.）}
\label{fig:linecal_residual_distribution}
\vspace{8pt}
\includegraphics[width=0.9\textwidth]{figure_10.pdf}
\captionof{figure}{共享太阳—大角星定标下六条太阳日心 FTS 谱线轮廓。波长以所标跃迁的目录实验室中心为原点。黑色为图集，灰虚线为原始谱线数据，品红点线为振子强度加三项阻尼的共享拟合。六例跨六种原子物种，线深与轮廓形状均有改进。（Six solar center-of-disk FTS line profiles.）}
\label{fig:linecal_profiles}
\end{figure}

\section{在 APOGEE 上的应用 Application to APOGEE}\label{sec:apogee}

我们现在问：同一套计算能否直接拟合约减后的巡天光谱。我们把物理拟合与保留的标准星谱线定标结合，应用到 APOGEE 第 14 次数据释放（DR14）。这台光纤馈源近红外光谱仪及其交付分辨率见 \citet{Wilson2019}。选用该释放是因为它提供了受控比较所需的合并光谱、仪器定标与巨星分析产物 \citep{Abolfathi2018,Holtzman2018}。APOGEE 恒星参数与化学丰度流水线（ASPCAP）的标签提供外部对照 \citep{GarciaPerez2016,Jonsson2018}。

\begin{EN}
We now ask whether the same calculation can fit reduced survey spectra
directly. We combine physical fitting with the retained standard-star line
calibration and apply it to APOGEE Data Release 14 (DR14).
The fiber-fed near-infrared spectrograph and its delivered resolution are
described by \citet{Wilson2019}. We use this release because it provides the combined spectra, instrument
calibrations, and giant-star analysis products needed for the controlled
comparison \citep{Abolfathi2018,Holtzman2018}. Labels from the APOGEE Stellar
Parameters and Chemical Abundances Pipeline (ASPCAP) provide the
external comparison
\citep{GarciaPerez2016,Jonsson2018}.
\end{EN}

\begin{physbox}{APOGEE 与 ASPCAP 背景}
APOGEE 是架设在阿帕奇点天文台的高分辨率（$R\sim22{,}500$）近红外（1.5--1.7 $\mu$m，H 波段）巡天，观测了数十万颗银河系恒星，以红巨星为主。近红外的优势是尘埃消光小，能看穿银河盘。ASPCAP 是其官方分析流水线：把观测光谱与预计算的合成光谱网格（由 FERRE 驱动）做多维 $\chi^2$ 匹配，输出恒星参数与二十余种元素丰度，最后再施加经验零点与温度定标。本文的做法与 ASPCAP 的根本区别：不预计算网格、不插值，每颗星都从物理模型实时合成；也不施加 ASPCAP 的经验零点——这正是后文 Si 丰度出现刚性偏移的原因。
\end{physbox}

在这一巡天演示中，DR14 把其原始光谱型 $T_{\rm eff}$ 定标到低消光光度测量标度，把巨星 $\log g$ 定标到星震学重力，并把微观湍流作为单独的巨星坐标拟合 \citep{Holtzman2018}。我们不推导新的光谱—光度温度关系；本概念验证保留已发表的 $T_{\rm eff}$、$\log g$、$\xi$，以单独考察多元素与仪器问题。生产级分析可以把光谱与宽波段能谱分布、视差联合拟合 \citep{Cargile2020}。附录~\ref{app:stellar_label_sensitivity} 重复了放开 $T_{\rm eff}$、$\log g$、$\xi$ 的合成搜索，检验巨星支、红团簇与丰度形态对该选择的敏感性。

\begin{EN}
For the survey demonstration, DR14 calibrates its raw spectroscopic $T_{\rm eff}$ to a
low-reddening photometric scale and giant $\log g$ to asteroseismic gravities
and fits microturbulence as a separate giant-star coordinate
\citep{Holtzman2018}. We do not derive a new relation between spectroscopic and
photometric temperature. For this proof of concept, we retain the published
$T_{\rm eff}$, $\log g$, and $\xi$ to isolate the many-element and instrument
problem. A production analysis could instead fit the spectrum jointly with a
broadband spectral energy distribution and parallax \citep{Cargile2020}.
Appendix~\ref{app:stellar_label_sensitivity} repeats the synthesis
search with $T_{\rm eff}$, $\log g$, and $\xi$ free to test the sensitivity of
the giant-branch, red-clump, and abundance morphology to this choice.
\end{EN}

我们拟合 $[\mathrm{Fe}/\mathrm{H}]$，C、N、O、Mg、Al、Si、S、Ca、Ti、Mn、Ni 的 $[\mathrm{X}/\mathrm{Fe}]$，残余速度 $v_r$ 与有效高斯致宽 $v_b$。每个保留元素满足 $[\mathrm{X}/\mathrm{H}]=[\mathrm{Fe}/\mathrm{H}]+[\mathrm{X}/\mathrm{Fe}]$，未选金属跟随 $[\mathrm{Fe}/\mathrm{H}]$。$v_b$ 合并自转、宏观致宽以及平均仪器模型未表示的残余轮廓宽度。

\begin{EN}
We fit $[\mathrm{Fe}/\mathrm{H}]$, $[\mathrm{X}/\mathrm{Fe}]$ for C, N, O,
Mg, Al, Si, S, Ca, Ti, Mn, and Ni, residual velocity $v_r$, and effective
Gaussian broadening $v_b$. For every retained element,
$[\mathrm{X}/\mathrm{H}]=[\mathrm{Fe}/\mathrm{H}]+[\mathrm{X}/\mathrm{Fe}]$,
while unselected metals follow $[\mathrm{Fe}/\mathrm{H}]$. The fitted width
$v_b$ combines rotation, macroscopic broadening, and residual profile width not
represented by the mean instrument model.
\end{EN}

合成返回原生网格上的总量与连续谱流量。归一化之前，先把两者映射到观测像素：
\begin{equation}
f_p(\boldsymbol{\psi})=
\frac{\left\{\mathcal{K}\mathcal{S}(v_r)\mathcal{B}(v_b)
F_\lambda(\boldsymbol{\theta})\right\}_p}
{\left\{\mathcal{K}\mathcal{S}(v_r)\mathcal{B}(v_b)
F_{\lambda,{\rm c}}(\boldsymbol{\theta})\right\}_p},
\qquad
\boldsymbol{\psi}=(\boldsymbol{\theta},v_r,v_b).
\end{equation}
\noindent 其中 $\mathcal{B}$ 作用于对数合成网格，$\mathcal{S}$ 施加多普勒平移，$\mathcal{K}$ 施加随波长变化的线扩展函数与 APOGEE 原生采样。先对两者卷积再相除，保留了非线性的归一化。

\begin{EN}
Synthesis returns native-grid total and continuum fluxes. We map both to the
observed pixels before normalization,
\begin{equation}
f_p(\boldsymbol{\psi})=
\frac{\left\{\mathcal{K}\mathcal{S}(v_r)\mathcal{B}(v_b)
F_\lambda(\boldsymbol{\theta})\right\}_p}
{\left\{\mathcal{K}\mathcal{S}(v_r)\mathcal{B}(v_b)
F_{\lambda,{\rm c}}(\boldsymbol{\theta})\right\}_p},
\qquad
\boldsymbol{\psi}=(\boldsymbol{\theta},v_r,v_b).
\label{eq:apogee_forward}
\end{equation}
\noindent where $\mathcal{B}$ acts on the logarithmic synthesis grid,
$\mathcal{S}$ applies the Doppler shift, and $\mathcal{K}$ applies the
wavelength-dependent line-spread function and native APOGEE sampling.
Convolving both before division preserves the nonlinear normalization.
\end{EN}

本实验使用 DR14 定标导出的 LSF \citep{Nidever2015,Holtzman2018}，对六根代表性光纤取平均。一次性构建把从 $R_{\rm grid}=300{,}000$ 合成网格的插值折叠成紧凑的波长依赖权重并常驻 GPU；该算子直接映射到 APOGEE 像素，无需构建稠密卷积矩阵。换用光纤、曝光或探测器专属的 LSF 只需更换一份预置 GPU 资产。致宽、平移、LSF 卷积与采样在 H100 上对每对“总量—连续”流量约 $240~\mu\mathrm{s}$，而原生合成为 $1.4$ 秒——仪器模型的额外成本可忽略。公开拟合器还提供 GPU 常驻自转 kernel，本演示保留单一高斯宽度。

\begin{EN}
For this experiment, we use a DR14 calibration-derived LSF
\citep{Nidever2015,Holtzman2018}, averaged over six representative fibers. A
one-time build folds interpolation from the $R_{\rm grid}=300{,}000$ synthesis
grid into compact wavelength-dependent weights that remain on the GPU. The
operator maps directly onto the APOGEE pixels without constructing a dense
convolution matrix. A fiber-, exposure-, or detector-specific LSF requires
only a different prepared GPU asset. Broadening, shifting, LSF convolution,
and sampling take about $240~\mu\mathrm{s}$ per total--continuum flux pair on
an H100, compared with $1.4$~s for native synthesis. The instrument model
therefore adds negligible cost. The public fitter also provides a GPU-resident
rotational kernel, although this demonstration retains one Gaussian width.
\end{EN}

伪连续谱归一化的 DR14 光谱中仍残留平滑的连续谱差异 \citep{Holtzman2018}。每次非线性试探时，加权线性最小二乘在每个探测器上独立轮廓化一组四阶勒让德基 \citep{GolubPereyra2003}。似然使用已发表的不确定度与掩模 \citep{Holtzman2018}，外加 0.5\% 的模型下限。由于合并光谱组装在公共静止波长网格上 \citep{Nidever2015}，拟合从 $v_r=0$ 出发，只解局部残余平移。

\begin{EN}
Smooth continuum differences remain in the pseudo-continuum-normalized DR14
spectra \citep{Holtzman2018}. At every nonlinear trial, weighted linear least
squares profiles a fourth-order Legendre basis independently on each detector
\citep{GolubPereyra2003}. The likelihood uses the published uncertainties and
masks \citep{Holtzman2018} with a 0.5 percent model floor.
Because the combined spectra are assembled on a common rest-wavelength grid
\citep{Nidever2015}, the fit begins at $v_r=0$ and solves only a local
residual shift.
\end{EN}

观测光谱搜索由初始化大气与直接合成驱动。每增加一个丰度只需一条响应光谱，成本随标签数近似线性增长而非组合爆炸。随后物理求解器在所选混合物处收敛大气，一条新的 GPU 光谱提供任何最终的局部修正。被接受的模型永远使用物理收敛大气加直接合成。这一 12 丰度拟合在 H100 上搜索约 40 秒/星，大气收敛在多核 CPU 上再加约 30 秒。

\begin{EN}
Initialized atmospheres and direct synthesis drive the observed-spectrum
search.
Each additional abundance requires one response spectrum, so this cost grows
approximately linearly rather than combinatorially with the label count. The
physical solver then converges an atmosphere at the selected mixture, and a
fresh GPU spectrum supplies any final local correction. The accepted model
always uses a physically converged atmosphere and direct synthesis. For this
12-abundance fit, the H100 search takes about 40~s per star, and atmosphere
convergence adds about 30~s on multicore CPUs.
\end{EN}

作为概念验证，我们选取 1,600 颗巨星，在高、低 $\alpha$ APOGEE 序列与金属丰度上均衡采样 \citep{Nidever2014,Hayden2015}，并对族群轮廓重加权以代表母样本。图~\ref{fig:apogee_spectrum} 首先在像素空间检验计算：跨全部三个探测器，前向模型跟随谱线森林，轮廓化连续谱只吸收平滑的残余结构。这确立了仪器算子与多元素光谱可以联合拟合。

\begin{EN}
As a proof of concept, we select 1,600 giants, balanced across the high- and
low-$\alpha$ APOGEE sequences and metallicity
\citep{Nidever2014,Hayden2015}, and reweight the population contours to
represent the parent sample. Figure~\ref{fig:apogee_spectrum} first tests the
calculation in pixel space.
Across all three detectors, the forward model follows the line forest while the
profiled continuum absorbs only smooth residual structure. This establishes
that the instrument operators and many-element spectrum can be fitted together.
\end{EN}

随后检验独立拟合的丰度是否保持族群结构。图~\ref{fig:apogee_structure} 在被采样的金属丰度范围内保留了两条主要的 $[\mathrm{Mg}/\mathrm{Fe}]$ 序列，而不是塌缩成单一整体金属丰度趋势。图~\ref{fig:apogee_elements} 把比较推广到逐个元素：O、Si、S、Ca、Ti 在两个输入族群之间保持了化学分离，但绝对零点不同。

\begin{EN}
We then test whether the independently fitted abundances preserve population
structure. Figure~\ref{fig:apogee_structure} retains the two principal
$[\mathrm{Mg}/\mathrm{Fe}]$ sequences over the sampled metallicity range rather
than collapsing them into one bulk-metallicity trend. Figure~\ref{fig:apogee_elements}
extends the comparison element by element. O, Si, S, Ca, and Ti retain chemical
separation between the two input populations, although their absolute zero
points differ.
\end{EN}

\begin{physbox}{高/低 $\alpha$ 双序列为什么重要}
银河系的恒星在 $[\alpha/\mathrm{Fe}]$--$[\mathrm{Fe}/\mathrm{H}]$ 平面上分成两条序列：高 $\alpha$ 星形成早、增丰快（以 II 型超新星为主，$\alpha$ 元素多）；低 $\alpha$ 星形成晚、增丰慢（Ia 型超新星贡献了大量铁，压低了 $\alpha/$铁比）。这两条序列是银河系化学演化与并合历史的化石记录。一个丰度分析流水线最重要的检验之一就是：独立拟合能否重现这条双序列，而不是把它涂抹成单一趋势。本文的直接物理拟合做到了——说明 12 个丰度坐标各自携带了真实的天体物理信息，而非流水线伪影。Si 的刚性偏移则反映零点约定差异（DR14 做了经验零点修正而本文没有），属于标度问题而非形态错误。
\end{physbox}

近乎刚性的 Si 偏移与标度差异一致：DR14 对单元素丰度施加了经验温度与零点定标，Si 的调整幅度位居前列，而我们不做任何 APOGEE 零点修正 \citep{Holtzman2018}。Mn 向高金属丰度上升，与依赖金属丰度的产率及可观的 Ia 型超新星贡献一致；Ni 大体跟随 Fe \citep{Weinberg2019}。C 与 N 保持分别受约束，保留了 APOGEE 红巨星中随质量与演化变化的 C/N 信息 \citep{CharbonnelLagarde2010,Masseron2015,Martig2016}。最终的丰度修正通常在百分之几 dex 量级，不改变这一族群形态。

\begin{EN}
The nearly rigid Si offset is consistent with a scale difference because DR14 applies empirical
temperature and zero-point calibrations to individual abundances, with Si
among its larger adjustments, while we apply no APOGEE zero-point correction
\citep{Holtzman2018}. Mn rises toward higher metallicity, consistent with
metallicity-dependent yields and a
substantial Type Ia supernova contribution, while Ni broadly tracks Fe
\citep{Weinberg2019}. C and N remain separately constrained, preserving the
mass- and evolution-dependent C/N information measured in APOGEE red giants
\citep{CharbonnelLagarde2010,Masseron2015,Martig2016}. The final
abundance corrections are typically at the hundredth-dex level and leave this
population morphology unchanged.
\end{EN}

谱线表的迁移效果是独立可测的。在匹配的对照样本中，定标降低了每颗星的每有效像素平均卡方 $Q$：中位数从 3.32 降到 1.56，成对下降的中位数为 53\%。

\begin{EN}
The line-list transfer is independently measurable. In a matched control
sample, the calibration lowers the mean chi-square per valid pixel, $Q$, for
every star. The median $Q$ falls from
3.32 to 1.56, with a median paired reduction of 53 percent.
\end{EN}

定标谱线表在不强迫与 ASPCAP 丰度标度一致的情况下改善了像素级目标。剩余零点差异可来自大气物理、太阳混合物、谱线选择、掩模与连续谱处理。光谱拟合与逐元素序列共同表明：直接物理前向模型无需光谱模拟器即可恢复自洽的巡天化学。要建立定标后的丰度标度，则需要超出本概念验证的更大标准星集合与巡天专属的谱线整编。

\begin{EN}
The calibrated line list improves the pixel-level objective without forcing
agreement with the ASPCAP abundance scale. The remaining zero-point differences
can reflect the atmosphere physics, solar mixture, line selection, masks, and
continuum treatment. Together, the spectral fit and the element-by-element
sequences show that a direct physical forward model can recover coherent survey
chemistry without a spectral emulator. Establishing a calibrated abundance
scale would require a broader standard-star set and survey-specific line
curation beyond this proof of concept.
\end{EN}

\begin{figure}[t]
\centering
\includegraphics[width=0.98\textwidth]{figure_11.pdf}
\caption{代表性的 APOGEE DR14 拟合（收敛大气计算之后）。黑色为观测归一化光谱，橙虚线为经速度平移、致宽、线扩展函数卷积与像素采样后的轮廓化 \paynezero\ 模型，灰带为采用的 $1\sigma$ 不确定度。上光谱与残差面板覆盖全部三个探测器；下三对放大每个探测器的一段富线区间。残差以像素不确定度为单位。（Representative APOGEE DR14 fit.）}
\label{fig:apogee_spectrum}
\end{figure}

\begin{figure}[t]
\centering
\includegraphics[width=0.98\textwidth]{figure_12.pdf}
\caption{均衡 1,600 目标 APOGEE 样本的化学结构。前两图为 DR14 定标星表与 \paynezero\ 拟合的 $[\mathrm{Mg}/\mathrm{Fe}]$--$[\mathrm{Fe}/\mathrm{H}]$，橙、蓝标记 APOGEE 高、低 $\alpha$ 样本；第三图比较两者的铁丰度，虚线为相等线。等值线包围逆选择权重密度的 50、80、95\%。直接拟合保留了两条化学序列，同时允许丰度零点平移。（Chemical structure of the balanced 1,600-target APOGEE sample.）}
\label{fig:apogee_structure}
\end{figure}

\begin{figure}[t]
\centering
\includegraphics[width=0.98\textwidth]{figure_13.pdf}
\caption{均衡 1,600 目标 \paynezero\ 拟合的逐元素丰度分布。每面板为一个 $[\mathrm{X}/\mathrm{Fe}]$ 对 $[\mathrm{Fe}/\mathrm{H}]$；橙、蓝点与等值线分别来自高、低 $\alpha$ 样本；灰虚等值线为同一批恒星与权重下 APOGEE DR14 的 50、80\% 密度。拟合的丰度平面呈现连贯的、随元素变化的序列，而非单一整体金属丰度坐标。（Elemental-abundance distributions from the balanced 1,600-target fit.）}
\label{fig:apogee_elements}
\end{figure}

\section{讨论 Discussion}\label{sec:discussion}

直接的从头算拟合改变了恒星推断中两类学习方法的角色。光谱模拟器对物理光谱库做插值，数据驱动流量模型从巡天观测中学习；二者都在恒星标签与被求值的光谱之间插入一个学习到的“标签到流量”关系。当合成以秒计、收敛大气以秒到数十秒计，物理前向模型就可以留在推断、谱线定标与验证内部。学习型大气模型可加速探索，而物理收敛仍决定被接受的大气与光谱。这改变了我们使用模拟器的方式，也改变了我们研究“观测—模型差距”的方式。

\begin{EN}
Direct ab initio fitting changes the role of two learned approaches in stellar
inference. Spectral emulators interpolate libraries of
physical spectra, while data-driven flux models learn from survey observations.
Both classes insert a learned label-to-flux relation between the stellar labels and the
spectrum being evaluated. When synthesis takes seconds and a converged
atmosphere takes seconds to tens of seconds, the physical forward model can
instead remain inside inference, line calibration, and validation. Learned
atmosphere models can accelerate exploration while physical convergence still
determines the accepted atmosphere and spectrum. This changes how we use
emulators and how we investigate the observation--model gap.
\end{EN}

\subsection{直接合成与光谱模拟器的角色 Direct synthesis and the role of spectral emulators}\label{sec:discussion_emulators}

物理自洽的丰度分析应该把相关恒星标签一起变化并拟合全光谱。一种元素改变的不只是归属它的谱线：镁改变电子供应从而改变 H$^-$ 连续谱；分子平衡与谱线覆盖在各物种间重新分配流量，并可改变大气结构 \citep{John1988,Gustafsson2008,Meszaros2012,Ting2018Oxygen,Matsuno2024}。只用某元素的具名特征测量该丰度、同时冻结模型的其余部分，会错过这些耦合响应。

\begin{EN}
A physically consistent abundance analysis should vary the relevant stellar
labels together and fit the full spectrum. An element changes more than the
lines assigned to it. Magnesium changes the electron supply and therefore the
H$^-$ continuum. Molecular balance and line blanketing redistribute flux among
species and can alter the atmospheric structure
\citep{John1988,Gustafsson2008,Meszaros2012,Ting2018Oxygen,Matsuno2024}.
Measuring each abundance only from its named features while holding the rest of
the model fixed misses these coupled responses.
\end{EN}

困难在于维度：每增加一个丰度，标签空间扩大一维，规则插值网格迅速稀疏 \citep{Ting2016}。Starfish、The Payne 及后续神经架构等概率与神经模拟器用物理光谱的平滑插值替代了网格 \citep{Czekala2015,Ting2019,Rozanski2025a}。我们当年提出 The Payne，是为了让“从自洽从头算模型做同时全光谱推断”在计算上可行——神经模拟器是加速该计算的手段，而非其科学目标。但这个名字后来成了模拟器本身的代称，而非物理拟合框架。

\begin{EN}
The difficulty is dimensionality. Each abundance expands the label space, and
regular interpolation grids become rapidly sparse \citep{Ting2016}.
Probabilistic and neural emulators such as Starfish, The Payne, and later neural
architectures replaced that grid with smooth interpolation of physical spectra
\citep{Czekala2015,Ting2019,Rozanski2025a}. We introduced The Payne to make
simultaneous full-spectrum inference from self-consistent ab initio models
computationally practical. A neural emulator was the means of accelerating
that calculation, not its scientific objective. The name nevertheless became
shorthand for the emulator rather than the physical fitting framework.
\end{EN}

\paynezero\ 回归这一最初目标，且不使用学习到的“标签到流量”插值器：“Zero”表示零层光谱模拟。每条试验光谱都针对所请求的标签与混合物直接合成，被接受的解用物理收敛大气检验。全光谱拟合因此可以保留所有标签的耦合响应，而无需构建高维光谱训练网格。

\begin{EN}
\paynezero\ returns to that original goal without a learned label-to-flux
interpolator. "Zero" denotes zero spectral-emulation layers. Every trial
spectrum is calculated by direct synthesis for the requested labels and
mixture, and the accepted solution is checked with a physically converged
atmosphere. Full-spectrum fitting can therefore retain the coupled response of
all labels without constructing a high-dimensional spectral training grid.
\end{EN}

\subsection{从学习流量关系到物理定标 From learned flux relations to physical calibration}

数据驱动的“标签到流量”关系从另一方向处理观测—模型差距。例如 The Cannon 直接从光谱与参考标签学习，不求值显式的物理前向模型。这条经验路线能复现合成库中缺失的可重复结构，其隐患是归因：训练样本中相关的标签可能产生有用的光谱预测器，却未必识别出某个特征背后的物理过程。

\begin{EN}
Data-driven label-to-flux relations address the observation--model gap from the
other direction. The Cannon, for example, learns directly from spectra and
reference labels rather than evaluating an explicit physical forward model.
This empirical route can reproduce repeatable structure missing from a
synthetic library. Its caveat is attribution. Correlated labels
in the training sample can produce a useful spectral predictor without
necessarily identifying the physical process responsible for a feature.
\end{EN}

当目标是改进前向模型时，直接从头算拟合检验的是被提议的物理原因。结构化残差可来自谱线数据、缺失不透明度以及更广的合成—观测失配 \citep{Czekala2015,OBriain2021}；大气、连续谱、地球大气与数据约减误差也能产生类似残差，但它们进入计算的位置不同。可微分的物理模型暴露这些输入，并通过不透明度与辐射转移检验被提议的原因。

\begin{EN}
When the goal is to improve the forward model, direct ab initio fitting tests
the proposed physical cause. Structured residuals can arise from line data,
missing opacity, and broader synthetic--observed mismatch
\citep{Czekala2015,OBriain2021}. Atmosphere, continuum, telluric, and reduction
errors can produce related residuals, but they enter the calculation in
different places. A
differentiable physical model exposes these inputs and tests a proposed cause
through opacity and radiative transfer.
\end{EN}

谱线表定标演示了这条路线：我们一次优化超过 $10^5$ 个耦合的振子强度与阻尼修正（包括混合跃迁），在单张 H100 上约一分钟得到共享太阳—大角星解。这些值是有效的标准星定标而非实验室测量；它们在冻结的 APOGEE 对照集上的表现检验了迁移效果，而未用 ASPCAP 标签来决定修正。

\begin{EN}
The line-list calibration demonstrates this route. We optimize more than
$10^5$ coupled oscillator-strength and damping corrections at once, including
blended transitions, and obtain the shared Sun--Arcturus solution in about one
minute on one H100. These values are effective standard-star calibrations rather
than laboratory measurements. Their performance on a frozen APOGEE control set
tests transfer without using ASPCAP labels to determine the corrections.
\end{EN}

生产级定标可以使用多个基准标准星 \citep{Heiter2015,Jofre2015,Smith2021}。我们建议联合求解它们的共享谱线参数，同时让每颗星的标签、连续谱与仪器响应自由变化。标准星之间的合适权重不由优化器决定：太阳与大角星探测不同的压强与分子区间，本文的等权只演示容量，而非定义一个普适定标目录。生产分析应针对被研究的恒星族群选择其标准星、权重、实验室先验与留出测试。

\begin{EN}
A production calibration could use several benchmark standards
\citep{Heiter2015,Jofre2015,Smith2021}. We propose solving their shared line
parameters while allowing each star's labels, continuum, and instrument
response to vary. The appropriate balance among standards is not determined by the
optimizer. The Sun and Arcturus probe different pressure and molecular regimes,
so the equal weighting used here demonstrates capacity rather than defining a
universal calibrated catalog. A production analysis should choose its standards,
weights, laboratory priors, and held-out tests for the stellar population being
analyzed.
\end{EN}

大幅度的不透明度修订随后会触发新的大气计算与新一轮谱线数据更新。这一外层迭代保留了物理耦合，并能暴露谱线定标无法修复的差异——包括大气迭代本身的失败。

\begin{EN}
Large opacity revisions would then trigger a new atmosphere
calculation and another line-data update. This outer iteration preserves
physical coupling and reveals discrepancies that line calibration cannot
repair, including failures of the atmosphere iteration itself.
\end{EN}

\subsection{大气边界与约化求解器 Atmosphere boundaries and reduced solvers}\label{sec:discussion_fixed_point}

一致性检验覆盖热矮星、太阳、红巨星与 K 矮星区间，巡天拟合则限于巨星。这些结果不意味着同一数值策略在每个恒星区间都保持高效。两个边缘案例指出了可能需要不同表示或求解器的位置。

\begin{EN}
The parity tests span hot-dwarf, solar, red-giant, and K-dwarf regimes, while
the survey fit is restricted to giants. These results do not imply that one
numerical strategy will remain efficient across every stellar regime. Two edge
cases identify where a different representation or solver may be needed.
\end{EN}

第一个是大气初始化器的采样维度。直接丰度模型把六条相关深度廓线压缩成紧凑 PCA 表示，但其输入仍是 3 个恒星参数加 81 个独立丰度，84 维输入相对其维度而言采样稀疏——同样的采样问题也限制高维模型网格。相关的大气响应仍让初始化器可用。未来的初始化器可以使用更贴近电子供给、分子平衡与谱线覆盖的坐标，或用化学感知的编码器学习这样的分组。初始化器只需提出初始状态，因为物理求解器仍在所请求的逐元素丰度下确立被接受的大气。

\begin{EN}
The first is the sampling dimension of the atmosphere initializer. The
direct-abundance model compresses six correlated depth profiles into a compact
PCA representation, but its input still spans three stellar parameters
together with 81 independently varied abundances. The resulting 84-coordinate
input is sparse relative to its dimension. The same sampling problem limits
high-dimensional model grids. Correlated atmospheric
responses nevertheless make the initializer useful. Future
initializers could use coordinates tied more closely to electron donation,
molecular balance, and line blanketing, or learn such groupings with a
chemistry-aware encoder. The initializer need only propose a starting state
because the physical solver still establishes the accepted atmosphere at the
requested individual abundances.
\end{EN}

第二个是冷分子大气的数值区间变化。一些更冷的高重力对照即使从邻近初始廓线出发也发展出增长的流量与加热残差，而另一些邻近起点却收敛——收敛性在该区间对初始结构变得更敏感。原 Kurucz 计算与 \paynezero\ 在其中一些对照上同样不通过收敛判据，这说明问题出在物理迭代本身，而非模拟器特有的误差。

\begin{EN}
The second is a change of numerical regime in cold molecular atmospheres. Some
cooler high-gravity controls develop growing flux and heating residuals even
from a nearby initial profile, while other nearby starts converge. Convergence
therefore becomes more sensitive to the initial structure across this regime.
Both the original Kurucz calculation and \paynezero\ fail the same convergence
criterion in some of these controls, which points to the physical iteration
rather than an emulator-specific error.
\end{EN}

\begin{physbox}{冷星大气为什么难收敛}
温度降到约 4000 K 以下，大气物理发生质变：分子（TiO、H$_2$O、CO 等）大量形成并主导不透明度；分子平衡对温度极其敏感（指数依赖）；对流层可向上延伸到直接形成光谱的层次。于是一个小的温度修正会同时重组分子平衡、不透明度、物态方程和绝热梯度，并移动对流边界——下一趟修正可能不减反增。用不动点语言说：该区间雅可比谱半径越过 1，经典逐次迭代失去局部稳定性。这是 Kurucz 类型求解器的已知固有边界，本文的重要贡献是诚实地标出了它，而不是声称学习型初值能绕过它。
\end{physbox}

该区间的小温度变化会一起重组分子平衡 \citep{BarklemCollet2016}、不透明度、物态方程与绝热梯度，对流也可能触及决定出射流量的层次 \citep{Gustafsson2008}。一趟迭代随后可能把罗斯兰不透明度、柱质量映射与对流边界移动得足以让下一步修正增大而非缩小——通常的不动点迭代失去了局部稳定性。

\begin{EN}
Small temperature changes in this regime reorganize molecular equilibrium
\citep{BarklemCollet2016}, opacity, the equation of state, and the adiabatic
gradient together. Convection can also reach the layers that set the emergent
flux \citep{Gustafsson2008}. One iteration
can then shift the Rosseland opacity, column-mass mapping, and convective
boundary enough that the next correction grows rather than shrinks. The usual
fixed-point iteration has lost local stability.
\end{EN}

在约化的平面平行表述中，我们选择温度 $T(\tau_{\rm R})$ 与单调柱质量 $m(\tau_{\rm R})$ 作为两条被迭代廓线。固定 $T_{\rm eff}$、$\log g$、$\xi$、丰度与罗斯兰网格后，流体静力学平衡与物态方程即可恢复压强、密度、电子密度与布居 \citep{HubenyMihalas2014}；不透明度设定柱质量与罗斯兰深度的关系，辐射转移设定加速度。约化求解器将强制两条耦合的根条件——能量平衡，以及柱质量与罗斯兰光学深度的一致性：
\begin{equation}
\begin{aligned}
\mathcal{R}_{T}[T,m]
&=F_{\rm rad}[T,m]+F_{\rm conv}[T,m]
-\sigma_{\rm SB}T_{\rm eff}^{4}=0,\\
\mathcal{R}_{m}[T,m]
&=\frac{\mathrm{d}\tau_{\rm R}}{\mathrm{d}m}-\kappa_{\rm R}[T,m]=0 .
\end{aligned}
\end{equation}
\noindent 这些残差提示在 $T$ 与 $m$ 中求解，而非为全部六条存档廓线做初始化。重算其余四个场并未让学习型冷矮星 $T$、$m$ 廓线恢复收敛，说明不稳定并非简单因为“预测的场太多”。

\begin{EN}
For a reduced plane-parallel formulation, we choose temperature
$T(\tau_{\rm R})$ and monotonic column mass $m(\tau_{\rm R})$ as the two
iterated profiles. At fixed $T_{\rm eff}$, $\log g$, $\xi$, abundances, and
Rosseland grid, hydrostatic balance and the equation of state then recover
pressure, density, electron density, and populations
\citep{HubenyMihalas2014}. Opacity
sets the relation between column mass and Rosseland depth, while radiative
transfer sets the acceleration. A reduced solver would enforce two coupled root
conditions, energy balance and consistency between column mass and Rosseland
optical depth,
\begin{equation}
\begin{aligned}
\mathcal{R}_{T}[T,m]
&=F_{\rm rad}[T,m]+F_{\rm conv}[T,m]
-\sigma_{\rm SB}T_{\rm eff}^{4}=0,\\
\mathcal{R}_{m}[T,m]
&=\frac{\mathrm{d}\tau_{\rm R}}{\mathrm{d}m}-\kappa_{\rm R}[T,m]=0 .
\end{aligned}
\end{equation}
\noindent These residuals motivate a solver in $T$ and $m$ rather than an
initializer for all six stored profiles. Recomputing the other four fields did
not restore convergence from learned cool-dwarf $T$ and $m$ profiles, so the
instability is not simply caused by predicting too many fields.
\end{EN}

原生的双廓线求解器必须对每组试探 $T$、$m$ 重算物态方程、不透明度、转移与对流，再联立求解两条残差。在探索性计算中，邻近物理模型估计两条残差随 $T$、$m$ 的变化并给出局部修正；残差确实下降了，但所需的完整大气求值次数多于“学习型六场初值 + 物理求解器迭代”。因此我们保留后者作为默认。若双廓线求解器的响应可以在不做多次完整物理求值的情况下估计，它仍是可选的备用方案。不同恒星区间最终可能需要不同的、经过验证的修正器。何时切换区间、同时保持同一套物理验收测试，成为一个工作流问题——它要求软件边界暴露物理状态、收敛性与验证，而不是藏在一个不透明计算内部。

\begin{EN}
A native two-profile solver must recompute the equation of state, opacity,
transfer, and convection for each trial $T$ and $m$, then solve both residuals
together. In exploratory calculations, nearby physical models estimated how
both residuals change with $T$ and $m$ and supplied local corrections. The
residuals fell, but the calculation required more complete atmosphere
evaluations than the learned six-field start followed by the physical solver iteration.
We therefore retain the latter as the default. A two-profile solver remains a
possible fallback if its response can be estimated without many full physical
evaluations. Different stellar regimes may ultimately require different
validated correctors. Choosing when to switch regimes, while preserving the
same physical acceptance tests, becomes a workflow problem. It motivates
software boundaries that expose physical state, convergence, and validation
rather than hiding them inside one opaque calculation.
\end{EN}

\subsection{作为科学接口的可读软件 Readable software as a scientific interface}

在数值区间之间安全切换，要求物理状态、残差与验收测试保持可见。原 Kurucz 工作流程把大量物理编码在由中转文件连接的紧凑 Fortran 程序中 \citep{Kurucz2005}。重实现过程中我们发现，简写的标识符、超长的例程与隐式共享状态让单个约定难以隔离：关于谱线放置、不透明度累积或状态传递的假设即使改变最终光谱，也难以审计与扩展。

\begin{EN}
Switching safely between numerical regimes requires the physical state,
residuals, and acceptance tests to remain visible. The original Kurucz workflow
encodes extensive physics in compact Fortran programs connected by staged files
\citep{Kurucz2005}. During reimplementation, we found that terse
identifiers, long routines, and implicit shared state made individual
conventions difficult to isolate. Assumptions about line placement, opacity
accumulation, or state transfer can therefore be difficult to audit and extend
even when they change the final spectrum.
\end{EN}

\paynezero\ 用具名物理数组、显式单位与阶段局部接口替代了简写标识符与隐式交接。这些选择把数值约定暴露给审计、复现、教学与扩展。

\begin{EN}
\paynezero\ replaces terse identifiers and implicit handoffs with named physical
arrays, explicit units, and stage-local interfaces. These choices expose
numerical conventions for audit, reproduction, teaching, and extension.
\end{EN}

同样的显式接口使 \paynezero\ 适合未来的智能体（agentic）工作流。Model Context Protocol 工具或任务专属技能可以把大气求解、光谱合成、谱线定标与验证暴露为带记录输入、输出与溯源的类型化操作。\footnote{见 Model Context Protocol 规范：\url{https://modelcontextprotocol.io/specification/2025-06-18/server/index}。}智能体可以检查收敛与流量残差，识别出冷的高重力模型进入了不同的数值区间，选择预验证的求解器或参数化，并启动相关诊断测试。智能体负责编排这些路径，而非为自己的结果背书；确定性的物理、数值与溯源门槛仍负责决定接受与否。

\begin{EN}
The same explicit interfaces make \paynezero\ suitable for future agentic
workflows. Model Context Protocol tools or task-specific skills could expose
atmosphere solving, spectrum synthesis, line calibration, and validation as
typed operations with recorded inputs, outputs, and provenance.
An agent could
inspect convergence and flux residuals, recognize that a cool high-gravity
model has entered a different numerical regime, select a prevalidated solver or
parameterization, and launch the relevant diagnostic tests. The agent would
orchestrate these routes rather than certify its own result. Deterministic
physical, numerical, and provenance gates would still decide acceptance.
\end{EN}


\section{结论 Conclusion}\label{sec:conclusion}

光谱模拟器与数据驱动流量模型解决了恒星光谱学的两个长期难题：物理计算太慢，无法在推断中反复执行；且其光谱不能复现所有观测特征。由此产生的“标签到流量”关系让巡天分析切实可行，但光谱模拟器带来插值误差，数据驱动模型则可能继承训练标签间的相关性。我们证明：在本文演示的一维 LTE 区间内，这一中间层对计算可行性已不再是必需品。

\begin{EN}
Spectral emulators and data-driven flux models addressed two longstanding
problems in stellar spectroscopy. Physical calculations were too slow to
repeat inside inference, and their spectra did not reproduce every observed
feature. The resulting label-to-flux relations made survey analysis practical,
but spectral emulators add interpolation error, while data-driven models can
inherit correlations among their training labels. We show that this
intermediate layer is no longer required for computational tractability in the
demonstrated one-dimensional LTE regime.
\end{EN}

\paynezero\ 通过面向现代处理器重组计算而不改变其物理目标，改变了这一权衡。GPU 原生合成把 $R_{\rm grid}=300{,}000$ 的 300--1000 nm 太阳计算从原 Kurucz 程序的约 760 秒降到单张 H100 上的 14 秒，APOGEE 区间约 1 秒。多核大气 kernel 使每趟物理迭代快 6--7 倍，学习型初始大气在受控测试中再把所需迭代次数压缩约 3 倍。所得光谱与原始计算保持实用一致。大气与合成计算因此可以留在推断内部，其假设与残差被直接评估。

\begin{EN}
\paynezero\ changes this tradeoff by reorganizing the calculation for modern
processors without changing its physical target. GPU-native synthesis reduces
a 300--1000~nm solar calculation at $R_{\rm grid}=300{,}000$ from about 760~s in
the original Kurucz programs to 14~s on one H100, while the APOGEE interval
takes about 1~s. Multicore atmosphere kernels make each physical iteration six
to seven times faster, and a learned starting atmosphere reduces the required
iteration count by about another factor of three in the controlled tests. The
resulting spectra remain in practical parity with the original calculations.
The atmosphere and synthesis calculation can therefore remain inside
inference, where their assumptions and residuals are evaluated directly.
\end{EN}

同一计算结构把“观测—模型差距”变成可处理的物理优化问题。自动微分联合更新超过 $10^5$ 个耦合的振子强度与阻尼参数；共享太阳—大角星定标在单张 H100 上约一分钟，把逐条跃迁的任务变成一次联合辐射转移计算。

\begin{EN}
The same computational structure turns the observation--model gap into a
tractable physical optimization problem. Automatic differentiation updates
more than $10^5$ coupled oscillator strengths and damping parameters together.
The shared Sun--Arcturus calibration takes about one minute on one H100,
replacing a transition-by-transition task with one joint radiative-transfer
calculation.
\end{EN}

这些能力在我们的 APOGEE 演示中汇合。GPU 常驻的速度平移、致宽、波长依赖 LSF 卷积与探测器采样相对合成的开销可忽略。对 APOGEE 巨星，拟合从约减后的一维光谱直接求解 12 个丰度坐标以及仪器与连续谱项：H100 搜索约 40 秒/星，物理大气收敛在多核 CPU 上独立运行。恢复的化学在跨越巨星支与红团簇的恒星上，对全部 11 个拟合丰度比呈现连贯的、随元素变化的趋势。

\begin{EN}
These capabilities meet in our APOGEE demonstration. GPU-resident velocity
shifting, broadening, wavelength-dependent LSF convolution, and detector
sampling add negligible cost to synthesis. For APOGEE giants, the fit solves
12 abundance coordinates together with the instrument and continuum terms
directly from reduced one-dimensional spectra. The H100 search takes about
40~s per star, while physical atmosphere convergence runs independently on
multicore CPUs. The recovered chemistry shows coherent element-dependent trends
across all eleven fitted abundance ratios for stars spanning the giant branch
and red clump.
\end{EN}

由此得到的计算规模改变了巡天中“可行”的含义：对 $10^5$ 条光谱做实测的直接合成搜索，对应约一千 GPU 小时，按常见学术云价格约合几千美元。独立恒星之间可以直接分发而不改变计算；大气收敛使用多核 CPU 工人，可独立于 GPU 合成任务运行。GPU 需求因此主要由直接合成搜索与直接光谱决定，而非大气迭代。

\begin{EN}
The resulting computational scale changes what is practical for a survey. The
measured direct-synthesis search for $10^5$ spectra corresponds to about
one thousand GPU-hours and a few thousand US dollars at typical academic cloud
rates.
Independent stars can be distributed without changing the calculation.
Atmosphere convergence uses multicore CPU workers and can run
independently of the GPU synthesis tasks. GPU demand is therefore set mainly by
the direct-synthesis search and direct spectra rather than by the atmosphere
iteration.
\end{EN}

恒星光谱学走过了从手工逐个操作物理程序，到预计算网格与学习近似，再到百万恒星巡天的历程。现代硬件现在允许物理计算在这一规模上回到推断循环之内。这并不能消除三维对流或非 LTE 谱线形成带来的偏离 \citep{Asplund2005,Nordlund2009,Magic2013,BergemannNordlander2014,Amarsi2019}，也不消除谱线数据或仪器建模的局限——但它让它们成为可见的改进目标。大气可以按请求的混合物重算，光谱可以在优化器内合成，原子数据可以通过同一辐射转移计算定标。物理前向模型可以重新成为恒星推断的工作基础，而不只是中间学习流量模型的训练来源。

\begin{EN}
Stellar spectroscopy has moved from individually operated physical programs,
through precomputed grids and learned approximations, to surveys of millions
of stars. Modern hardware now allows the physical calculation to return to the
inference loop at that scale. This does not remove departures associated with
three-dimensional convection or non-LTE line formation
\citep{Asplund2005,Nordlund2009,Magic2013,BergemannNordlander2014,Amarsi2019}, nor limitations
in line data or instrument modeling. It makes them
visible targets for improvement. Atmospheres can be recomputed at the requested
mixture, spectra can be synthesized inside the optimizer, and atomic data can
be calibrated through the same radiative-transfer calculation. Physical
forward models can again be the working foundation of stellar inference rather
than only the training source for an intermediate learned flux model.
\end{EN}

\begin{physbox}{本文未覆盖的物理：3D 与非 LTE}
最后一段的限定值得强调。本文全部计算基于一维、平面平行、LTE、混合长对流的 Kurucz 框架。真实恒星大气是三维的：对流元胞的涨落会系统性地改变温度结构与谱线形状（3D 效应）；在低密度高层大气，碰撞不足以维持局部热动平衡，能级布居由辐射场逐能级决定（NLTE 效应），对 Li、O、Mn 等元素的丰度影响可达零点几 dex。这些是模型本身的系统误差，与计算速度无关。\paynezero\ 的意义在于：当正演足够快，这些系统误差不再被埋在模拟器插值误差之下，而是可以被直接测量、检验和逐项改进。
\end{physbox}

\section*{致谢 Acknowledgments}

对 YST 而言，本项目同时是一项业余爱好，也是对现代大语言模型智能体用于科学软件的帕累托前沿的一次非正式测试。2024 年开始的尝试可以翻译孤立文件，但无法保持大型耦合代码库的逻辑。实际的转折点出现在 2025 年 12 月，Claude Opus 4.5 使 Kurucz 例程到 Python 的细致翻译变得可行。EK 完成了大量最初的、艰辛的 PyKurucz 翻译并开发了其早期 Numba 实现，即 Kim \& Ting (2026) 的 pyKurucz 研究。YST 主导了随后的 GPU 原生重设计、优化实验与 Payne Zero 开发。超越蛮力重写的进展依赖 YST 对 Kurucz 代码的经验，以识别常驻数据、kernel 边界、谱线筛选与并行结构。后期开发使用了 Claude Opus 4.8 与持续运行的 GPT-5.5 Codex 智能体，并有来自 Gemini 3.1 Pro 与 GLM 5.2 的独立批评。没有这些工具，这一规模的重写在同一时间表内是不切实际的。那些未成功的实验（包括备选的双廓线大气求解器）同样是有益的提醒：当前的智能体仍需要密切的领域指导。通用人工智能（AGI）本会让这项业余项目短得多。

\begin{EN}
For YST, this project was also a hobby and an informal test of the Pareto
frontier of modern large-language-model agents for scientific software.
Attempts beginning in
2024 could translate isolated files
but could not preserve the logic of a large coupled codebase. A practical
turning point came in December 2025, when Claude Opus 4.5 made careful
translation of Kurucz routines into Python tractable. EK carried out much of
the initial painstaking translation into PyKurucz and developed its early
Numba implementation, leading to the pyKurucz study of Kim \& Ting (2026).
YST led the subsequent GPU-native
redesign, optimization experiments, and Payne Zero development. Progress
beyond a brute-force rewrite depended on YST's experience with the Kurucz codes
to identify the resident data, kernel boundaries, line selection, and parallel
structure. Later development used Claude Opus 4.8 and sustained GPT-5.5 Codex
agents, with independent criticism from Gemini 3.1 Pro and GLM 5.2. A rewrite
at this scale would not have been practical on the same timescale without these
tools. The unsuccessful experiments, including the alternate two-profile
atmosphere solver, were equally useful reminders that current agents still
require close domain guidance. Artificial general intelligence (AGI) would
have made this hobby considerably shorter.
\end{EN}

智能体贡献了代码翻译、重构、测试、诊断、研究日志维护与文字修订。科学框架、引言、讨论与结论始于人类草稿；部分技术段落始于智能体对代码与研究日志的总结。作者审阅了保留的代码、计算、图、引文与文字，并对结果负责。一些意外的个人境遇也给了 YST 完成本项目所需的不被打扰的时间。

\begin{EN}
The agents contributed code translation, refactoring, tests, diagnostics,
research-log maintenance, and prose revision. The scientific framing,
Introduction, Discussion, and Conclusion began from human drafts. Some
technical passages began from agent summaries of the code and research logs.
The authors reviewed the retained code, calculations, figures, citations, and
text and take responsibility for the results. Unexpected personal
circumstances also gave YST the uninterrupted time needed to finish this
project.
\end{EN}

YST 感谢亚历山大·冯·洪堡基金会的洪堡研究奖学金资助。

\begin{EN}
YST acknowledges support from a Humboldt Research Fellowship from the
Alexander von Humboldt Foundation.
\end{EN}


\appendix

\section{对拟合恒星参数的敏感性 Sensitivity to Fitting the Stellar Parameters}
\label{app:stellar_label_sensitivity}

主 APOGEE 实验保留了 DR14 发表的、经光度定标的 $T_{\rm eff}$ 与经星震定标的巨星 $\log g$，并固定 $\xi$，以单独考察化学与仪器拟合 \citep{Holtzman2018}。我们用相同的官方光谱、样本、谱线定标、仪器算子、掩模、不确定度与连续谱处理来检验该选择：从主丰度解出发，把三个恒星参数与同样的 12 个丰度坐标一起拟合，并收敛最终的物理大气。图~\ref{fig:appendix_stellar_label_sensitivity} 与图~\ref{fig:appendix_abundance_sensitivity} 把所得恒星与化学分布同固定参数解及 APOGEE DR14 比较。

\begin{EN}
The main APOGEE experiment retains the photometrically calibrated $T_{\rm eff}$
and asteroseismically calibrated giant $\log g$ published in DR14, and fixes
$\xi$ to isolate the chemical and instrumental fit \citep{Holtzman2018}. We
test that choice with the same official spectra, sample, line calibration,
instrument operator, masks, uncertainties, and continuum treatment. Starting
from the main abundance solution, we fit those three stellar parameters
together with the same 12 abundance coordinates and converge the final physical
atmosphere. Figures
\ref{fig:appendix_stellar_label_sensitivity} and
\ref{fig:appendix_abundance_sensitivity} compare the resulting stellar and
chemical distributions with the fixed-parameter solution and APOGEE DR14.
\end{EN}

完整物理重拟合给出的逆选择权重中位变化为：$T_{\rm eff}$ $-200$ K、$\log g$ $+0.17$ dex、$\xi$ $-0.32$ km s$^{-1}$、$[\mathrm{Fe}/\mathrm{H}]$ $-0.06$ dex。在初始红团簇聚集区的 694 颗星中，中位变化为 $-220$ K 与 $+0.17$ dex。巨星支脊线与红团簇聚集仍清晰可见，但其绝对位置随所采用的恒星参数标度移动。

\begin{EN}
The full physical refit gives inverse-selection-weighted median changes of
$-200$~K in $T_{\rm eff}$,
$+0.17$~dex in $\log g$, $-0.32$~km~s$^{-1}$ in $\xi$, and $-0.06$~dex in
$[\mathrm{Fe}/\mathrm{H}]$. Among the 694 stars in the initial red-clump
concentration, the median changes are $-220$~K and $+0.17$~dex. The giant-branch
ridge and red-clump concentration remain visible, but their absolute location
moves with the adopted stellar-parameter scale.
\end{EN}

丰度中位变化从氮的 $-0.23$ dex、氧的 $-0.22$ dex 到硫的 $+0.22$ dex 不等；镁变化 $-0.03$ dex。丰度序列保持存在，而其零点响应恒星参数的变化。因此，更冷的光谱温度标度会传播为随元素变化的丰度零点。

\begin{EN}
The median abundance changes range from $-0.23$~dex for nitrogen and
$-0.22$~dex for oxygen to $+0.22$~dex for sulfur. Magnesium changes by
$-0.03$~dex. The abundance sequences persist while their zero points respond
to the stellar parameters. The cooler spectroscopic temperature scale therefore
propagates into element-dependent abundance zero points.
\end{EN}

生产级分析可以对光谱、宽波段能谱分布与视差做联合拟合，锚定 $T_{\rm eff}$ 与 $\log g$ 并吸收这一标度差异 \citep{Cargile2020}。本概念验证不包含这些外部信息。

\begin{EN}
A joint fit to the spectrum, broadband spectral energy distribution, and parallax can
anchor $T_{\rm eff}$ and $\log g$ and absorb this scale difference in a
production analysis \citep{Cargile2020}. We do not include that external
information in this proof of concept.
\end{EN}

\begin{physbox}{为什么恒星参数与丰度会“耦合”}
附录展示的是光谱分析中著名的\textbf{参数耦合}（或退化）现象：降低 $T_{\rm eff}$ 会让模型光谱整体变化，拟合器会调整丰度去补偿——例如温度影响电离平衡，从而改变电离与中性谱线的相对强度，进而影响推出的元素丰度。本附录定量说明：采用更冷的自由光谱温度（中位 $-200$ K）后，丰度零点整体移动最多约 0.2 dex，但族群的相对形态（双 $\alpha$ 序列、红团簇结构）不变。这说明本文主文结论对“固定 DR14 恒星参数”这一选择是稳健的，但绝对丰度标度需要外部信息（光度、视差）来锚定。
\end{physbox}

\begin{figure}[t]
\centering
\includegraphics[width=0.98\textwidth]{figure_14.pdf}
\captionof{figure}{$T_{\rm eff}$、$\log g$、$\xi$ 与 12 个丰度坐标联合拟合的敏感性。左：合并样本的逆选择权重等值线，灰虚线保留 DR14 恒星参数，深色实线为最终物理收敛拟合。右：固定参数的铁丰度为灰色空心圆，最终值按 APOGEE 高/低 $\alpha$ 样本着色。（Sensitivity to fitting the stellar parameters.）}
\label{fig:appendix_stellar_label_sensitivity}
\vspace{8pt}
\includegraphics[width=0.98\textwidth]{figure_15.pdf}
\captionof{figure}{$T_{\rm eff}$、$\log g$、$\xi$ 与丰度最终物理收敛拟合的逐元素丰度分布。布局与标度同图~\ref{fig:apogee_elements}；彩色点与等值线为两个 APOGEE 样本的最终值，灰虚等值线为 APOGEE DR14 的 50、80\% 密度水平。（Elemental-abundance distributions from the final physically converged fit.）}
\label{fig:appendix_abundance_sensitivity}
\end{figure}



\bibliographystyle{plainnat}
\bibliography{references}

\end{document}
